%% This is a LaTeX document.  Hey, Emacs, -*- latex -*- , get it?
\newif\ifbeamermode
\beamermodetrue
\ifbeamermode
\documentclass[mathserif,a4paper,aspectratio=169]{beamer}
\else
\documentclass[a4paper]{article}
\usepackage{beamerarticle}
\usepackage[a4paper,margin=2.5cm]{geometry}
\fi
\usepackage[shorthands=off,french]{babel}
\usepackage[utf8]{inputenc}
\usepackage[T1]{fontenc}
\usepackage{lmodern}
\DeclareUnicodeCharacter{00A0}{~}
\DeclareUnicodeCharacter{2026}{...}
% Beamer theme:
\usetheme{Goettingen}
%\usecolortheme{albatross}
%\usecolortheme{lily}
%\setbeamercovered{transparent}
% A tribute to the worthy AMS:
\usepackage{amsmath}
\usepackage{amsfonts}
\usepackage{amssymb}
\usepackage{amsthm}
%
\usepackage{mathrsfs}
\usepackage{mathpartir}
\usepackage{flagderiv}
%
\usepackage{graphicx}
\usepackage{tikz}
\usetikzlibrary{arrows,automata,calc}
\usepackage{hyperref}
%
\newcommand{\itempoint}{\strut\hbox{\color{beamerstructure}\donotcoloroutermaths$\blacktriangleright$}\nobreak\hskip.5em plus.5em\relax}
\renewcommand{\thefootnote}{\textdagger}
\newcommand{\dbllangle}{\mathopen{\langle\!\langle}}
\newcommand{\dblrangle}{\mathclose{\rangle\!\rangle}}
\newcommand{\cps}{\mathrm{CPS}}
\newcommand{\dottedlimp}{\mathbin{\dot\Rightarrow}}
\newcommand{\dottedland}{\mathbin{\dot\land}}
\newcommand{\dottedlor}{\mathbin{\dot\lor}}
\newcommand{\dottedtop}{\mathord{\dot\top}}
\newcommand{\dottedbot}{\mathord{\dot\bot}}
\newcommand{\dottedneg}{\mathop{\dot\neg}}
\mathchardef\emdash="07C\relax
\newcommand{\mpdotsabove}[1]{\inferrule*{\vdots}{#1}}
\setlength{\derivskip}{4pt}
%
\mode<article>{
    \newcommand{\donotcoloroutermaths}{\relax}
    \renewcommand{\frametitle}[1]{\subsubsection*{#1 {\footnotesize[transp. \insertframenumber]}}}
    \newcommand{\myitemmark}{\hbox{$\blacktriangleright$}}
    \renewcommand{\itempoint}{\strut\myitemmark\nobreak\hskip.5em plus.5em\relax}
    \setlength{\parindent}{0pt}
    \renewcommand{\thepage}{II–\arabic{page}}
}
%
%
%
\title{{\footnotesize Partie II:} Typage simple et calcul propositionnel}
\subtitle{CSC-3TC34-TP / INF110 (Logique et Fondements de l'Informatique)}
\author[David Madore]{David A. Madore\\
{\footnotesize Télécom Paris}\\
\texttt{david.madore@}(nom de l'école)\texttt{.fr}}
\date{2023–2026}
\mode<presentation>{%
\beamertemplatenavigationsymbolsempty
\usenavigationsymbolstemplate{\vbox{\hbox{\footnotesize(II) \hyperlinkslideprev{$\leftarrow$}\insertframenumber/\inserttotalframenumber\hyperlinkslidenext{$\rightarrow$}}}}
}
\setbeamercolor{myhighlight}{fg=black,bg=white!90!green}
\colorlet{mydarkgreen}{green!50!black}
\begin{document}
\mode<article>{\maketitle\renewcommand{\labelitemi}{\myitemmark}}
%
\setlength\abovedisplayskip{2pt plus 2pt minus 2pt}
\setlength\belowdisplayskip{2pt plus 2pt minus 2pt}
%
\begin{frame}
\titlepage
{\footnotesize\center{\url{http://perso.enst.fr/madore/inf110/transp-inf110.pdf}}\par}
{\tiny
\immediate\write18{sh ./vc > vcline.tex}
\begin{center}
Git: \input{vcline.tex}
\end{center}
\immediate\write18{echo ' (stale)' >> vcline.tex}
\par}
\end{frame}
%
\section*{Plan}
\begin{frame}
\frametitle{Plan}
\tableofcontents
\end{frame}
%
\section{Généralités sur le typage}
\begin{frame}
\frametitle{Qu'est-ce que le typage ?}

{\footnotesize Philosophie opposée du « codage de Gödel » en
  calculabilité, lequel représente toute donnée finie par un
  entier.\par}

\medskip

\itempoint Informellement, un \textbf{système de typage} est une façon
d'affecter à toute \alert{expression} et/ou \alert{valeur} manipulée
par un langage informatique un \textbf{type} qui contraint son usage
{\footnotesize (p.ex., « entier », « fonction », « chaîne de
  caractères »)}.

\medskip

\textbf{Buts :}
\begin{itemize}
\item attraper plus tôt des \textbf{erreurs de programmation}
  {\footnotesize (« ajouter un entier et une chaîne de caractères »
    est probablement une erreur, ou demande une conversion
    explicite)},
\item éviter des \textbf{plantages ou problèmes de sécurité}
  {\footnotesize (exécution d'un entier)},
\item garantir certaines \textbf{propriétés du comportement} des
  programmes ($\rightarrow$ analyse statique), forcément de façon
  limitée (cf. théorème de Rice), p.ex., la \alert{terminaison},
\item aider à l'\textbf{optimisation} {\footnotesize (fonction pure :
  sans effet de bord)}.
\end{itemize}

\medskip

{\footnotesize Cf. systèmes d'unités (homogénéité) en physique.\par}

\end{frame}
%
\begin{frame}
\frametitle{Variétés de typage}

Il y a autant de systèmes de typage que de langages de programmation !

\medskip

\itempoint Typage \textbf{statique} (à compilation) vs
\textbf{dynamique} (à exécution) ou mixte.

{\footnotesize Mixte = « graduel ».  On préfère : erreur de
  compilation $>$ erreur à l'exécution $>$ plantage.\par}

\medskip

\itempoint Typage \textbf{inféré} par le langage ou \textbf{annoté}
par le programmeur.

{\footnotesize Type = promesse donnée par le langage au
  programmeur ou par le programmeur au langage ?\par}

\medskip

\itempoint Typage « sûr » ou partiel/contournable (\textit{cast} de
la valeur, manipulation de la représentation mémoire, suppression ou
report à l'exécution de vérification).

\medskip

\itempoint Typage superficiel {\footnotesize (« ceci est une liste »)}
ou complexe {\footnotesize (« ceci est une liste d'entiers »)}.

\medskip

\itempoint Diverses annotations possibles {\footnotesize (« cette
  fonction est pure », « soulève une exception »)}.

\medskip

\itempoint Liens avec les mécanismes de sécurité (qui peut faire
quoi ?), de gestion de la mémoire (système de typage linéaire),
l'évaluation (exceptions), la mutabilité.

\medskip

\itempoint Les types sont-ils citoyens de première classe
           {\footnotesize (:= manipulables dans le langage)} ?

{\footnotesize Juvénal : « Quis custodiet ipsos custodes? » Quel est
  le type des types eux-mêmes ?\par}

\end{frame}
%
\begin{frame}
\frametitle{Opérations de base sur les types}

En plus de types de base (p.ex. $\texttt{Nat}$ = entiers,
$\texttt{Bool}$ = booléens), les opérations suivantes sur les types
sont \emph{souvent} proposées par les systèmes de typage :

\bigskip

\itempoint Types \textbf{produits} (= couples, triplets, $k$-uplets).
P.ex. $\texttt{Nat} \times \texttt{Bool}$ = type des paires formées
d'un entier et d'un booléen.

\smallskip

Composantes éventuellement nommées $\rightarrow$ structures
(= enregistrements).

\smallskip

Produit vide = type trivial, $\texttt{Unit}$ (une seule valeur).

\bigskip

\itempoint Types \textbf{sommes} (= unions).  P.ex. $\texttt{Nat} +
\texttt{Bool}$ = type pouvant contenir un entier \emph{ou} un booléen,
avec un \alert{sélecteur} de cas.

\smallskip

Cas particulier : $\texttt{Unit} + \cdots + \texttt{Unit}$ = type
« énumération » (pur sélecteur).

\smallskip

Somme vide = type inhabité (impossible : aucune valeur).

\bigskip

\itempoint Types \textbf{fonctions} (= exponentielles).
P.ex. $\texttt{Nat} \rightarrow \texttt{Bool}$ (f\textsuperscript{n}
de $\texttt{Nat}$ vers $\texttt{Bool}$).

\bigskip

\itempoint Types \textbf{listes}.  P.ex. $\texttt{List}~\texttt{Nat}$
= type des listes d'entiers.

\end{frame}
%
\begin{frame}
\frametitle{Quelques fonctionnalités fréquentes}

\itempoint \textbf{Sous-typage} = les valeurs d'un type sont
automatiquement des valeurs possibles d'un autre.

\bigskip

\itempoint \textbf{Polymorphisme} = utilisation de plusieurs types
possibles, voire de n'importe quel type (cf. transp. suivant).
P.ex. la fonction « identité » $(\forall\mathtt{t})\; \mathtt{t}
\rightarrow \mathtt{t}$.

\bigskip

\itempoint \textbf{Familles de types} = fabriquent un type à
partir d'un (ou plusieurs) autres.  P.ex. $\texttt{List}$ (fabrique le
type « liste de $\mathtt{t}$ » à partir de $\mathtt{t}$).

\bigskip

\itempoint \textbf{Types récursifs} = construits par les opérateurs
(produits, sommes, fonctions, familles…) à partir des types
définis eux-mêmes.  P.ex. $\texttt{Tree} =
\texttt{List}~\texttt{Tree}$.

\bigskip

\itempoint \textbf{Types dépendants} = un type à partir d'une valeur.
P.ex. $k \mapsto \texttt{Nat}^k$.

\bigskip

\itempoint \textbf{Types opaques} (abstraits, privés…) = types dont
les valeurs sont cachées, l'usage est limité à une interface publique.

\end{frame}
%
\begin{frame}
\frametitle{Polymorphisme}

On distingue deux (trois ?) sortes de polymorphismes :

\medskip

\itempoint Polymorphisme \textbf{paramétrique} (ou « génériques ») :
la même fonction \alert{s'applique à l'identique} à une donnée de
n'importe quel type.

\smallskip

Exemples :
\begin{itemize}
\item $\texttt{head} : (\forall\mathtt{t})\; \texttt{List}~\mathtt{t} \to \mathtt{t}$ (renvoie le premier élément d'une liste)
\item $\lambda xy.\langle x,y\rangle :
  (\forall\mathtt{u},\mathtt{v})\; \mathtt{u} \to \mathtt{v} \to
  \mathtt{u}\times \mathtt{v}$ (fabrique un couple)
\item $\lambda xyz.xz(yz) :
  (\forall\mathtt{u},\mathtt{v},\mathtt{w})\; (\mathtt{u} \to
  \mathtt{v} \to \mathtt{w}) \to (\mathtt{u} \to
  \mathtt{v}) \to \mathtt{u} \to
  \mathtt{w}$
\end{itemize}

\smallskip

Pas seulement pour les fonctions !  $\texttt{[]} :
(\forall\mathtt{t})\; \texttt{List}~\mathtt{t}$ (liste vide)

\smallskip

{\footnotesize Et même : $\texttt{while~true~do~pass~done} :
  (\forall\mathtt{t})\; \mathtt{t}$ (boule infinie)\par}

\bigskip

\itempoint Polymorphisme \textbf{ad hoc} (ou « surcharge » /
« \textit{overloading} ») : la fonction \alert{agit différemment} en
fonction du type de son argument (connu à la compilation !).

\bigskip

{\footnotesize Le sous-typage est parfois considéré comme une forme de
  polymorphisme, voire la coercion (\textit{cast}).  Les limites de
  ces notions sont floues.\par}

\end{frame}
%
\begin{frame}
\frametitle{Tâches d'un système de typage}

\itempoint \textbf{Vérification} de type : vérifier qu'une expression
\emph{annotée} a bien le type prétendu par les annotations.

\bigskip

\itempoint \textbf{Inférence} (« reconstruction » / « assignation »)
de type : calculer le type de l'expression \alert{en l'absence
  d'annotations} (ou avec annotations partielles).

\medskip

Algorithme important : \textbf{Hindley-Milner} (inférence de type dans
les langages fonctionnels à polymorphisme paramétrique).  Utilisé
dans OCaml, Haskell, etc.

\medskip

\textbf{N.B.} Dans un système de typage trop complexe, l'inférence
(voire la vérification !) \textcolor{orange}{peut devenir indécidable}
(notamment si types de première classe / dépendant de valeurs
arbitraires à l'exécution).

\bigskip

{\footnotesize Rarement utile en informatique mais essentiel en
  logique ($\cong$ recherche de preuves) :\par}

\itempoint \textbf{Habitation} de type : trouver un \alert{terme}
(=expression) ayant un type donné.

{\footnotesize P.ex. : y a-t-il un terme de type
  $(\forall\mathtt{p},\mathtt{q})\; ((\mathtt{p} \to \mathtt{q}) \to
  \mathtt{p}) \to \mathtt{p}$ ?\par}

{\footnotesize \textbf{N.B.} Dans les langages usuels, \alert{tous}
  les types sont habités par une boucle infinie (« \texttt{while true
    do pass done} » ou « \texttt{let rec f () = f () in f ()} »),
  \alert{même} le type vide.\par}

\end{frame}
%
\begin{frame}
\frametitle{Utilisations du typage au-delà des valeurs stockées}

\itempoint Annotation des \textbf{exceptions soulevables} (fréquent,
p.ex. Java).

\bigskip

\itempoint Annotation de la \textbf{mutabilité} par le type.
P.ex. $\texttt{Nat}$ = type d'un entier (immuable) mais
$\texttt{Ref~Nat}$ = type d'une \alert{référence} vers un entier
(mutable).

\bigskip

\itempoint Annotation des \textbf{effets de bord} par le type.  P.ex.,
en Haskell, $\texttt{Char}$ = caractère = fonction de zéro argument
renvoyant un caractère (fonction pure : toujours le même retour), mais
$\texttt{IO~Char}$ = \alert{action} avec effets de bord renvoyant un
caractère ($\texttt{IO}$ est une famille de type appelé « monade »
I/O).

\bigskip

\itempoint Typage \textbf{linéaire} (forme de typage
« sous-structurel ») ou types à unicité : assure qu'une valeur est
utilisée \alert{une et une seule fois} dans un calcul (ni duplication
ni perte sauf manœuvre spéciale).

\smallskip

Permet d'optimiser la gestion de la mémoire (Rust) et/ou d'annoter les
effets de bord (Clean).

\end{frame}
%
\begin{frame}
\frametitle{Quelques exemples (1)}

Les langages impératifs \emph{tendent} à avoir des systèmes de typage
moins complexes que les langages fonctionnels.

\medskip

{\footnotesize Éviter les termes de typage « faible » et « fort », qui
  veulent tout (ou rien) dire.\par}

\medskip

\itempoint Assembleur (langage machine) : aucun typage (tout est
donnée binaire).

{\footnotesize Idem : machine de Turing, fonctions générales
  récursives (tout est entier), $\lambda$-calcul non typé (tout est
  fonction).\par}

\medskip

\itempoint C : annoté, vérifié à la compilation (\emph{aucune}
vérification à l'exécution), moyennement complexe et contournable
(pointeurs génériques \texttt{void*}).

\medskip

\itempoint Python, JavaScript, Perl, Scheme, etc. : typage vérifié à
l'exécution, superficiel.  Suffisant pour éviter les comportements
indéfinis.

\medskip

\itempoint Java : annoté, mixte compilation/exécution (double
vérification), initialement superficiel (listes non typées), puis
introduction de « génériques » ($\rightarrow$ polymorphisme
paramétrique) avec Java 5, puis diverses sortes d'inférence.

\medskip

\itempoint Rust : interaction avec la gestion de la mémoire
($\approx$ typage linéaire/affine).

\end{frame}
%
\begin{frame}
\frametitle{Quelques exemples (2)}

{\footnotesize Qqs exemples de langages, généralement fonctionnels,
  ayant un système de typage (très) complexe, mélangeant plusieurs
  fonctionnalités évoquées (interactions parfois délicates !) :\par}

\medskip

\itempoint OCaml : inférence de type à la H-M, types récursifs,
polymor\textsuperscript{sme} paramétrique.

\smallskip

Système de « modules » (« signatures » $\approx$ interfaces
abstraites, « foncteurs » entre signatures…), comparable aux
« classes » des langages orientés objet.

\medskip

\itempoint Haskell : beaucoup de similarité avec OCaml :

\smallskip

+ polymorphisme ad hoc : « classes de type », comparable aux
« modules » de OCaml, aux « classes » des langages orientés objet.

\smallskip

+ purement fonctionnel (toutes les fonctions sont « pures ») : les
effets de bord sont enrobés dans des « monades ».

\medskip

\itempoint Mercury (langage de type fonctionnel+logique, inspiré de
Haskell+Prolog) : typage comparable à Haskell ; + sous-typage,
linéarité.

\medskip

\itempoint Idris : langage fonctionnel + assistant de preuve.

\end{frame}
%
\begin{frame}
\frametitle{Quelques curiosités}

\itempoint En F\#, les unités de mesure sont intégrées au système de
typage, qui peut donc vérifier l'homogénéité physique.

\bigskip

\itempoint En Rust, le système de typage évite non seulement les
fuites de mémoire mais aussi certains problèmes de concurrence
(absence de \textit{race condition} sur les données).

\bigskip

\itempoint Certains langages spécialisés assurent, éventuellement en
lien avec leur système de typage, des propriétés diverses de leurs
programmes : terminaison garantie en temps polynomial (Bellantoni \&
Cook 1992), voire constant (langage Usuba), validation XML (langage
CDuce), absence de fuite d'information (langage Flow Caml), etc.

\bigskip

\itempoint En Idris, le système de typage est aussi un système de
preuve (cf. Curry-Howard après) et permet de certifier des invariants
quelconques d'un programme (p.ex., correction d'un algorithme de tri).

\end{frame}
%
\begin{frame}
\label{typing-and-termination}
\frametitle{Typage et terminaison}

Un système de typage \alert{peut} garantir que \textbf{tout programme
  bien typé termine}.

\bigskip

Ce \alert{n'est pas le cas} pour les langages de programmation usuels
(en OCaml, Haskell, etc., un programme bien typé peut faire une boucle
infinie).

\bigskip

Si on veut que le système de typage soit décidable, ceci met forcément
des \alert{limites sur l'expressivité} du langage
($\leftarrow$ indécidabilité du problème de l'arrêt).

\bigskip

Notamment, le langage ne permettra pas d'écrire son propre
interpréteur (même argument « diagonal » que pour les fonctions p.r.
par thm. de réc\textsuperscript{sion} de Kleene).

\bigskip

Notamment aussi : pas de boucle illimitée, pas d'appels récursifs
\emph{non contraints}, pas de type tel que $\mathtt{t} \,\cong\,
(\mathtt{t}\rightarrow \mathtt{t})$.  Aucun terme de type vide !

\bigskip

Néanmoins, le langage peut être \alert{beaucoup plus puissant} que les
fonctions p.r.

\bigskip

Exemple de tel langage : Coq (maintenant appelé Rocq).

\end{frame}
%
\begin{frame}
\frametitle{La correspondance de Curry-Howard (avant-goût)}

Idée générale : établir une analogie, voire une correspondance précise
entre
\begin{itemize}
\item \textbf{types} et \textbf{termes} (=programmes) de ces types
  dans un système de typage,
\item \textbf{propositions} et \textbf{preuves} de ces propositions
  dans un système logique.
\end{itemize}

\smallskip

P.ex., la preuve évidente de l'affirmation $(A \land B) \Rightarrow A$
(« si $A$ et $B$ sont vraies alors $A$ est vraie ») correspond à la
« première projection » de type $a \times b \rightarrow a$.

\smallskip

{\footnotesize Ceci « explique » que le type d'un terme tel que
  $\lambda xy.x$, soit $u \to v \to u$, soit une vérité logique, en
  l'occurrence $U \Rightarrow V \Rightarrow U$ (« si $U$ est vrai
  alors, si $V$ est vrai alors $U$ est vrai »).\par}

\bigskip

Beaucoup de variations sur cette correspondance, mais il y a des
restrictions :
\begin{itemize}
\item le langage informatique doit garantir la terminaison (pas
  d'appels récursifs !), sinon cela reviendrait à permettre les
  preuves circulaires,
\item la logique concernée est plutôt « intuitionniste » (sans tiers
  exclu),
\item diverses subtilités notamment dans la correspondance entre types
  sommes paramétriques (types $\Sigma$) et $\exists$ côté logique.
\end{itemize}

\end{frame}
%
\begin{frame}
\frametitle{La correspondance de Curry-Howard (illustration graphique)}

\begin{center}
\begin{tabular}{c@{\hskip 30bp}c}
\includegraphics[height=150bp]{images/sean-connery-in-name-of-the-rose.jpg}
&\includegraphics[height=150bp]{images/sean-connery-in-zardoz.jpg}
\\
Preuves mathématiques&Programmes informatiques
\\
{\footnotesize (contemplatives ?)}&{\footnotesize (dynamiques ?)}
\end{tabular}

\bigskip

On a du mal à le croire, mais c'est la même chose !

{\footnotesize (N.B. : ne pas prendre ce résumé simpliste trop au sérieux.)\par}
\end{center}

\end{frame}
%
\begin{frame}
\frametitle{La correspondance de Curry-Howard (premier aperçu)}

{\footnotesize\textcolor{gray}{Divulgâchis pour la suite !}\par}

\smallskip

La correspondance de Curry-Howard fait notamment correspondre :
\begin{itemize}
\item type \textbf{fonction} $A \to B$ avec \textbf{implication}
  logique $A \Rightarrow B$,
\item \textbf{application} d'une fonction ($A \to B$ à un argument de
  type $A$) avec \textbf{\textit{modus ponens}} (si $A \Rightarrow B$
  et $A$, alors $B$),
\item \textbf{abstraction} (= création d'une fonction) avec ouverture
  d'une \textbf{hypothèse} (« supposons $A$ : alors (…) donc $B$ ;
  ceci prouve $A \Rightarrow B$ »),
\item type \textbf{produit} (= couples) $A \times B$ avec
  \textbf{conjonction} (« et ») logique $A \land B$,
\item type \textbf{somme} $A + B$ avec \textbf{disjonction} (« ou »)
  logique $A \lor B$,
\item types \textbf{trivial} ($\texttt{Unit}$) et \textbf{vide} avec
  \textbf{vrai} $\top$ et \textbf{faux} $\bot$ logiques,
\item mais pour la \textbf{négation}…  c'est plus délicat.
\end{itemize}

\bigskip

{\footnotesize Les détails demandent un système de typage et un
  système logique précisément définis.\par}

\end{frame}
%
\begin{frame}
\frametitle{Paradoxe de Curry (interlude distrayant)}

{\footnotesize Variante plus sophistiquée de « cette phrase est
  fausse ».  C'est un exemple de preuve circulaire
  ($\leftrightarrow$ combinateur $\mathsf{Y}$ de Curry par la
  correspondance de C-H), invalide en logique.\par}

\medskip

Je tiens l'affirmation : « \textcolor{teal}{si j'ai raison, alors je
  suis un grand génie} ».

\begin{itemize}
\item clairement, si j'ai raison, je suis un grand génie ;
\item mais c'était justement mon affirmation : donc j'ai raison ;
\item donc je suis un grand génie. \qedsymbol
\end{itemize}

{\footnotesize Remplacer « j'ai raison » par « cette phrase est
  vraie » (voire utiliser l'astuce de Quine pour éviter « cette
  phrase ») et « je suis un grand génie » par absolument n'importe
  quoi.\par}

\bigskip

\itempoint Ce qui est correct : \alert{si on peut construire} un
énoncé $A$ tel que $A \Leftrightarrow (A \Rightarrow B)$ (ici $A$ est
l'affirmation tenue et $B$ est la conclusion voulue) alors $B$ vaut :
\[
(A \Leftrightarrow (A \Rightarrow B)) \, \Rightarrow \, B
\]

{\footnotesize (Preuve correcte : supposons $A$ : par hypothèse, ceci
  signifie $A \Rightarrow B$ ; on a donc $B$ ; tout ceci prouve $A
  \Rightarrow B$.  Bref on a $A \Rightarrow B$.  Mais par hypothèse
  c'est $A$.  Donc $A$ vaut.  Donc $B$ aussi.)\par}

\medskip

\itempoint Ce qui est \textcolor{orange}{fallacieux} : la supposition
tacite qu'on peut construire un tel $A$.

{\footnotesize L'astuce de Quine permet de faire une auto-référence
  sur la syntaxe, pas sur la vérité.\par}

\end{frame}
%
\begin{frame}
\frametitle{Théories des types pour la logique (très bref aperçu)}

\itempoint Langages spécialisés pour servir à la fois à l'écriture de
programmes informatiques et de preuves mathématiques ; ils peuvent
être soit sous forme de systèmes abstraits, soit sous forme
d'implémentation informatique (« assistants de preuve » : Coq/Rocq, Agda,
Lean…).

\bigskip

\itempoint Dans tous les cas il s'agit de langages assurant la
terminaison des programmes (cf. transp. \ref{typing-and-termination}),
ou, puisqu'il s'agit de variantes du $\lambda$-calcul, la
\emph{normalisation forte}.  Le type vide doit être inhabité !

\bigskip

Quelques grandes familles (chacune avec énormément de variantes) :
\begin{itemize}
\item théories des types à la Martin-Löf (« MLTT ») : suit jusqu'au
  bout la correspondance de Curry-Howard (\emph{aucune} distinction
  entre types et propositions ; {\footnotesize pas de $\exists$ mais
    plutôt un $\Sigma$}), $\rightarrow$ Agda ;
\item variantes du « calcul des constructions » (« CoC »),
  $\rightarrow$ Coq ;
\item théorie homotopique des types (« HoTT ») : « égalité = isomorphisme ».
\end{itemize}

\end{frame}
%
\begin{frame}
\frametitle{Le $\lambda$-cube de Barendregt (très bref aperçu)}

Il s'agit d'un ensemble de $8 = 2^3$ théories des types pour la
logique : la plus faible est le « $\lambda$-calcul simplement typé »
($\lambda_\rightarrow$, décrit plus loin), la plus forte le « calcul
des constructions » (« CoC »).

\medskip

Chaque théorie du cube est caractérisée par l'absence ou la présence
de chacune des $3$ fonctionnalités suivantes ($\lambda_\rightarrow$
n'a aucune des trois, CoC a les trois) :

\begin{itemize}
\item \textcolor{blue}{termes pouvant dépendre de types}, ou
  polymorphisme paramétrique « explicitement quantifié » (p.ex. $\prod
  t.\; (t \rightarrow t)$) ;
\item \textcolor{blue}{types pouvant dépendre de types}, ou
  familles de types ;
\item \textcolor{blue}{types pouvant dépendre de termes}, ou « types
  dépendants », correspondant côté logique à la quantification sur les
  individus.
\end{itemize}

\medskip

{\footnotesize On peut ensuite encore ajouter des fonctionnalités :
  par exemple, Coq/Rocq est basé sur le « calcul des constructions
  inductives » qui ajoute au calcul des constructions des mécanismes
  systématiques pour former des types (positivement !) récursifs.\par}

\end{frame}
%
\section{Le \texorpdfstring{$\lambda$}{lambda}-calcul simplement typé}
\begin{frame}
\frametitle{Le $\lambda$-calcul simplement typé : description sommaire}

\itempoint Le $\lambda$-calcul simplement typé (=: $\lambda$CST ou
$\lambda_\rightarrow$) est une \alert{variante typée} du
$\lambda$-calcul, assurant la propriété de terminaison (normalisation
forte).

\medskip

\itempoint Il a une seule opération sur les types, le type
\textbf{fonction} : donnés deux types $\sigma,\tau$, on a un type
$\sigma\to\tau$ pour les fonctions de l'un vers l'autre.

\medskip

\itempoint L'application et l'abstraction doivent respecter le
typage :
\begin{itemize}
\item si $P$ a pour type $\sigma\to\tau$ et $Q$ a type $\sigma$ alors $PQ$
  a pour type $\tau$,
\item si $E$ a pour type $\tau$ en faisant intervenir une variable $v$
  libre de type $\sigma$ alors $\lambda(v:\sigma).E$ {\footnotesize
    (« fonction prenant $v$ de type $\sigma$ et renvoyant $E$ »)} a
  pour type $\sigma\to\tau$.
\end{itemize}

\medskip

\itempoint Typage \textbf{annoté} (=« à la Church ») : on écrit
$\lambda(v:\sigma).E$ pas juste $\lambda v.E$.

\medskip

\itempoint On notera « $x_1:\sigma_1,\ldots,x_k:\sigma_k \vdash
E:\tau$ » pour dire « $E$ est bien typé avec pour type $\tau$ lorsque
$x_1,\ldots,x_k$ sont des variables libres de types
$\sigma_1,\ldots,\sigma_k$ :
\begin{itemize}
\item $x_1:\sigma_1,\ldots,x_k:\sigma_k$ s'appelle un
  \textbf{contexte} de typage (souvent $\Gamma$),
\item $\Gamma \vdash E:\tau$ s'appelle un \textbf{jugement} de typage.
\end{itemize}

\end{frame}
%
\begin{frame}
\frametitle{Exemples de termes et de typages}

{\footnotesize On donnera les règles précises plus tard, commençons
  par quelques exemples.\par}

\begin{itemize}
\item $f:\alpha\to\beta,\; x:\alpha \;\vdash\; fx:\beta$\newline Lire : « dans le
  contexte où $f$ est une variable de type $\alpha\to\beta$ et $x$ une
  variable de type $\alpha$, alors $fx$ est de type $\beta$ ».
\item $f:\beta\to\alpha\to\gamma,\; x:\alpha,\; y:\beta \;\vdash\; fyx
  : \gamma$\newline {\footnotesize (Parenthéser $\beta\to\alpha\to\gamma$
  comme $\beta\to(\alpha\to\gamma)$ et $fyx$ comme $(fy)x$.)}
\item $f:\beta\to\alpha\to\gamma,\; x:\alpha \;\vdash\;
  \lambda(y:\beta).fyx \;:\; \beta\to\gamma$\newline {\footnotesize
    Comprendre $\lambda(y:\beta).fyx$ comme « fonction prenant $y$ de
    type $\beta$ et renvoyant $fyx$ ».}
\item $x:\alpha \;\vdash\;
  \lambda(f:\alpha\to\beta).fx\;:\;(\alpha\to\beta)\to\beta$\newline « Si
  $x$ est de type $\alpha$ alors le terme
  $\lambda(f:\alpha\to\beta).fx$ est de type
  $(\alpha\to\beta)\to\beta$. »
\item $\vdash\; \lambda(x:\alpha).
  \lambda(f:\alpha\to\beta).fx\;:\;\alpha\to(\alpha\to\beta)\to\beta$\newline « Le
  terme $\lambda(x:\alpha).\lambda(f:\alpha\to\beta).fx$ a type
  $\alpha\to(\alpha\to\beta)\to\beta$ dans le contexte
  vide. »\newline {\footnotesize (Parenthéser
    $\alpha\to(\alpha\to\beta)\to\beta$ comme
    $\alpha\to((\alpha\to\beta)\to\beta)$.)}
\end{itemize}

\end{frame}
%
\begin{frame}
\frametitle{Types et prétermes}

\itempoint Un \textbf{type} du $\lambda$CST est (inductivement) :
\begin{itemize}
\item une \textbf{variable de type} ($\alpha$, $\beta$, $\gamma$... en nombre illimité),
\item un \textbf{type fonction} $(\sigma\to\tau)$ où $\sigma$ et $\tau$
  sont deux types.
\end{itemize}

\bigskip

\itempoint Un \textbf{préterme} du $\lambda$CST est (inductivement) :
\begin{itemize}
\item une \textbf{variable de terme} ($a$, $b$, $c$... en nombre illimité),
\item une \textbf{application} $(PQ)$ où $P$ et $Q$ sont deux prétermes,
\item une \textbf{abstraction} $\lambda(v:\sigma).E$ avec $v$
  variable, $\sigma$ type et $E$ préterme.
\end{itemize}

\medskip

\itempoint Conventions d'écriture :
\begin{itemize}
\item « $\rho\to\sigma\to\tau$ » signifie
  « $(\rho\to(\sigma\to\tau))$ » ; « $xyz$ » signifie
  « $((xy)z)$ » ;
\item {\footnotesize on note « $\lambda (x:\sigma,t:\tau).E$ » ou
  « $\lambda (x:\sigma)(t:\tau).E$ » pour « $\lambda
  (x:\sigma). \lambda (t:\tau). E$ » ;}
\item l'abstraction est moins prioritaire que l'application ;
\item on considère les termes à renommage près des variables liées
  {\footnotesize ($\alpha$-conversion)}.
\end{itemize}

\end{frame}
%
\begin{frame}
\frametitle{Règles de typage}

\itempoint Un \textbf{typage} est la donnée d'un préterme $M$ et d'un
type $\sigma$.  On note « $M:\sigma$ ».

\medskip

\itempoint Un \textbf{contexte} est un ensemble fini $\Gamma$ de
typages $x_i:\sigma_i$ où $x_1,\ldots,x_k$ sont des \alert{variables}
de terme \alert{distinctes}.  On le note
« $x_1:\sigma_1,\ldots,x_k:\sigma_k$ ».

\medskip

\itempoint Un \textbf{jugement} de typage est la donnée d'un contexte
$\Gamma$ et d'un typage $E:\tau$, sujet aux règles ci-dessous.  On
note $\Gamma \vdash E:\tau$ (ou juste $\vdash E:\tau$ si $\Gamma =
\varnothing$), et on dit que $E$ est un « \textbf{terme} (= bien typé)
de type $\tau$ dans le contexte $\Gamma$ ».

\medskip

\textbf{Règles} de typage du $\lambda$CST : {\footnotesize (quels que
  soient $\Gamma,x,\sigma,\ldots$,)}
\begin{itemize}
\item \textcolor{olive}{(« variable »)} si $(x:\sigma) \in \Gamma$
  alors $\Gamma \vdash x:\sigma$ ;
\item \textcolor{olive}{(« application »)} si $\Gamma \vdash
  P:\sigma\to\tau$ et $\Gamma \vdash Q:\sigma$ alors $\Gamma \vdash
  PQ:\tau$ ;
\item \textcolor{olive}{(« abstraction »)} si $\Gamma, v:\sigma \vdash
  E:\tau$ alors $\Gamma \vdash \lambda(v:\sigma).E : \sigma\to\tau$.
\end{itemize}

\smallskip

{\footnotesize (Comprendre : l'ensemble des jugements de typage est
  l'ensemble engendré par les règles ci-dessus, i.e., le plus petit
  ensemble qui les respecte.)\par}

\medskip

\itempoint Une \textbf{dérivation} d'un jugement est un arbre
(d'instances de) règles qui aboutit au jugement considéré.

\end{frame}
%
\begin{frame}
\frametitle{Représentation des règles de typage}

Les trois règles de typage du $\lambda$CST :

\[
\inferrule*[Left={Var}]{ }{\Gamma,x:\sigma \vdash x:\sigma}
\]

\[
\inferrule*[Left={App},sep=4em]{\Gamma \vdash P:\sigma\to\tau\\ \Gamma
  \vdash Q:\sigma}{\Gamma \vdash PQ:\tau}
\]

\[
\inferrule*[Left={Abs}]{\Gamma, v:\sigma \vdash E:\tau}{\Gamma
  \vdash \lambda(v:\sigma).E : \sigma\to\tau}
\]

\bigskip

\itempoint Chaque « fraction » indique que le jugement écrit
\alert{en-dessous} découle par la règle inscrite \alert{à côté} à
partir des hypothèses portées \alert{au-dessus}.

\end{frame}
%
\begin{frame}
\label{example-typing-derivation}
\frametitle{Exemple de dérivation}

Montrons le jugement selon lequel
$\lambda(f:\beta\to\alpha\to\gamma). \lambda(x:\alpha). \lambda(y:\beta). fyx$
est un terme de type $(\beta\to\alpha\to\gamma) \to
(\alpha\to\beta\to\gamma)$ dans le contexte vide :

\vskip-2ex

{\footnotesize
\[
\inferrule*[Left={Abs}]{
\inferrule*[Left={Abs}]{
\inferrule*[Left={Abs}]{
\inferrule*[Left={App}]{
\inferrule*[Left={App}]{
\inferrule*[Left={Var}]{ }{\mbox{$\begin{array}{cc}\scriptscriptstyle f:\beta\to\alpha\to\gamma,\\ \scriptscriptstyle x:\alpha,y:\beta\\ \end{array}$} \vdash f : \beta\to\alpha\to\gamma}
\\
\inferrule*[Right={Var}]{ }{\mbox{$\begin{array}{cc}\scriptscriptstyle f:\beta\to\alpha\to\gamma,\\ \scriptscriptstyle x:\alpha,y:\beta\\ \end{array}$} \vdash y : \beta}
}{\mbox{$\begin{array}{cc}\scriptscriptstyle f:\beta\to\alpha\to\gamma,\\ \scriptscriptstyle x:\alpha,y:\beta\\ \end{array}$} \vdash fy : \alpha\to\gamma}
\\
\inferrule*[Right={Var}]{ }{\mbox{$\begin{array}{cc}\scriptscriptstyle f:\beta\to\alpha\to\gamma,\\ \scriptscriptstyle x:\alpha,y:\beta\\ \end{array}$} \vdash x : \alpha}
}{f:\beta\to\alpha\to\gamma, x:\alpha, y:\beta \vdash fyx : \gamma}
}{f:\beta\to\alpha\to\gamma, x:\alpha \vdash \lambda(y:\beta). fyx : \beta\to\gamma}
}{f:\beta\to\alpha\to\gamma \vdash \lambda(x:\alpha). \lambda(y:\beta). fyx : \alpha\to\beta\to\gamma}
}{\vdash \lambda(f:\beta\to\alpha\to\gamma). \lambda(x:\alpha). \lambda(y:\beta). fyx : (\beta\to\alpha\to\gamma) \to (\alpha\to\beta\to\gamma)}
\]
\par}

Chaque barre horizontale justifie le jugement écrit \alert{en-dessous}
par la règle inscrite \alert{à côté} à partir des hypothèses portées
\alert{au-dessus}.

\medskip

Ceci est typographiquement abominable et hautement redondant, ce qui
explique qu'on écrive rarement de tels arbres complètement.

\end{frame}
%
\begin{frame}
\frametitle{Propriétés du typage}

Les propriétés suivantes sont faciles mais utiles :

\medskip

\itempoint \textbf{Affaiblissement :} si $\Gamma \subseteq \Gamma'$
sont deux contextes et $\Gamma \vdash M:\sigma$ alors $\Gamma' \vdash
M:\sigma$ aussi.

\smallskip

{\footnotesize On pouvait présenter les règles en limitant la règle
  « variable » à « $x:\sigma \vdash x:\sigma$ » et en prenant
  l'affaiblissement comme règle.  C'est peut-être préférable.\par}

\medskip

\itempoint \textbf{Duplication :} si $\Gamma, x:\rho, x':\rho \vdash
M:\sigma$ alors $\Gamma, x:\rho \vdash M[x'\backslash x]:\sigma$ aussi
(où $M[x'\backslash x]$ désigne la substitution correcte de $x'$
par $x$).

\smallskip

{\footnotesize \emph{Remarque :} Le typage linéaire supprime notamment
  l'affaiblissement et la duplication.\par}

\medskip

\itempoint \textbf{Variables libres :} si
$x_1:\sigma_1,\ldots,x_k:\sigma_k \vdash E:\tau$ alors toute variable
libre de $E$ est une des $x_1,\ldots,x_k$.

\medskip

\itempoint \textbf{Variables inutiles :} si $\Gamma \vdash E:\tau$
alors $\Gamma|_{\operatorname{free}(E)} \vdash E:\tau$ où
$\Gamma|_{\operatorname{free}(E)}$ est la partie de $\Gamma$
concernant les variables libres de $E$.

\medskip

\itempoint \textbf{Renommage :} si $\Gamma \vdash E:\tau$ alors ceci
vaut encore après tout renommage des variables libres.

\end{frame}
%
\begin{frame}
\label{constructing-the-typing-derivation}
\frametitle{Construction de la dérivation}

La dérivation du jugement de typage $\Gamma \vdash M : \sigma$ d'un
(pré)terme $M$ du $\lambda$CST se construit systématiquement et
évidemment à partir de $M$ (\alert{aucun choix} à faire !) :
\begin{itemize}
\item si $M = x$ est une variable, alors un $x:\sigma$ doit être dans
  le contexte $\Gamma$ (sinon : échouer), et la dérivation consiste en
  la règle « variable » ;
\item si $M = PQ$ est une application, construire des dérivations de
  jugements $\Gamma \vdash P : \rho$ et $\Gamma \vdash Q : \sigma$, on
  doit avoir $\rho = (\sigma\to\tau)$ (sinon : échouer), et la
  dérivation finit par la règle « application » et donne $\Gamma
  \vdash M : \tau$ ;
\item si $M = \lambda(v:\sigma).E$ est une abstraction, construire une
  dérivation de jugement $\Gamma' \vdash E : \tau$ dans $\Gamma' :=
  \Gamma \cup \{v:\sigma\}$ {\footnotesize (quitte à renommer $v$)},
  et la dérivation finit par la règle « abstraction » et donne $\Gamma
  \vdash M : \sigma\to\tau$.
\end{itemize}

\bigskip

\itempoint \textbf{Donc :} aussi bien la vérification de type
(vérifier $\Gamma \vdash M : \sigma$) que l'assignation de type
(trouver $\sigma$ à partir de $\Gamma,M$) sont (très facilement)
\alert{décidables} dans le $\lambda$CST.  De plus, $\sigma$ est
\alert{unique} (donnés $\Gamma$ et $M$).

\end{frame}
%
\begin{frame}
\frametitle{Problèmes faciles et moins faciles}

\itempoint Facile (transp. précédent) : donné un (pré)terme $M$,
trouver son (seul possible) type $\sigma$ est facile (notamment :
vérifier que $M$ est un terme = bien typé, ou vérifier un jugement de
type).

\smallskip

\textcolor{blue}{Le terme donne directement l'arbre de dérivation.}

\bigskip

\itempoint Facile aussi : réciproquement, donné un arbre de dérivation
où on a effacé les termes pour ne garder que les types, retrouver un
terme (transp. suivant).

\smallskip

\textcolor{blue}{L'arbre de dérivation redonne directement le terme.}

\bigskip

\itempoint Moins facile : donné un terme « désannoté », i.e., un terme
du $\lambda$-calcul non typé, p.ex. $\lambda xyz.xz(yz)$, savoir s'il
correspond à un terme (typable) du $\lambda$CST.

\smallskip

\textcolor{brown}{$\rightarrow$ Algorithme de Hindley-Milner
  (\emph{inférence} de type).}

\bigskip

\itempoint Moins facile : donné un type, savoir s'il existe un terme
ayant ce type.

\smallskip

\textcolor{brown}{$\rightarrow$ Décidabilité du calcul propositionnel
  intuitionniste.}

\end{frame}
%
\begin{frame}
\label{reconstructing-typed-term}
\frametitle{Reconstruction du terme typé}

Peut-on retrouver les termes manquants dans un arbre de dérivation
dont on n'a donné que les types et les règles appliquées ?

{\footnotesize
\[
\inferrule*[Left={Abs}]{
\inferrule*[Left={Abs}]{
\inferrule*[Left={Abs}]{
\inferrule*[Left={App}]{
\inferrule*[Left={App}]{
\inferrule*[Left={Var}]{ }{\mbox{$\begin{array}{cc}\scriptscriptstyle \textbf{?}:\beta\to\alpha\to\gamma,\\ \scriptscriptstyle \textbf{?}:\alpha,\textbf{?}:\beta\\ \end{array}$} \vdash \textbf{?} : \beta\to\alpha\to\gamma}
\\
\inferrule*[Right={Var}]{ }{\mbox{$\begin{array}{cc}\scriptscriptstyle \textbf{?}:\beta\to\alpha\to\gamma,\\ \scriptscriptstyle \textbf{?}:\alpha,\textbf{?}:\beta\\ \end{array}$} \vdash \textbf{?} : \beta}
}{\mbox{$\begin{array}{cc}\scriptscriptstyle \textbf{?}:\beta\to\alpha\to\gamma,\\ \scriptscriptstyle \textbf{?}:\alpha,\textbf{?}:\beta\\ \end{array}$} \vdash \textbf{?} : \alpha\to\gamma}
\\
\inferrule*[Right={Var}]{ }{\mbox{$\begin{array}{cc}\scriptscriptstyle \textbf{?}:\beta\to\alpha\to\gamma,\\ \scriptscriptstyle \textbf{?}:\alpha,\textbf{?}:\beta\\ \end{array}$} \vdash \textbf{?} : \alpha}
}{\textbf{?}:\beta\to\alpha\to\gamma, \textbf{?}:\alpha, \textbf{?}:\beta \vdash \textbf{?} : \gamma}
}{\textbf{?}:\beta\to\alpha\to\gamma, \textbf{?}:\alpha \vdash \textbf{?} : \beta\to\gamma}
}{\textbf{?}:\beta\to\alpha\to\gamma \vdash \textbf{?} : \alpha\to\beta\to\gamma}
}{\vdash \textbf{?} : (\beta\to\alpha\to\gamma) \to (\alpha\to\beta\to\gamma)}
\]
\par}

\textbf{Oui} (aux noms des variables près) à condition, en cas de
types identiques dans le contexte, d'identifier l'élément utilisé pour
chaque règle « variable ».

\bigskip

\itempoint\textcolor{blue}{Le terme typé représente exactement l'arbre
  de dérivation} du jugement le concernant.

\end{frame}
%
\begin{frame}
\label{substitution-lemma}
\frametitle{Lemme de substitution}

\itempoint\textbf{Lemme :} Si $\Gamma, v:\sigma \vdash E : \tau$ et
$\Gamma \vdash T : \sigma$ alors $\Gamma \vdash E[v\backslash T] :
\tau$, où $E[v\backslash T]$ désigne le terme obtenu par substitution
correcte de $v$ par $T$ dans $E$.

\bigskip

\underline{Esquisse de preuve :} reprendre l'arbre de dérivation de
$\Gamma, v:\sigma \vdash E : \tau$ en supprimant $v:\sigma$ du
contexte et en substituant $v$ par $T$ partout à droite du ‘$\vdash$’.
Les règles s'appliquent à l'identique, sauf la règle « variable »
introduisant $v:\sigma$, qui est remplacée par l'arbre de dérivation
de $\Gamma \vdash T : \sigma$.\qed

\medskip

\centerline{*}

\medskip

Remarquer la similarité entre
\[
\inferrule*[left=App]{
\inferrule*[left=Abs]{\Gamma, v:\sigma \vdash E : \tau}{
\Gamma \vdash \lambda(v:\sigma).E : \sigma\to\tau}
\\
\Gamma \vdash T : \sigma
}{\Gamma \vdash (\lambda(v:\sigma).E)T : \tau}
\]
et
\[
\inferrule*[left=Subs]{\Gamma, v:\sigma \vdash E : \tau
\\
\Gamma \vdash T : \sigma
}{\Gamma \vdash E[v\backslash T] : \tau}
\]

\end{frame}
%
\begin{frame}
\label{typing-vs-beta-reduction}
\frametitle{Typage et $\beta$-réduction}

On rappelle que la « $\beta$-réduction » désigne le remplacement en
sous-expression d'un « redex » $(\lambda (v:\sigma). E)T$ par son
« réduit » $E[v\backslash T]$.

\medskip

On note $X \rightarrow_\beta X'$ pour une $\beta$-réduction, et $X
\twoheadrightarrow_\beta X'$ pour une suite de $\beta$-réductions.

\bigskip

\itempoint D'après le lemme de substitution, si $\Gamma \vdash
(\lambda (v:\sigma). E)T : \tau$ alors on a aussi $\Gamma \vdash
E[v\backslash T] : \tau$.

\bigskip

\itempoint Toujours d'après ce lemme, si $X \twoheadrightarrow_\beta
X'$ et $\Gamma \vdash X:\tau$ alors $\Gamma \vdash X':\tau$.

\bigskip

\textcolor{blue}{\textbf{Moralité :}} le typage est compatible avec
l'évaluation du $\lambda$-calcul (un terme bien typé reste bien typé
quand on effectue la $\beta$-réduction et le type ne change pas).

\end{frame}
%
\begin{frame}
\frametitle{Normalisation forte}

\itempoint\textbf{Théorème} (Turing, Tait, Girard) : le
$\lambda$-calcul simplement typé est \textbf{fortement normalisant},
c'est-à-dire que toute suite de $\beta$-réductions sur un terme
(bien typé !) termine.

\bigskip

{\footnotesize

\textcolor{blue}{Idée cruciale de la preuve :} une induction directe
échoue ($P,Q$ fortement normalisables $\not\Rightarrow$ $PQ$ fortement
normalisable).  À la place, on introduit une notion \emph{plus forte}
permettant l'induction :

\smallskip

\itempoint\textbf{Définition} technique : un terme $P$ de type $\rho$
est dit « fortement calculable » lorsque :
\begin{itemize}
\item si $\rho = \alpha$ est une variable de type (type
  « atomique ») : $P$ est fortement normalisable,
\item si $\rho = (\sigma\to\tau)$ est un type fonction : pour
  \alert{tout} $Q$ de type $\sigma$ fortement calculable, $PQ$ est
  fortement calculable.
\end{itemize}

\smallskip

On prouve alors successivement :
\begin{itemize}
\item tout terme fortement calculable est fortement normalisable
  (induction facile sur le type),
\item si $x_1:\sigma_1,\ldots,x_k:\sigma_k \vdash E:\tau$ et si
  $P_1,\ldots,P_k$ sont de types $\sigma_1,\ldots,\sigma_k$, fortement
  calculables et n'impliquant pas $x_1,\ldots,x_k$, alors la
  substitution de $P_i$ pour $x_i$ dans $E$ est fortement calculable
  (preuve par induction sur la dérivation du jugement, i.e., induction
  sur $E$ : le cas délicat est l'abstraction).
\end{itemize}

\par}

\end{frame}
%
\section{Le calcul propositionnel intuitionniste}
\begin{frame}
\frametitle{Le calcul propositionnel : présentation sommaire}

\itempoint Le \textbf{calcul propositionnel} est une forme de logique
qui ne parle que de « propositions » (énoncés logiques) : pas de
variables d'individus (p.ex. entiers naturels, ensembles…), \alert{pas
  de quantification} (pas de $\forall,\exists$).

\medskip

\itempoint Les \textbf{connecteurs logiques} sont : $\Rightarrow$
(implication), $\land$ (conjonction logique : « et »), $\lor$
(disjonction : « ou »), $\top$ (« vrai ») et $\bot$ (« faux » ou
« absurde »).

{\footnotesize Tous sont binaires sauf $\top$, $\bot$ (nullaires).\par}

\smallskip

On peut y ajouter $\neg$ (négation logique) unaire : $\neg P$ est une
abréviation de $P \Rightarrow \bot$ (soit : « $P$ est absurde »).

\medskip

On distingue {\footnotesize (les règles précises seront décrites plus
  loin)} :
\begin{itemize}
\item la logique \textbf{classique} ou \textbf{booléenne} qui admet la
  règle du \alert{tiers exclu} (« tertium non datur ») sous une forme
  ou une autre : logique usuelle des maths ; on peut la modéliser par
  $2$ valeurs de vérité (« vrai » et « faux ») ;
\item la logique \textbf{intuitionniste}, plus \emph{faible}, qui
  n'admet pas cette règle : il n'y a pas vraiment de « valeur de
  vérité ».
\end{itemize}

\end{frame}
%
\begin{frame}
\frametitle{Le calcul propositionnel : syntaxe}

\itempoint Une \textbf{formule} du calcul propositionnel est
(inductivement) :
\begin{itemize}
\item une \textbf{variable propositionnelle} ($A$, $B$, $C$...),
\item l'application d'un connecteur : $(P\Rightarrow Q)$, $(P\land
  Q)$, $(P\lor Q)$ où $P,Q$ sont deux formules, ou encore $\top$,
  $\bot$.
\end{itemize}

\medskip

\itempoint Conventions d'écriture :
\begin{itemize}
\item $\neg P$ abrège « $(P \Rightarrow \bot)$ » ;
\item on omet certaines parenthèses : $\neg$ est prioritaire sur
  $\land$ qui est prioritaire sur $\lor$ qui est prioritaire sur
  $\Rightarrow$ {\footnotesize (\textcolor{purple}{ne pas abuser de
    l'omission des parenthèses, merci})} ;
\item $\land$ et $\lor$ associent à gauche disons (ça n'aura pas
  d'importance) ; mais $\Rightarrow$ associe à droite : $P \Rightarrow
  Q\Rightarrow R$ signifie $(P\Rightarrow(Q\Rightarrow R))$ ;
\item parfois : $P \Leftrightarrow Q$ pour $(P\Rightarrow Q) \land
  (Q\Rightarrow P)$ {\footnotesize (priorité encore plus basse ?)}.
\end{itemize}

\medskip

\itempoint Si $P_1,\ldots,P_r, Q$ sont des formules, $P_1,\ldots,P_r
\vdash Q$ s'appelle un \textbf{séquent}, à comprendre comme « sous les
hypothèses $P_1,\ldots,P_r$, on a $Q$ ».

\medskip

{\footnotesize \itempoint Notation historique : « $\supset$ » pour
  notre « $\Rightarrow$ », et « $\Rightarrow$ » pour notre
  « $\vdash$ ».\par}

\end{frame}
%
\begin{frame}
\frametitle{Le calcul propositionnel intuitionniste : connecteurs}

\itempoint Chaque connecteur logique a des règles
d'\textbf{introduction} permettant de \alert{démontrer} ce connecteur,
et des règles d'\textbf{élimination} permettant de l'\alert{utiliser}.

\bigskip

En français, et un peu informellement :
\begin{itemize}
\item pour introduire $\Rightarrow$ : on démontre $Q$ \alert{sous
  l'hypothèse} $P$, ce qui donne $P\Rightarrow Q$ (sans hypothèse :
  « hypothèse déchargée ») ;
\item pour éliminer $\Rightarrow$ : si on a $P\Rightarrow Q$ et $P$,
  on a $Q$ (\textit{modus ponens}) ;
\item pour introduire $\land$ : on démontre $Q_1$ et $Q_2$, ce qui
  donne $Q_1\land Q_2$ ; pour l'éliminer : si on a $Q_1 \land Q_2$ on
  en tire $Q_1$ resp. $Q_2$ ;
\item pour introduire $\lor$ : on démontre $Q_1$ ou $Q_2$, ce qui
  donne $Q_1\lor Q_2$ ;
\item pour éliminer $\lor$ : si on a $P_1 \lor P_2$, on démontre $Q$
  successivement sous les hypothèses $P_1$ et $P_2$, ce qui donne $Q$
  (hypothèses déchargées) ;
\item pour introduire $\top$, c'est trivial, et on ne peut pas
  l'éliminer ;
\item pour éliminer $\bot$, on tire la conclusion qu'on veut
  (\textit{ex falso quodlibet}), et on ne peut pas l'introduire.
\end{itemize}

\end{frame}
%
\begin{frame}
\label{natural-deduction-rules}
\frametitle{Le calcul propositionnel intuitionniste : règles}

Axiomes : si $Q \in \{P_1,\ldots,P_r\}$ alors $P_1,\ldots,P_r \vdash Q$.

\medskip

Introduction et élimination des connecteurs en style « déduction naturelle » :

\smallskip

\begin{tabular}{c|c|c}
&Intro&Élim\\\hline
$\Rightarrow$
&$\inferrule{\Gamma,{\color{mydarkgreen}P}\vdash Q}{\Gamma\vdash P\Rightarrow Q}$
&$\inferrule{\Gamma\vdash P\Rightarrow Q\\\Gamma\vdash P}{\Gamma\vdash Q}$
\\\hline
$\land$
&$\inferrule{\Gamma\vdash Q_1\\\Gamma\vdash Q_2}{\Gamma\vdash Q_1\land Q_2}$
&$\inferrule{\Gamma\vdash Q_1\land Q_2}{\Gamma\vdash Q_1}$
\quad$\inferrule{\Gamma\vdash Q_1\land Q_2}{\Gamma\vdash Q_2}$
\\\hline
$\lor$
&$\inferrule{\Gamma\vdash Q_1}{\Gamma\vdash Q_1\lor Q_2}$
\quad$\inferrule{\Gamma\vdash Q_2}{\Gamma\vdash Q_1\lor Q_2}$
&$\inferrule{\Gamma\vdash P_1\lor P_2\\\Gamma,{\color{mydarkgreen}P_1}\vdash Q\\\Gamma,{\color{mydarkgreen}P_2}\vdash Q}{\Gamma\vdash Q}$
\\\hline
$\top$
&$\inferrule{\strut}{\Gamma\vdash \top}$
&(néant)
\\\hline
$\bot$
&(néant)
&$\inferrule{\Gamma\vdash \bot}{\Gamma\vdash Q}$
\\
\end{tabular}

{\footnotesize Dans tout ça, $\Gamma$ désigne un jeu (\alert{non
    ordonné}, mais avec multiplicités) d'hypothèses.\par}

{\footnotesize Les hypothèses en \textcolor{mydarkgreen}{vert} sont
  « déchargées », c'est-à-dire qu'elles disparaissent.\par}

\end{frame}
%
\begin{frame}
\label{example-propositional-proofs}
\frametitle{Exemples de démonstrations}

Indiquant à côté de chaque barre une abréviation de la règle utilisée
(« $\Rightarrow$Int » pour introduction de $\Rightarrow$,
« $\land$Élim$_2$ » pour la 2\textsuperscript{e} règle d'élimination
de $\land$, etc.) :

\bigskip

{\footnotesize
\[
\inferrule*[left={$\Rightarrow$Int},right=\hphantom{Int}]{
\inferrule*[Left={$\land$Int}]{
\inferrule*[Left={$\land$Élim$_2$}]{
\inferrule*[Left={Ax}]{ }{A\land B \vdash A\land B}
}{A\land B \vdash B}
\\
\inferrule*[Right={$\land$Élim$_1$}]{
\inferrule*[Right={Ax}]{ }{A\land B \vdash A\land B}
}{A\land B \vdash A}
}{A\land B \vdash B\land A}
}{\vdash A\land B \Rightarrow B\land A}
\]
\par}

\bigskip

{\footnotesize
\[
\inferrule*[left={$\Rightarrow$Int},right=\hphantom{Int}]{
\inferrule*[Left={$\lor$Élim}]{
\inferrule*[Left={Ax}]{ }{A\lor B \vdash A\lor B}
\\
\inferrule*[right={$\lor$Int$_2$}]{
\inferrule*[Right={Ax}]{ }{A\lor B,A \vdash A}
}{A\lor B,A \vdash B\lor A}
\\
\inferrule*[Right={$\lor$Int$_1$}]{
\inferrule*[Right={Ax}]{ }{A\lor B,B \vdash B}
}{A\lor B,B \vdash B\lor A}
}{A\lor B \vdash B\lor A}
}{\vdash A\lor B \Rightarrow B\lor A}
\]
\par}

\end{frame}
%
\begin{frame}
\frametitle{Présentation différente}

On peut aussi n'écrire que les conclusions (partie droite du
signe « $\vdash$ »), à condition d'indiquer pour chaque hypothèse
déchargée à quel endroit elle l'a été (ceci sacrifie de la lisibilité
pour un gain de place) :

\medskip

{\footnotesize
\[
\inferrule*[left={$\Rightarrow$Int(\textcolor{mydarkgreen}{$u$})},right=\hphantom{Int}]{
\inferrule*[Left={$\land$Int}]{
\inferrule*[Left={$\land$Élim$_2$}]{
\inferrule*[Left={\textcolor{mydarkgreen}{$u$}}]{ }{A\land B}
}{B}
\\
\inferrule*[Right={$\land$Élim$_1$}]{
\inferrule*[Right={\textcolor{mydarkgreen}{$u$}}]{ }{A\land B}
}{A}
}{B\land A}
}{A\land B \Rightarrow B\land A}
\]
\par}

\medskip

{\footnotesize
\[
\inferrule*[left={$\Rightarrow$Int(\textcolor{mydarkgreen}{$u$})},right=\hphantom{Int}]{
\inferrule*[Left={$\lor$Élim(\textcolor{mydarkgreen}{$v$},\textcolor{mydarkgreen}{$v'$})}]{
\inferrule*[Left={\textcolor{mydarkgreen}{$u$}}]{ }{A\lor B}
\\
\inferrule*[right={$\lor$Int$_2$}]{
\inferrule*[Right={\textcolor{mydarkgreen}{$v$}}]{ }{A}
}{B\lor A}
\\
\inferrule*[Right={$\lor$Int$_1$}]{
\inferrule*[Right={\textcolor{mydarkgreen}{$v'$}}]{ }{B}
}{B\lor A}
}{B\lor A}
}{A\lor B \Rightarrow B\lor A}
\]
\par}

\medskip

{\footnotesize\textbf{N.B.} : une même hypothèse \emph{peut} être
  déchargée sur plusieurs endroits ($u$ dans le 1\textsuperscript{er}
  exemple).\par}

\end{frame}
%
\begin{frame}
\frametitle{Encore une présentation différente}

Présentation « drapeau », plus proche de l'écriture naturelle, commode
à vérifier :

\begin{footnotesize}
\begin{flagderiv}[example-proof1]
\assume{conj}{A\land B}{}
\step{partb}{B}{$\land$Élim$_2$ sur \ref{conj}}
\step{parta}{A}{$\land$Élim$_1$ sur \ref{conj}}
\step{newconj}{B\land A}{$\land$Int sur \ref{partb}, \ref{parta}}
\conclude{}{A\land B \Rightarrow B\land A}{$\Rightarrow$Int de \ref{conj} dans \ref{newconj}}
\end{flagderiv}
\end{footnotesize}

\vskip-2ex

\begin{footnotesize}
\begin{flagderiv}[example-proof2]
\assume{disj}{A\lor B}{}
\assume{parta}{A}{}
\step{fromparta}{B\lor A}{$\lor$Int$_2$ sur \ref{parta}}
\done
\assume{partb}{B}{}
\step{frompartb}{B\lor A}{$\lor$Int$_1$ sur \ref{partb}}
\done
\step{newdisj}{B\lor A}{$\lor$Élim sur \ref{disj} de \ref{parta} dans \ref{fromparta} et de \ref{partb} dans \ref{frompartb}}
\conclude{}{A\lor B \Rightarrow B\lor A}{$\Rightarrow$Int de \ref{disj} dans \ref{newdisj}}
\end{flagderiv}
\end{footnotesize}

\end{frame}
%
\begin{frame}
\frametitle{Quelques théorèmes du calcul propositionnel intuitionniste}

\itempoint Lorsque $\vdash P$ (sans hypothèse), on dit que $P$ est
\textbf{valable} dans le calcul propositionnel intuitionniste, ou est
un \textbf{théorème} / une \textbf{tautologie}.

\medskip

{\footnotesize

Quelques théorèmes importants (« $P\Leftrightarrow Q$ » abrège
« $P \Rightarrow Q$ et $Q \Rightarrow P$ ») :

\begin{tabular}{|c|c|}
\hline
$A\Rightarrow A$\quad\quad\quad $A\Rightarrow B\Rightarrow A$&$(A\Rightarrow B\Rightarrow C)
\Rightarrow (A\Rightarrow B)\Rightarrow A\Rightarrow C$\\
\hline
$(A\Rightarrow B\Rightarrow C) \Leftrightarrow (A\land B \Rightarrow C)$
&$A\land A \mathrel{\color{olive}\Leftrightarrow} A$\quad\quad\quad$A\lor A \mathrel{\color{olive}\Leftrightarrow} A$\\
\hline
$A\land B \Rightarrow A$\quad\quad\quad$A\land B \Rightarrow B$&$A\Rightarrow A\lor B$\quad\quad\quad$B \Rightarrow A\lor B$\\
\hline
$A\Rightarrow B\Rightarrow A\land B$&$(A\Rightarrow C)\Rightarrow(B\Rightarrow C)\Rightarrow A\lor B\Rightarrow C$\\
\hline
$A\land B \Leftrightarrow B\land A$&$A\lor B \Leftrightarrow B\lor A$\\
\hline
$(A\land B)\land C \Leftrightarrow A\land(B\land C)$&$(A\lor B)\lor C \Leftrightarrow A\lor(B\lor C)$\\
\hline
$(A\Rightarrow B)\land (A\Rightarrow C) \Leftrightarrow (A\Rightarrow B\land C)$
&$(A\Rightarrow C)\land (B\Rightarrow C) \Leftrightarrow (A\lor B \Rightarrow C)$\\
\hline
$(A\Rightarrow B)\lor (A\Rightarrow C) \mathrel{\color{red}\Rightarrow} (A\Rightarrow B\lor C)$
&$(A\Rightarrow C)\lor (B\Rightarrow C) \mathrel{\color{red}\Rightarrow} (A\land B \Rightarrow C)$\\
\hline
$A\land(B\lor C) \Leftrightarrow (A\land B)\lor (A\land C)$
&$A\lor(B\land C) \mathrel{\color{olive}\Leftrightarrow} (A\lor B)\land (A\lor C)$\\
\hline
$\top$&$\bot \Rightarrow C$\\
\hline
$\top\land A \Leftrightarrow A$\quad\quad\quad$\bot \land A \Leftrightarrow \bot$
&$\top\lor A \mathrel{\color{olive}\Leftrightarrow} \top$\quad\quad\quad$\bot \lor A \Leftrightarrow A$\\
\hline
$\neg A \land \neg B \Leftrightarrow \neg(A\lor B)$&$\neg A \lor \neg B \mathrel{\color{red}\Rightarrow} \neg(A\land B)$\\
\hline
$A \mathrel{\color{red}\Rightarrow} \neg\neg A$&$(\neg\neg A\Rightarrow \neg\neg B) \Leftrightarrow \neg\neg (A\Rightarrow B)$\\
\hline
$(\neg\neg A\land \neg\neg B) \Leftrightarrow \neg\neg (A\land B)$
&$(\neg\neg A\lor \neg\neg B) \mathrel{\color{red}\Rightarrow} \neg\neg (A\lor B)$\\
\hline
$(A\Rightarrow B) \mathrel{\color{red}\Rightarrow} (\neg B\Rightarrow \neg A)$
&$\neg A \Leftrightarrow \neg\neg\neg A$\\
\hline
$\neg(A\land\neg A)$&$\neg\neg(A\lor\neg A)$\\
\hline
\end{tabular}

\par}

\end{frame}
%
\begin{frame}
\frametitle{Quelques non-théorèmes}

{\footnotesize Rappelons qu'\textcolor{blue}{on ne permet pas de
    raisonnement par l'absurde} dans notre logique.\par}

\bigskip

Les énoncés suivants \alert{ne sont pas} des théorèmes du calcul
propositionnel intuitionniste (même s'ils \emph{sont} valables en
calcul propositionnel \emph{classique}) :
\begin{itemize}
\item $((A\Rightarrow B)\Rightarrow A)\Rightarrow A$ (« loi de Peirce »)
\item $A \lor \neg A$ (« tiers exclu »)
\item $\neg\neg A \Rightarrow A$ (« élimination de la double
  négation »)
\item $\neg(A\land B) \Rightarrow \neg A\lor\neg B$ (une des « lois de
  De Morgan »)
\item $(A\Rightarrow B) \Rightarrow \neg A\lor B$ {\footnotesize (la
  réciproque est bien valable intuitionnistement)}
\item $(\neg B\Rightarrow \neg A) \Rightarrow (A\Rightarrow B)$
  {\footnotesize (même remarque)}
\item $(A\Rightarrow B) \lor (B\Rightarrow A)$ (« loi de Dummett »)
\item $\neg A \lor \neg\neg A$ (« tiers exclu faible »)
\item $(\neg\neg A \Rightarrow A) \Rightarrow (A \lor \neg A)$ {\tiny
  (même si $(\neg\neg A \Rightarrow A)$ pour \emph{tout} $A$ donne $A
  \lor \neg A$ pour tout $A$)}
\end{itemize}

\bigskip

{\footnotesize\textcolor{brown}{Mais comment sait-on que quelque chose
    \emph{n'est pas} un théorème, au juste ?}\par}

\end{frame}
%
\begin{frame}
\frametitle{Preuve de la négation vs raisonnement par l'absurde}

\itempoint Rappel : la négation $\neg A$ abrège $A \Rightarrow \bot$
{\footnotesize(« si $A$ est vrai alors \textsc{absurdité} »)}.

\bigskip

Bien distinguer :

\medskip

\itempoint d'une part la \textbf{preuve de la négation} de $A$ :

\smallskip

\textcolor{teal}{« Supposons $A$ [vrai].  Alors (...), ce qui est
  absurde.  Donc $A$ est faux [i.e., $\neg A$]. »}

\smallskip

qui \alert{est valable} intuitionnistement : c'est prouver $A
\Rightarrow \bot$ par la règle $\Rightarrow$Intro,

\medskip

\itempoint et la \textbf{preuve par l'absurde} de $A$ d'autre part :

\smallskip

\textcolor{teal}{« Supposons $A$ faux [i.e., $\neg A$].  Alors (...),
  ce qui est absurde.  Donc $A$ est vrai. »}

\smallskip

qui \alert{n'est pas valable} intuitionnistement : tout ce qu'on peut
en tirer est $\neg\neg A$.

\bigskip

\itempoint On peut lire $\neg A$ comme « $A$ est faux » et $\neg\neg
A$ comme « $A$ n'est pas faux ».

\smallskip

{\footnotesize

\itempoint Noter que $\neg\neg\neg A$ équivaut à $\neg A$ (donc
$\neg\neg\neg\neg A$ à $\neg\neg A$, etc.).

\par}

\end{frame}
%
\begin{frame}
\frametitle{Logique classique}

La logique \textbf{classique} ou \textbf{booléenne} modifie la règle
\[
\inferrule*[left={$\bot$Élim}]{\Gamma\vdash \bot}{\Gamma\vdash Q}
\quad\quad\quad
\text{~en~}
\quad\quad\quad
\inferrule*[left={Absurde}]{\Gamma,{\color{mydarkgreen}\neg Q}\vdash \bot}{\Gamma\vdash Q}
\]

\bigskip

Elle est donc \alert{plus forte} (= a plus de théorèmes) que la
logique intuitionniste.

\bigskip

\itempoint\textbf{Théorème} {\footnotesize (« complétude de la
  sémantique booléenne du calcul propositionnel classique »)} : $P$
est un théorème de la logique classique \alert{ssi} pour toute
substitution de $\bot$ ou $\top$ à chacune des variables
propositionnelles de $P$ a la valeur de vérité $\top$.  (Voir
transp. \ref{truth-table-semantics} plus loin pour précisions.)

\bigskip

Tableaux de vérité :

\smallskip

\begin{tabular}{c@{\hskip 30bp}c@{\hskip 30bp}c}
$
\begin{array}{c|cc}
\land&\bot&\top\\\hline
\bot&\bot&\bot\\
\top&\bot&\top\\
\end{array}
$
&
$
\begin{array}{c|cc}
\lor&\bot&\top\\\hline
\bot&\bot&\top\\
\top&\top&\top\\
\end{array}
$
&
$
\begin{array}{c|cc}
A \Rightarrow B&B=\bot&B=\top\\\hline
A=\bot&\top&\top\\
A=\top&\bot&\top\\
\end{array}
$
\end{tabular}

\end{frame}
%
\begin{frame}
\frametitle{Un peu d'histoire des maths : Brouwer et l'intuitionnisme}

\itempoint L'\textbf{intuitionnisme} est à l'origine une philosophie
des mathématiques développée (à partir de 1912) par L. E. J. Brouwer
(1881–1966) en réaction/opposition au \textbf{formalisme} promu par
D. Hilbert (1862–1943).

\medskip

Quelques principes de l'intuitionnisme \emph{à l'origine} :

\begin{itemize}
\item rejet de la loi du tiers exclu (« principe d'omniscience ») pour
  demander des preuves \alert{constructives} (pour Brouwer, montrer
  que $x$ ne peut pas ne pas exister n'est pas la même chose que
  montrer que $x$ existe),
\item refus de la logique formelle (pour Brouwer, les maths sont
  « créatives »),
\item principes de continuité (notamment : toute fonction $\mathbb{R}
  \to \mathbb{R}$ est continue).
\end{itemize}

\medskip

\itempoint Idées rejetées par Hilbert (et la plupart des
mathématiciens).

$\rightarrow$ « controverse Hilbert-Brouwer » {\footnotesize
  (rapidement devenue querelle personnelle)}

\medskip

Hilbert : « Retirer au mathématicien le principe du tiers exclu
serait comme empêcher à un boxeur d'utiliser ses poings. »

\end{frame}
%
\begin{frame}
\frametitle{Les maths constructives après Brouwer}

\itempoint L'intuitionnisme a été ensuite reformulé et modifié par les
élèves de Brouwer et leurs propres élèves (notamment A. Heyting,
A. Troelstra) et d'autres écoles (analyse constructive d'E. Bishop,
constructivisme russe de A. Markov fils) :

\begin{itemize}
\item dans le cadre de la logique formelle (ou compatible avec elle),
\item sans impliquer de principe \emph{contredisant} les maths
  classiques,
\item mais toujours \alert{sans la loi du tiers exclu}.
\end{itemize}

\bigskip

\itempoint Les termes de « \textbf{mathématiques constructives} » et
« intuitionnisme » sont devenus plus ou moins interchangeables pour
« maths sans le tiers exclu ».

\bigskip

\itempoint Dans \emph{certains} tels systèmes, prouver $P\lor Q$
resp. $\exists x.P(x)$ passe forcément par la preuve de $P$ ou de $Q$,
resp. par la construction d'un $x$ vérifiant $P(x)$.

\bigskip

\itempoint Les maths « normales » se font en logique classique, mais
certains bouts (explicit\textsuperscript{t} signalés comme tels) de la
littérature sont en logique intuitionniste.  On peut y faire de
l'algèbre, de l'analyse, etc., constructives.

\end{frame}
%
\begin{frame}
\frametitle{Pourquoi vouloir « boxer sans ses poings » ?}

{\footnotesize Évidence : l'\alert{écrasante majorité des maths}
  se fait en logique classique (loi du tiers exclu admise) !\par}

\medskip

Une preuve « constructive » (sans tiers exclu) est plus restrictive
qu'une preuve classique, donc \alert{apporte plus} :

\begin{itemize}
\item principe de parcimonie ; \quad \itempoint curiosité purement théorique ;
\item intérêt philosophique (pas de « principe d'omniscience ») ;
\item validité dans un cadre mathématique plus large
  ($\rightarrow$ « topos ») ;
\item lien avec les systèmes de typage (via Curry-Howard) ;
\item extraction d'algorithme (p.ex., d'une preuve de $\forall
  m. \exists n. P(m,n)$ dans certains systèmes on peut extraire un
  algo qui \alert{calcule} $n$ à partir de $m$) ;
\item compatibilité avec des axiomes qui contredisent les maths
  classiques (p.ex. « toute fonction $\mathbb{R}\to\mathbb{R}$ est
  continue », « toute fonction $\mathbb{N}\to\mathbb{N}$ est
  calculable » [$\leftarrow$ « calculabilité synthétique »]) qu'on
  peut vouloir étudier.
\end{itemize}

\end{frame}
%
\begin{frame}
\frametitle{Preuves classiques pas toujours satisfaisantes}

Que penser de la preuve suivante ?

\medskip

\textbf{Affirmation :} Il existe un algo qui donné $k \in
\mathbb{N}$ \alert{termine en temps fini} et renvoie
\begin{itemize}
\item soit le rang de la première occurrence de $k$ chiffres $7$
  consécutifs dans l'écriture décimale de $\pi$,
\item soit « $\infty$ » si une telle occurrence n'existe pas.
\end{itemize}

\medskip

{\footnotesize

\underline{Preuve} (classique !) : soit $f\colon \mathbb{N} \to
\mathbb{N} \cup \{\infty\}$ la fonction qui à $k$ associe le rang
recherché.

\alert{De deux choses l'une :}
\begin{itemize}
\item soit il existe $k_0$ tel que $f(k) = \infty$ pour $k\geq k_0$ :
  alors $f$ est calculable car $f$ est déterminée par $k_0$ et un
  nombre \emph{fini} de valeurs $(f(0),\ldots,f(k_0-1))$ ;
\item soit $f(k)$ est fini pour tout $k \in \mathbb{N}$ : alors $f$
  est calculable par l'algorithme qui, donné $k$, calcule indéfiniment
  des décimales de $\pi$ jusqu'à en trouver $k$ consécutives
  valant $7$, et renvoie le rang (cet algorithme termine toujours par
  hypothèse).
\end{itemize}

Dans tous les cas $f$ est calculable.\qed
\par}

\medskip

\textcolor{blue}{Objection :} cette preuve (correcte en maths
classiques) ne nous donne pas du tout d'algorithme !  Elle montre
juste qu'il « existe » classiquement.

\end{frame}
%
\begin{frame}
\label{bhk-interpretation}
\frametitle{L'interprétation de Brouwer-Heyting-Kolmogorov}

Interprétation \alert{informelle/intuitive} des connecteurs de la
logique intuitionniste, due à A. Kolmogorov, A. Heyting, G. Kreisel,
A. Troelstra et d'autres :

\begin{itemize}
\item un témoignage de $P\land Q$, est un témoignage de $P$ et un de $Q$,
\item un témoignage de $P\lor Q$, est un témoignage de $P$ ou un de
  $Q$, et la donnée duquel des deux on a choisi,
\item un témoignage de $P\Rightarrow Q$ est un moyen de transformer
  un témoignage de $P$ en un témoignage de $Q$,
\item un témoignage de $\top$ est trivial, \quad \itempoint un
  témoignage de $\bot$ n'existe pas,
\item {\color{darkgray} un témoignage de $\forall x.P(x)$ est un moyen
  de transformer un $x$ quelconque en un témoignage de $P(x)$,}
\item {\color{darkgray} un témoignage de $\exists x.P(x)$ est la
  donnée d'un certain $x_0$ et d'un témoignage de $P(x_0)$.}
\end{itemize}

\medskip

{\footnotesize J'écris « témoignage », mais Kolmogorov parlait de
  « solution » d'un problème, Heyting de « preuve », etc.\par}

\end{frame}
%
\begin{frame}
\frametitle{Correspondance de Curry-Howard : implication seule}

\begin{tabular}{c|c}
Typage du $\lambda$-calcul simplement typé
&Calcul propositionnel intuitionniste\\[1ex]\hline\\
$\inferrule*[left={Var}]{ }{\Gamma,x:\sigma \vdash x:\sigma}$
&
$\inferrule*[left={Ax}]{ }{\Gamma,P \vdash P}$
\\[2ex]
$\inferrule*[left={App},sep=3em]{\Gamma \vdash f:\sigma\to\tau\\ \Gamma
  \vdash x:\sigma}{\Gamma \vdash fx:\tau}$
&
$\inferrule*[left={$\Rightarrow$Élim},sep=3em]{\Gamma \vdash P\Rightarrow Q\\ \Gamma
  \vdash P}{\Gamma \vdash Q}$
\\[2ex]
$\inferrule*[left={Abs}]{\Gamma, v:\sigma \vdash t:\tau}{\Gamma
  \vdash \lambda(v:\sigma).t : \sigma\to\tau}$
&
$\inferrule*[left={$\Rightarrow$Int}]{\Gamma, P \vdash Q}{\Gamma
  \vdash P\Rightarrow Q}$
\end{tabular}

\bigskip

\itempoint Ce sont \alert{exactement les mêmes règles}, aux
notations/nommage près…

\smallskip

\itempoint …sauf que la colonne de droite n'a pas le terme typé, mais
on sait qu'on peut le reconstruire
(transp. \ref{reconstructing-typed-term}), à partir de l'arbre de
dérivation.

\smallskip

\itempoint On peut donc \alert{identifier} termes du $\lambda$CST et
arbres de preuve en déduction naturelle du calcul propositionnel
intuitionniste restreint au seul connecteur $\Rightarrow$.

\end{frame}
%
\begin{frame}
\label{example-proposition-proof-with-implication}
\frametitle{Correspondance de Curry-Howard : exemple avec implication}

\itempoint Transformons en démonstration le terme
\[\lambda(f:\beta\to\alpha\to\gamma). \lambda(x:\alpha). \lambda(y:\beta). fyx\]
qu'on a typé (transp. \ref{example-typing-derivation}) comme
$(\beta\to\alpha\to\gamma) \to (\alpha\to\beta\to\gamma)$ :

{\footnotesize
\[
\inferrule*[Left={$\Rightarrow$Int}]{
\inferrule*[Left={$\Rightarrow$Int}]{
\inferrule*[Left={$\Rightarrow$Int}]{
\inferrule*[Left={$\Rightarrow$Élim}]{
\inferrule*[Left={$\Rightarrow$Élim}]{
\inferrule*[Left={Ax}]{ }{\mbox{$\begin{array}{cc}\scriptscriptstyle B \Rightarrow A \Rightarrow C ,\\ \scriptscriptstyle A ,B \\ \end{array}$} \vdash B \Rightarrow A \Rightarrow C }
\\
\inferrule*[Right={Ax}]{ }{\mbox{$\begin{array}{cc}\scriptscriptstyle B \Rightarrow A \Rightarrow C ,\\ \scriptscriptstyle A ,B \\ \end{array}$} \vdash B }
}{\mbox{$\begin{array}{cc}\scriptscriptstyle B \Rightarrow A \Rightarrow C ,\\ \scriptscriptstyle A ,B \\ \end{array}$} \vdash A \Rightarrow C }
\\
\inferrule*[Right={Ax}]{ }{\mbox{$\begin{array}{cc}\scriptscriptstyle B \Rightarrow A \Rightarrow C ,\\ \scriptscriptstyle A ,B \\ \end{array}$} \vdash A }
}{B \Rightarrow A \Rightarrow C , A , B  \vdash C }
}{B \Rightarrow A \Rightarrow C , A  \vdash B \Rightarrow C }
}{B \Rightarrow A \Rightarrow C  \vdash A \Rightarrow B \Rightarrow C }
}{\vdash (B \Rightarrow A \Rightarrow C ) \Rightarrow  (A \Rightarrow B \Rightarrow C )}
\]
\par}

\end{frame}
%
\begin{frame}
\frametitle{Correspondance de Curry-Howard : conjonction}

On veut \alert{étendre} le $\lambda$CST avec un \textbf{type produit}
pour refléter les règles de la conjonction logique :

\medskip

\begin{tabular}{c|c}
Typage du $\lambda$-calcul
&Calcul propositionnel intuitionniste\\[1ex]\hline\\
$\inferrule*[left=Proj$_1$]{\Gamma \vdash t:\tau_1\times\tau_2}{\Gamma \vdash \pi_1 t:\tau_1}$
&$\inferrule*[left={$\land$Élim$_1$}]{\Gamma\vdash Q_1\land Q_2}{\Gamma\vdash Q_1}$
\\[2ex]
$\inferrule*[left=Proj$_2$]{\Gamma \vdash t:\tau_1\times\tau_2}{\Gamma \vdash \pi_2 t:\tau_2}$
&$\inferrule*[left={$\land$Élim$_2$}]{\Gamma\vdash Q_1\land Q_2}{\Gamma\vdash Q_2}$
\\[2ex]
$\inferrule*[left=Pair]{\Gamma \vdash t_1:\tau_1\\\Gamma \vdash t_2:\tau_2}{\Gamma \vdash \langle t_1,t_2\rangle : \tau_1\times\tau_2}$
&$\inferrule*[left={$\land$Int}]{\Gamma\vdash Q_1\\\Gamma\vdash Q_2}{\Gamma\vdash Q_1\land Q_2}$
\end{tabular}

\medskip

\itempoint Ici, $\langle\emdash,\emdash\rangle$ sert à construire des
couples, et $\pi_1,\pi_2$ à les déconstruire.

\smallskip

{\footnotesize

\itempoint Nouvelle règle de réduction : $\pi_i\langle t_1,t_2\rangle$
se réduit en $t_i$ (pour $i\in\{1,2\}$).

\par}

\end{frame}
%
\begin{frame}
\label{curry-howard-example-with-conjunction}
\frametitle{Correspondance de Curry-Howard : exemple avec conjonction}

\itempoint Transformons en programme la démonstration qu'on a donnée
de $A\land B \Rightarrow B\land A$
(transp. \ref{example-propositional-proofs} et suivants) :

{\footnotesize
\[
\inferrule*[left={Abs},right=\hphantom{Abs}]{
\inferrule*[Left={Pair}]{
\inferrule*[Left={Proj$_2$}]{
\inferrule*[Left={Var}]{ }{u:\alpha\times \beta \vdash u:\alpha\times \beta}
}{u:\alpha\times \beta \vdash \pi_2 u:\beta}
\\
\inferrule*[Right={Proj$_1$}]{
\inferrule*[Right={Var}]{ }{u:\alpha\times \beta \vdash u:\alpha\times \beta}
}{u:\alpha\times \beta \vdash \pi_1 u:\alpha}
}{u:\alpha\times \beta \vdash \langle \pi_2 u, \pi_1 u\rangle:\beta\times \alpha}
}{\vdash \lambda(u:\alpha\times\beta).\,\langle \pi_2 u, \pi_1 u\rangle : \alpha\times \beta \rightarrow \beta\times \alpha}
\]
\par}

\bigskip

\itempoint Il s'agit de la fonction
$\lambda(u:\alpha\times\beta).\,\langle \pi_2 u, \pi_1 u\rangle$
(polymorphe de type $\alpha\times\beta \to \beta\times\alpha$) qui
échange les deux termes d'un couple.

\end{frame}
%
\begin{frame}
\frametitle{Correspondance de Curry-Howard : disjonction}

On veut étendre le $\lambda$CST avec un \textbf{type somme}
pour refléter les règles de la disjonction logique :

\medskip

\begin{tabular}{c|c}
Typage du $\lambda$-calcul
&Calcul propositionnel intuitionniste\\[1ex]\hline\\
$\inferrule*[left={Inj$_1$}]{\Gamma\vdash t:\tau_1}{\Gamma\vdash \iota_1^{(\tau_1,\tau_2)} t: \tau_1 + \tau_2}$
&$\inferrule*[left={$\lor$Int$_1$}]{\Gamma\vdash Q_1}{\Gamma\vdash Q_1\lor Q_2}$
\\[2ex]
$\inferrule*[left={Inj$_2$}]{\Gamma\vdash t:\tau_2}{\Gamma\vdash \iota_2^{(\tau_1,\tau_2)} t: \tau_1 + \tau_2}$
&$\inferrule*[left={$\lor$Int$_2$}]{\Gamma\vdash Q_2}{\Gamma\vdash Q_1\lor Q_2}$
\\[2ex]
\textcolor{gray}{ci-dessous $\downarrow$}
&{\footnotesize $\inferrule*[left={$\lor$Élim}]{\Gamma\vdash P_1\lor P_2\\\Gamma,P_1\vdash Q\\\Gamma,P_2\vdash Q}{\Gamma\vdash Q}$}
\end{tabular}

\[
\inferrule*[left={Match}]{\Gamma\vdash r:\sigma_1+ \sigma_2\\\Gamma,v_1:\sigma_1\vdash t_1:\tau\\\Gamma,v_2:\sigma_2\vdash t_2:\tau}{\Gamma\vdash
(\texttt{match~}r\texttt{~with~}\iota_1 v_1 \mapsto t_1,\; \iota_2 v_2 \mapsto t_2) : \tau}
\]

{\footnotesize\textbf{N.B.} : $v_1,v_2$ sont des
  \alert{variables} qui sont \alert{liées} par le \texttt{match} ; et
  $r,t_1,t_2$ sont des \alert{termes}.\par}

\end{frame}
%
\begin{frame}
\frametitle{Remarques sur types produits et sommes}

\itempoint $\iota_1,\iota_2$ sont des \emph{constructeurs} dans la
terminologie OCaml.

\medskip

\itempoint À côté de la $\beta$-réduction usuelle du $\lambda$-calcul,
$(\lambda (v:\sigma). e)t \rightsquigarrow e[v\backslash t]$, on
introduit des nouvelles règles de réduction pour la conjonction et la
disjonction :
\begin{itemize}
\item $\pi_i\langle t_1,t_2\rangle \;\rightsquigarrow\; t_i$ (pour
  $i\in\{1,2\}$)
\item $(\texttt{match~}\iota_i^{(\tau_1,\tau_2)}s\texttt{~with~}\iota_1
  v_1 \mapsto t_1,\; \iota_2 v_2 \mapsto t_2) \;\rightsquigarrow\;
  t_i[v_i\backslash s]$ (pour $i\in\{1,2\}$)
\end{itemize}

\medskip

\itempoint Côté démonstrations : cette réduction court-circuite une
règle \textsc{Intro} immédiatement suivie par sa règle \textsc{Élim}
(ce qu'on appelle un « détour »).

\medskip

\itempoint Le $\lambda$-calcul simplement typé \alert{reste fortement
  normalisant} avec ces extensions par types produits et sommes et les
réductions ci-dessus {\footnotesize (et types $1$ et $0$ après)}.

\medskip

\itempoint Les injections $\iota_i : \tau_i \to \tau_1+\tau_2$ portent
l'exposant $(\tau_1,\tau_2)$ pour que les annotations de type soient
complètes (mais il est inutile dans le matching).

\medskip

\itempoint Aucune des notations
$\pi_i,\ \langle,\rangle,\ \iota_i,\ \texttt{match}...\texttt{with}$
n'est standardisée (bcp de variations existent), mais il n'y a aucun
doute sur la correspondance elle-même.

\end{frame}
%
\begin{frame}
\frametitle{Correspondance de Curry-Howard : exemple avec disjonction}

\itempoint Transformons en programme la démonstration qu'on a donnée
de $A\lor B \Rightarrow B\lor A$
(transp. \ref{example-propositional-proofs} et suivants) :

{\footnotesize
\[
\inferrule*[left={Abs},right=\hphantom{Int}]{
\inferrule*[Left={Match}]{
\inferrule*[Left={Var}]{ }{u:\alpha+\beta \vdash u:\alpha+\beta}
\\
\inferrule*[right={Inj$_2$}]{
\inferrule*[Right={Var}]{ }{\ldots,v:\alpha \vdash v:\alpha}
}{{\scriptstyle\cdots} \vdash \iota_2^{(\beta,\alpha)} v : \beta+\alpha}
\\
\inferrule*[Right={Inj$_1$}]{
\inferrule*[Right={Var}]{ }{\ldots,v':\beta \vdash v':\beta}
}{{\scriptstyle\cdots} \vdash \iota_1^{(\beta,\alpha)} v' : \beta+\alpha}
}{u:\alpha+\beta \vdash (\texttt{match~} u\texttt{~with~}\iota_1
  v \mapsto \iota^{(\beta,\alpha)}_2 v,\; \iota_2 v' \mapsto \iota_1^{(\beta,\alpha)} v') : \beta+\alpha}
}{\vdash \lambda (u:\alpha+\beta) .
(\texttt{match~} u\texttt{~with~}\iota_1
  v \mapsto \iota^{(\beta,\alpha)}_2 v,\; \iota_2 v' \mapsto \iota_1^{(\beta,\alpha)} v') : \alpha+\beta \to \beta+\alpha}
\]
\par}

\bigskip

\itempoint Il s'agit de la fonction
\[
\lambda (u:\alpha+\beta) .
(\texttt{match~} u\texttt{~with~}\iota_1
  v \mapsto \iota^{(\beta,\alpha)}_2 v,\; \iota_2 v' \mapsto \iota_1^{(\beta,\alpha)} v')
\]
(polymorphe de type $\alpha+\beta \to \beta+\alpha$) qui échange les
deux cas d'une somme.

\end{frame}
%
\begin{frame}
\frametitle{Correspondance de Curry-Howard : vrai et faux}

On veut étendre le $\lambda$CST avec un \textbf{type unité} et un
\textbf{type vide} pour refléter les règles du vrai et du faux :

\medskip

\begin{tabular}{c|c}
Typage du $\lambda$-calcul
&Calcul propositionnel intuitionniste\\[1ex]\hline\\
$\inferrule*[left={Unit}]{ }{\Gamma\vdash \bullet:1}$
&$\inferrule*[left={$\top$Int}]{ }{\Gamma\vdash \top}$
\\[2ex]
$\inferrule*[left={Void}]{\Gamma\vdash r:0}{\Gamma\vdash
\texttt{exfalso}^{(\tau)} r : \tau}$
&$\inferrule*[left={$\bot$Élim}]{\Gamma\vdash \bot}{\Gamma\vdash Q}$
\end{tabular}

\medskip

\itempoint Ici, $\bullet$ désigne la valeur triviale de type unité
(\texttt{()} en OCaml), et $\texttt{exfalso}$ est un matching vide.
{\footnotesize (Notations pas standardisées du tout.)}

\smallskip

\itempoint Pas de nouvelle règle de réduction à ajouter.

\medskip

{\footnotesize

En OCaml :

{\tt
type void = | ;;
let exfalso = fun (r:void) -> match r with \_ -> . ;;
\par}

\par}

\end{frame}
%
\begin{frame}
\frametitle{Correspondance de Curry-Howard : exemples divers}

\itempoint L'implication $(A\Rightarrow B\Rightarrow C) \Rightarrow
(A\Rightarrow B)\Rightarrow A\Rightarrow C$ est démontrée par le terme
(« combinateur $\mathsf{S}$ ») :
\[
\lambda (x:\alpha\to\beta\to\gamma).\,
\lambda (y:\alpha\to\beta).\,
\lambda (z:\alpha).\,
xz(yz)
\]

\medskip

\itempoint L'équivalence $(A\land B \Rightarrow C) \Leftrightarrow
(A\Rightarrow B\Rightarrow C)$ est démontrée par les fonctions de
« curryfication »
\[
\lambda (f:\alpha\times\beta\to\gamma).\,
\lambda (x:\alpha).\, \lambda(y:\beta).\,
f\langle x,y\rangle
\]
et « décurryfication »
\[
\lambda (f:\alpha\to\beta\to\gamma).\,
\lambda (z:\alpha\times\beta).\,
f(\pi_1 z)(\pi_2 z)
\]

\medskip

\itempoint L'équivalence $A\lor(B\land C) \Leftrightarrow (A\lor
B)\land (A\lor C)$ est démontrée par les termes
{\footnotesize
\[
\lambda (u:\alpha+(\beta\times\gamma)).\,
(\texttt{match~}u\texttt{~with~}\iota_1 v \mapsto \langle \iota^{(\alpha,\beta)}_1 v, \iota^{(\alpha,\gamma)}_1 v\rangle,\;
\iota_2 w \mapsto \langle \iota^{(\alpha,\beta)}_2 (\pi_1 w), \iota^{(\alpha,\gamma)}_2 (\pi_2 w)\rangle)
\]
\vskip-1ex et\vskip-5ex
\[
\begin{aligned}
\lambda (u:(\alpha+\beta)\times(\alpha+\gamma)).\,
&(\texttt{match~}\pi_1 u\texttt{~with~}\iota_1 v \mapsto \iota_1^{(\alpha,\beta\times\gamma)} v,\\
\iota_2 v' \mapsto
&(\texttt{match~}\pi_2 u\texttt{~with~}\iota_1 w \mapsto \iota_1^{(\alpha,\beta\times\gamma)} w,\;
\iota_2 w' \mapsto \iota_2^{(\alpha,\beta\times\gamma)}\langle v',w'\rangle ))
\end{aligned}
\]
}

\end{frame}
%
\begin{frame}
\frametitle{Correspondance de Curry-Howard : exemple en OCaml}

Reprise de l'équivalence $A\lor(B\land C) \Leftrightarrow (A\lor
B)\land (A\lor C)$ typée par OCaml :

\bigskip

{\footnotesize\tt
let pi1 = fun (x,y) -> x ;; (* $\pi_1$ *)\newline
{\color{purple}val pi1 : 'a * 'b -> 'a = <fun>}\newline
let pi2 = fun (x,y) -> y ;; (* $\pi_2$ *)\newline
{\color{purple}val pi2 : 'a * 'b -> 'b = <fun>}\newline
type ('a, 'b) sum = Inj1 of 'a | Inj2 of 'b ;; (* $\alpha+\beta$ *)\newline
{\color{purple}type ('a, 'b) sum = Inj1 of 'a | Inj2 of 'b}\newline
fun u -> match u with Inj1 v -> (Inj1 v, Inj1 v)\newline
\ \ | Inj2 w -> (Inj2 (pi1 w), Inj2 (pi2 w)) ;;\newline
{\color{purple}- : ('a, 'b * 'c) sum -> ('a, 'b) sum * ('a, 'c) sum = <fun>}\newline
fun u -> (match (pi1 u) with Inj1 v -> Inj1 v\newline
\ \ | Inj2 v\_ -> (match (pi2 u) with Inj1 w -> Inj1 w\newline
\ \ \ \ | Inj2 w\_ -> Inj2 (v\_,w\_))) ;;\newline
{\color{purple}- : ('a, 'b) sum * ('a, 'c) sum -> ('a, 'b * 'c) sum = <fun>}
\par}

\end{frame}
%
\begin{frame}
\frametitle{Les preuves ne sont pas uniques}

{\footnotesize La question de l'égalité est compliquée, on ne
  l'abordera guère.\par}

\itempoint Une même proposition peut avoir des preuves
(fonctionnellement) \alert{différentes}.

\medskip

Par exemple, les entiers de Church typés :
\begin{itemize}
\item $\overline{0}_{\alpha} :=
  \lambda(f:\alpha\to\alpha).\lambda(x:\alpha).x$
\item $\overline{1}_{\alpha} :=
  \lambda(f:\alpha\to\alpha).\lambda(x:\alpha).fx$
\item $\overline{2}_{\alpha} :=
  \lambda(f:\alpha\to\alpha).\lambda(x:\alpha).f(fx)$ etc.
\end{itemize}
tous de type $(\alpha\to\alpha)\to(\alpha\to\alpha)$ prouvent tous
$(A\Rightarrow A)\Rightarrow(A\Rightarrow A)$.

\medskip

{\footnotesize

($\overline{0}$) : « [\textbf{(1)} Supposons $A\Rightarrow A$.
    [\textbf{(2)} Supposons $A$.]  \textbf{(3)} On a $A\Rightarrow A$
    par $\Rightarrow$Int de (2) dans (2).]  \textbf{(4)} On a
  $(A\Rightarrow A)\Rightarrow(A\Rightarrow A)$ par $\Rightarrow$Int
  de (1) dans (3). »

\medskip

($\overline{1}$) : « [\textbf{(1)} Supposons $A\Rightarrow A$.
  [\textbf{(2)} Supposons $A$.  \textbf{(3)} On a $A$ par
    $\Rightarrow$Élim sur (1) et (2).]  \textbf{(4)} On a
  $A\Rightarrow A$ par $\Rightarrow$Int de (2) dans (3).]
\textbf{(5)} On a $(A\Rightarrow A)\Rightarrow(A\Rightarrow A)$ par
$\Rightarrow$Int de (1) dans (4). »

\medskip

($\overline{2}$) : « [\textbf{(1)} Supposons $A\Rightarrow A$.
  [\textbf{(2)} Supposons $A$.  \textbf{(3)} On a $A$ par
    $\Rightarrow$Élim sur (1) et (2).  \textbf{(4)} On a $A$ par
    $\Rightarrow$Élim sur (1) et (3).]  \textbf{(5)} On a
  $A\Rightarrow A$ par $\Rightarrow$Int de (2) dans (4).]
\textbf{(6)} On a $(A\Rightarrow A)\Rightarrow(A\Rightarrow A)$ par
$\Rightarrow$Int de (1) dans (5). »

\par}

\end{frame}
%
\begin{frame}
\frametitle{Curry-Howard : récapitulation}

\itempoint La correspondance de Curry-Howard permet
d'\alert{identifier}
\begin{itemize}
\item\textcolor{blue}{types} du $\lambda$-calcul simplement typé,
  éventuellement enrichi de constructions de types produit ($\times$),
  somme ($+$), trivial ($1$) et vide ($0$), et
\item\textcolor{blue}{propositions} (=formules logiques) du calcul
  propositionnel intuitionniste avec pour connecteurs
  l'implication ($\Rightarrow$) et éventuellement la
  conjonction ($\land$), disjonction ($\lor$), vrai ($\top$) et
  faux ($\bot$) respectivement,
\end{itemize}
en identifant aussi
\begin{itemize}
\item\textcolor{blue}{termes} ayant les types en question,
\item\textcolor{blue}{preuves} des propositions en question, dans le
  style « déduction naturelle », en calcul propositionnel
  intutionniste.
\end{itemize}

\medskip

\itempoint Noter aussi : abstraction$\leftrightarrow$ouverture
d'hypothèse ; application$\leftrightarrow$\textit{modus ponens} ;
variables$\leftrightarrow$hypothèses ; variables
liées$\leftrightarrow$hypothèses déchargées.

\bigskip

{\footnotesize

\itempoint On se permettra maintenant parfois des abus de notation
justifiés par Curry-Howard, p.ex., traiter $\rightarrow$ et
$\Rightarrow$ comme interchangeables.

\par}

\end{frame}
%
\begin{frame}
\frametitle{La négation et la double négation}

\textcolor{orange}{Qu'est-ce qui correspond à la proposition $\neg P$
  par Curry-Howard ?}

\medskip

C'est le type $\sigma \rightarrow 0$ des fonctions prenant un argument
de type $\sigma$ et renvoyant une valeur impossible, i.e., ne peuvent
jamais renvoyer.

\medskip

Mais on parle d'un langage ($\lambda$CST enrichi) où \alert{tous les
  programmes terminent} (« normalisation forte ») !  Donc une telle
fonction prouve que $\sigma$ est lui-même vide ; et la fonction est
triviale : c'est une « pure preuve » de vacuité de $\sigma$ :

\smallskip

\begin{itemize}
\item le type $\sigma \rightarrow 0$ (ou « $\neg\sigma$ ») est le type
  des « témoignages de vacuité » de $\sigma$ \textcolor{teal}{(si on
    me fournit un truc de type $\sigma$, je soulève une exception
    parce que c'est impossible)},
\item le type $(\sigma \rightarrow 0) \rightarrow 0$ (ou
  « $\neg\neg\sigma$ ») est une sorte de type des témoignages de
  \alert{non-vacuité} de $\sigma$,
\item mais ce dernier ne permet pas « magiquement » d'en tirer une
  valeur :
\item pas de terme de type $((\alpha \rightarrow 0) \rightarrow 0) \to
  \alpha$ ou bien $\alpha + (\alpha\rightarrow 0)$ dans le
  $\lambda$CST : on est bien en \alert{logique intuitionniste}.
\end{itemize}

\end{frame}
%
\begin{frame}
\label{variables-are-effectively-polymorphic}
\frametitle{Un embryon de polymorphisme}

\itempoint Les types du $\lambda$CST sont écrits avec des
\alert{variables de types} $\alpha,\beta,\gamma$…  En principe ce sont
des \alert{types opaques}.  En pratique, une fonction comme
$\lambda(x:\alpha).\lambda(y:\beta).x$ de type
$\alpha\to\beta\to\alpha$ se comporte \alert{comme polymorphe} : on
peut imaginer un $\forall\alpha,\beta$ (\alert{implicite}) devant :

\medskip

\itempoint En effet, substituer n'importe quel type $\sigma$ à une
variable de type $\alpha$ dans un terme du $\lambda$CST (enrichi)
donne encore un terme correct (la dérivation de typage est la même,
après substitution).

\medskip

\itempoint Conséquence côté logique : substituer des propositions
quelconques aux \alert{variables propositionnelles} d'un théorème
donne encore un théorème.

\medskip

P.ex. : $(A\land B\Rightarrow C) \Rightarrow (A\Rightarrow
B\Rightarrow C)$ est un théorème, donc
\[
((D\Rightarrow E)\land (E\Rightarrow D) \Rightarrow D) \Rightarrow ((D\Rightarrow E)\Rightarrow (E\Rightarrow D)\Rightarrow D)
\]
en est (par subst\textsuperscript{ion} de $D\Rightarrow E$ pour $A$,
de $E\Rightarrow D$ pour $B$, et de $D$ pour $C$).

\medskip

{\footnotesize

\itempoint Attention, ce polymorphisme s'applique aux
\alert{conclusions, pas aux hypothèses} (l'hypothèse $A$ ne permet pas
de tout déduire !).

\par}

\end{frame}
%
\begin{frame}
\frametitle{« Fonctorialité » des connecteurs logiques}

\begin{itemize}
\item Donnés $t_1 : \tau_1 \to \tau_1'$ et $t_2 : \tau_2 \to \tau_2'$,
  le terme $\lambda (u : \tau_1\times\tau_2).\langle t_1(\pi_1 u),
  t_2(\pi_2 u)\rangle$ est de type $\tau_1\times\tau_2 \to
  \tau_1'\times\tau_2'$.
\item Donnés $t_1 : \tau_1 \to \tau_1'$ et $t_2 : \tau_2 \to \tau_2'$,
  le terme $\lambda (u :
  \tau_1+\tau_2). (\texttt{match~}u\texttt{~with~}\iota_1 v_1 \mapsto
  \iota_1^{(\tau_1',\tau_2')} (t_1v_1),\; \iota_2 v_2 \mapsto
  \iota_2^{(\tau_1',\tau_2')} (t_2v_2))$\\ est de type $\tau_1+\tau_2 \to
  \tau_1'+\tau_2'$.
\item Donnés $s : \sigma' \to \sigma$ (\textcolor{olive}{attention au
  sens !}) et $t : \tau \to \tau'$, le terme $\lambda (u :
  \sigma\to\tau). \lambda (x:\sigma'). t(u(sx))$ est de type
  $(\sigma\to\tau) \to (\sigma'\to\tau')$.
\end{itemize}

\bigskip

Donc, par Curry-Howard :

\begin{itemize}
\item Si $Q_1\Rightarrow Q_1'$ et $Q_2\Rightarrow Q_2'$ alors
  $(Q_1\land Q_2)\Rightarrow (Q_1'\land Q_2')$.
\item Si $Q_1\Rightarrow Q_1'$ et $Q_2\Rightarrow Q_2'$ alors
  $(Q_1\lor Q_2)\Rightarrow (Q_1'\lor Q_2')$.
\item Si $P'\Rightarrow P$ (\textcolor{olive}{attention au sens !}) et
  $Q\Rightarrow Q'$ alors $(P\Rightarrow Q)\Rightarrow (P'\Rightarrow
  Q')$.
\end{itemize}

\bigskip

{\footnotesize

\itempoint On dit que $\land$ et $\lor$ sont
\textcolor{purple}{\textbf{croissants}} (ou
\textcolor{purple}{covariants}) en leurs deux arguments, et que
$\Rightarrow$ l'est dans son argument de droite, mais qu'il est
\textcolor{olive}{\textbf{décroissant}} (ou
\textcolor{olive}{contravariant}) dans son argument de gauche.

\par}

\end{frame}
%
\begin{frame}
\frametitle{Sous-formules positives et négatives}

Dans une formule propositionnelle, on définit les sous-formules
\textcolor{purple}{\textbf{positives}} (voire \textbf{strictement
  positives}) et \textcolor{olive}{\textbf{négatives}} par induction :
\begin{itemize}
\item les sous-formules \textcolor{purple}{positives} de $Q_1\land Q_2$ et $Q_1\lor Q_2$
  sont la formule tout entière, et les sous-formules \textcolor{purple}{positives} de
  $Q_1$ et celles de $Q_2$,
\item les sous-formules \textcolor{olive}{négatives} de $Q_1\land Q_2$ et $Q_1\lor Q_2$
  sont les sous-formules \textcolor{olive}{négatives} de $Q_1$ et celles de $Q_2$,
\item les sous-formules \textcolor{purple}{positives} de $P \Rightarrow Q$ sont la formule
  tout entière, les sous-formules \textcolor{purple}{positives} de $Q$ et les
  \emph{\textcolor{olive}{négatives}} de $P$,
\item les sous-formules \textcolor{olive}{négatives} de $P \Rightarrow
  Q$ sont les sous-formules \textcolor{olive}{négatives} de $Q$ et les
  \emph{\textcolor{purple}{positives}} de $P$.
\end{itemize}

\smallskip

{\footnotesize

\itempoint Les sous-formules strict\textsuperscript{t} positives de
$Q_1\land Q_2$ et $Q_1\lor Q_2$, resp. $P\Rightarrow Q$ sont la
formule tout entière et les sous-formules strict\textsuperscript{t}
positives de $Q_1$ et de $Q_2$ (resp. $Q$).

\par}

\medskip

{\footnotesize

P.ex. dans\vskip-5ex
\[
({\color{olive}A}\Rightarrow ({\color{olive}B}\Rightarrow {\color{purple}C})) \land (({\color{purple}D}\Rightarrow {\color{olive}E})\Rightarrow {\color{purple}F})
\]
les occurrences
${\color{purple}C},{\color{purple}D},{\color{purple}F}$ sont positives
($C$ et $F$ strictement), tandis que
${\color{olive}A},{\color{olive}B},{\color{olive}E}$ sont négatives.

\par}

\end{frame}
%
\begin{frame}
\frametitle{« Fonctorialité » des formules}

Conséquence de la fonctorialité des connecteurs logiques :

\smallskip

\itempoint Si $S$ est une formule propositionnelle et $S'$ obtenue en
remplaçant une sous-formule positive $Q$ par $Q'$ telle que $Q
\Rightarrow Q'$, alors $S \Rightarrow S'$.

\smallskip

\itempoint Si $S$ est une formule propositionnelle et $S'$ obtenue en
remplaçant une sous-formule négative $P$ par $P'$ telle que $P'
\Rightarrow P$, alors $S \Rightarrow S'$.

\bigskip

Comme toute sous-formule est soit positive soit négative, on en
déduit :

\smallskip

\itempoint Corollaire : Si $S$ est une formule propositionnelle et
$S'$ obtenue en remplaçant une sous-formule quelcone $R$ par $R'$
telle que $R \Leftrightarrow R'$, alors $S \Leftrightarrow S'$.

\medskip

P.ex., on peut remplacer $Q_1\land Q_2$ par $Q_2\land Q_1$ ou
$(Q_1\land Q_2)\land Q_3$ par $Q_1\land (Q_2\land Q_3)$ n'importe où
dans une formule et on obtient ainsi une formule équivalente.

\end{frame}
%
\begin{frame}
\frametitle{Calcul des séquents}

\itempoint Le \textbf{calcul des séquents} est une autre présentation
de la (même) logique propositionnelle intuitionniste (mêmes
théorèmes, mêmes séquents).

\medskip

\itempoint Cette présentation est peut-être moins « naturelle », mais
elle est plus symétrique, et a divers intérêts théoriques.

\medskip

\itempoint La différence principale porte sur les \alert{règles
  d'élimination} (transp. suivant) : au lieu d'une règle
d'introduction et d'une règle d'élimination, on a une \textbf{règle
  droite} (= introduction) et une \textbf{règle gauche}
($\leftrightarrow$ élimination), qui introduisent le connecteur à
droite \alert{ou à gauche} du symbole « $\vdash$ ».

\medskip

\itempoint On garde la règle $\inferrule*[left=Ax]{ }{\Gamma,Q \vdash
  Q}$ ; on va reparler des règles de « contraction » et « coupure » :

\begin{center}
\begin{tabular}{c|c}
$\inferrule*[left=Contr]{\Gamma,P,P \vdash Q}{\Gamma,P\vdash Q}$
&$\inferrule*[left=Cut]{\Gamma \vdash P\\\Gamma,P\vdash Q}{\Gamma\vdash Q}$
\end{tabular}
\end{center}

\end{frame}
%
\begin{frame}
\label{sequent-calculus-connector-rules}
\frametitle{Calcul des séquents : règles (intuitionnistes) des connecteurs}

{\footnotesize

La colonne de droite est la même que la colonne de gauche (intro) du
transp. \ref{natural-deduction-rules}.  C'est la colonne de gauche qui
est « nouvelle » :

\par}

\bigskip

\begin{tabular}{c|c|c}
&Gauche&Droite\\\hline
$\Rightarrow$
&$\inferrule{\Gamma\vdash M\\\Gamma,P\vdash R}{\Gamma,M\Rightarrow P\vdash R}$
&$\inferrule{\Gamma,P\vdash Q}{\Gamma\vdash P\Rightarrow Q}$
\\\hline
$\land$
&$\inferrule{\Gamma,P_1\vdash R}{\Gamma,P_1\land P_2\vdash R}$
\quad$\inferrule{\Gamma,P_2\vdash R}{\Gamma,P_1\land P_2\vdash R}$
&$\inferrule{\Gamma\vdash Q_1\\\Gamma\vdash Q_2}{\Gamma\vdash Q_1\land Q_2}$
\\\hline
$\lor$
&$\inferrule{\Gamma,P_1\vdash R\\\Gamma,P_2\vdash R}{\Gamma,P_1\lor P_2\vdash R}$
&$\inferrule{\Gamma\vdash Q_1}{\Gamma\vdash Q_1\lor Q_2}$
\quad$\inferrule{\Gamma\vdash Q_2}{\Gamma\vdash Q_1\lor Q_2}$
\\\hline
$\top$
&(néant)
&$\inferrule{\strut}{\Gamma\vdash \top}$
\\\hline
$\bot$
&$\inferrule{\strut}{\Gamma,\bot\vdash R}$
&(néant)
\\
\end{tabular}

\end{frame}
%
\begin{frame}
\frametitle{Exemples de démonstrations en calcul des séquents}

On reprend les mêmes exemples que dans le
transp. \ref{example-propositional-proofs} :
\vskip-2ex\leavevmode
\begin{center}
\begin{tabular}{c|c}
%
{\footnotesize
$
\inferrule*[left={$\Rightarrow$R},right=\hphantom{Int}]{
\inferrule*[Left={$\land$R}]{
\inferrule*[Left={$\land$L$_2$}]{
\inferrule*[Left={Ax}]{ }{B \vdash B}
}{A\land B \vdash B}
\\
\inferrule*[Right={$\land$L$_1$}]{
\inferrule*[Right={Ax}]{ }{A \vdash A}
}{A\land B \vdash A}
}{A\land B \vdash B\land A}
}{\vdash A\land B \Rightarrow B\land A}
$
}
&
{\footnotesize
$
\inferrule*[left={$\Rightarrow$R},right=\hphantom{Int}]{
\inferrule*[Left={$\lor$L}]{
\inferrule*[Left={$\lor$R$_2$}]{
\inferrule*[Left={Ax}]{ }{A \vdash A}
}{A \vdash B\lor A}
\\
\inferrule*[Right={$\lor$R$_1$}]{
\inferrule*[Right={Ax}]{ }{B \vdash B}
}{B \vdash B\lor A}
}{A\lor B \vdash B\lor A}
}{\vdash A\lor B \Rightarrow B\lor A}
$
}
\end{tabular}
\end{center}

\medskip

Et que dans le transp. \ref{example-proposition-proof-with-implication} :

\vskip-2ex
{\footnotesize
\[
\inferrule*[Left={$\Rightarrow$R}]{
\inferrule*[Left={$\Rightarrow$R}]{
\inferrule*[Left={$\Rightarrow$R}]{
\inferrule*[Left={$\Rightarrow$L}]{
\inferrule*[Left={Ax}]{ }{A, B \vdash B }
\\
\inferrule*[Right={$\Rightarrow$L}]{
\inferrule*[Left={Ax}]{ }{A , B  \vdash A }
\\
\inferrule*[Right={Ax}]{ }{C , A , B  \vdash C  }
}{A \Rightarrow C , A , B  \vdash C }
}{B \Rightarrow A \Rightarrow C , A , B  \vdash C }
}{B \Rightarrow A \Rightarrow C , A  \vdash B \Rightarrow C }
}{B \Rightarrow A \Rightarrow C  \vdash A \Rightarrow B \Rightarrow C }
}{\vdash (B \Rightarrow A \Rightarrow C ) \Rightarrow  (A \Rightarrow B \Rightarrow C )}
\]
\par}

\end{frame}
%
\begin{frame}
\label{sequent-calculus-structural-rules}
\frametitle{Calcul des séquents : règles « structurales »}

\itempoint On peut choisir de séparer la règle $\inferrule*[left=Ax]{
}{\Gamma,Q \vdash Q}$ en deux, une règle d'axiome strict et une règle
d'« affaiblissement » :

\begin{center}
\begin{tabular}{c|c}
$\inferrule*[left=Ax]{ }{Q \vdash Q}$
&$\inferrule*[left=Weak]{\Gamma \vdash Q}{\Gamma,P\vdash Q}$
\end{tabular}
\end{center}

\medskip

\itempoint Comme en déduction naturelle, les hypothèses n'ont pas
d'ordre.

\medskip

\itempoint Si elles ont des multiplicités, la règle de contraction est
nécessaire (sans elle, on ne peut pas prouver $(A\Rightarrow B) \land A
\vdash B$) :
\[
\inferrule*[left=Contr]{\Gamma,P,P \vdash Q}{\Gamma,P\vdash Q}
\]
On peut s'en dispenser en décidant que les hypothèses n'ont pas de
multiplicité (forment un ensemble, pas un multi-ensemble), ou en
modifiant les règles $\Rightarrow$L et $\land$L pour ne pas perdre
d'hypothèse en remontant.

\end{frame}
%
\begin{frame}
\frametitle{Calcul des séquents : équivalence avec la déduction naturelle}

L'équivalence générale entre déduction naturelle et calcul des
séquents n'est pas difficile si on utilise librement la règle de
coupure.

\medskip

Montrons l'exemple de la conversion entre $\Rightarrow$Élim et
$\Rightarrow$Left :

\medskip

\itempoint Dans le sens déduction naturelle $\to$ calcul des
séquents :
\vskip-3ex\leavevmode
\begin{center}
\scalebox{0.8}{%
\begin{tabular}{c|c}
Déduction naturelle&Calcul des séquents\\\hline
$\inferrule*[left=$\Rightarrow$Élim]{\Gamma\vdash P\Rightarrow Q\\\Gamma\vdash P}{\Gamma\vdash Q}$
&
$\inferrule*[left=Cut]{\Gamma\vdash P\Rightarrow Q
\\
\inferrule*[left=$\Rightarrow$L]{\Gamma\vdash P\\
\inferrule*[left=Ax]{ }{\Gamma,Q\vdash Q}
}{\Gamma,P\Rightarrow Q\vdash Q}
}{\Gamma\vdash Q}$
\end{tabular}
}
\end{center}

\itempoint Dans le sens calcul des séquents $\to$ déduction
naturelle :
\vskip-6ex\leavevmode
\begin{center}
\scalebox{0.8}{%
\begin{tabular}{c|c}
Calcul des séquents&Déduction naturelle\\\hline
$\inferrule*[left=$\Rightarrow$L]{\Gamma\vdash M\\\Gamma,P\vdash R}{\Gamma,M\Rightarrow P\vdash R}$
&
$\inferrule*[left=$\Rightarrow$Élim]{
\inferrule*[Left={$\Rightarrow$Intro}]{\Gamma,P\vdash R}{\Gamma\vdash P\Rightarrow R}
\\
\inferrule*[left=$\Rightarrow$Élim]{
\inferrule*[Left={Ax}]{ }{\Gamma,M\Rightarrow P\vdash M\Rightarrow P}
\\\Gamma\vdash M}{\Gamma,M\Rightarrow P\vdash P}
}{\Gamma,M\Rightarrow P\vdash R}$
\end{tabular}
}
\end{center}

\end{frame}
%
\begin{frame}
\frametitle{Calcul des séquents : élimination des coupures}

La règle de coupure exprime une forme de transitivité des démonstrations :
\[
\inferrule*[left=Cut]{\Gamma \vdash P\\\Gamma,P\vdash Q}{\Gamma\vdash Q}
\]

\smallskip

(On dira que $P$ est la « formule coupée ».)

\medskip

\itempoint \textbf{Théorème} (Gentzen) : la règle de « coupure » n'est
pas nécessaire en calcul des séquents : tout séquent prouvable avec
elle est encore prouvable sans elle.

\smallskip

La preuve est même constructive et se fait par des transformations
« locales ».

\medskip

\itempoint En déduction naturelle, l'élimination des coupures est
facile : pour éliminer une coupure il suffit de reprendre la
démonstration de $\Gamma,P\vdash Q$ avec l'hypothèse $P$ en moins, et
d'utiliser celle de $\Gamma\vdash P$ partout où $P$ est invoquée
(cf. transp. \ref{substitution-lemma}).  {\footnotesize (La difficulté
  est plutôt l'élimination de « détours » Intro+Élim $\approx$
  normal\textsuperscript{ion}.)}

\medskip

\itempoint En calcul des séquents, la difficulté supplémentaire est
que $P$ peut faire intervenir un connecteur qui est introduit à droite
dans $\Gamma\vdash P$ et à gauche dans $\Gamma,P\vdash Q$
(cf. transp. \ref{cut-elimination-interesting-cases}).

\end{frame}
%
\begin{frame}
\label{cut-elimination-boring-cases}
\frametitle{Exemples d'étape d'élimination des coupures (1)}

\itempoint Cas « déplaçants » : la dernière règle d'une des branches
de la coupure n'opère pas sur la formule coupée : on fait
\alert{remonter} la coupure sur cette branche (quitte à choisir).

\begin{center}
\scalebox{0.65}{%
$\inferrule*[left={\color{purple}Cut}]{
\inferrule*[Left={$\land$R}]{\mpdotsabove{\Gamma \vdash P_1}\\\mpdotsabove{\Gamma \vdash P_2}}{\Gamma \vdash P_1\land P_2}
\\
\inferrule*[Right={?}]{\mpdotsabove{\Gamma',P_1\land P_2\vdash Q'}}{\Gamma,P_1\land P_2\vdash Q}
}{\Gamma\vdash Q}$
}
~devient~
\scalebox{0.65}{%
$\inferrule*[Right={?}]{
\inferrule*[Left={\color{purple}Cut}]{
\inferrule*[Left={$\land$R}]{\mpdotsabove{\Gamma \vdash P_1}\\\mpdotsabove{\Gamma \vdash P_2}}{\Gamma \vdash P_1\land P_2}
\\
\mpdotsabove{\Gamma',P_1\land P_2\vdash Q'}
}{\Gamma,\Gamma'\vdash Q'}
}{\Gamma\vdash Q}$
}
\end{center}

{\footnotesize (On n'a pas montré ici les règles structurales
  (contraction+affaiblissement) permettant d'avoir $\Gamma$ et/ou
  $\Gamma'$ comme hypothèses.)\par}

\bigskip

\itempoint Cas final : une des branches de la coupure est la règle
« axiome » sur la formule coupée : la coupure disparaît :

\begin{center}
\scalebox{0.8}{%
$\inferrule*[left={\color{purple}Cut}]{\inferrule*[left=Ax]{ }{\Gamma \vdash P}\\\mpdotsabove{\Gamma,P\vdash Q}}{\Gamma\vdash Q}$
}
~devient~
\scalebox{0.8}{%
$\mpdotsabove{\Gamma\vdash Q}$
}
\end{center}

\end{frame}
%
\begin{frame}
\label{cut-elimination-interesting-cases}
\frametitle{Exemples d'étape d'élimination des coupures (2)}

\itempoint Cas « principaux » : la dernière règle de chaque branche de
la coupure concerne le connecteur logique de la formule coupée : la
coupure ne remonte pas et peut même se multiplier, mais le « degré »
de complexité de la formule coupée diminue :

\begin{center}
\scalebox{0.8}{%
$\inferrule*[left={\color{purple}Cut}]{
\inferrule*[Left={$\land$R}]{\mpdotsabove{\Gamma \vdash P_1}\\\mpdotsabove{\Gamma \vdash P_2}}{\Gamma \vdash P_1\land P_2}
\\
\inferrule*[Right={$\land$L$_1$}]{\mpdotsabove{\Gamma,P_1\vdash Q}}{\Gamma,P_1\land P_2\vdash Q}
}{\Gamma\vdash Q}$
}
~devient~
\scalebox{0.8}{%
$\inferrule*[left={\color{purple}Cut}]{
\mpdotsabove{\Gamma \vdash P_1}
\\
\mpdotsabove{\Gamma,P_1\vdash Q}
}{\Gamma\vdash Q}$
}
\end{center}

\begin{center}
\scalebox{0.7}{%
$\inferrule*[left={\color{purple}Cut}]{
\inferrule*[Left={$\Rightarrow$R}]{\mpdotsabove{\Gamma, M\vdash P}}{\Gamma \vdash M\Rightarrow P}
\\
\inferrule*[Right={$\Rightarrow$L}]{\mpdotsabove{\Gamma \vdash M}\\\mpdotsabove{\Gamma,P\vdash Q}}{\Gamma,M\Rightarrow P\vdash Q}
}{\Gamma\vdash Q}$
}
~devient~
\scalebox{0.7}{%
$\inferrule*[left={\color{purple}Cut}]{
\mpdotsabove{\Gamma \vdash M}
\\
\inferrule*[Right={{\color{purple}Cut}}]{\mpdotsabove{\Gamma,M \vdash P}\\\mpdotsabove{\Gamma,P\vdash Q}}{\Gamma,M\vdash Q}
}{\Gamma\vdash Q}$
}
\end{center}

\end{frame}
%
\begin{frame}
\frametitle{Élimination des coupures : récurrence}

\itempoint Le \textbf{degré} $\deg P$ d'une formule propositionnelle
est : $0$ pour une variable propositionnelle ou $\top$, $\bot$, et
$\max(\deg P_1, \deg P_2) + 1$ si $P = P_1 \circ P_2$ pour un
connecteur $\circ \in \{\Rightarrow, \land, \lor\}$.  C'est donc la
profondeur de l'arbre de cette formule.

\smallskip

{\footnotesize P.ex., $\deg(((A\Rightarrow B)\Rightarrow A)\Rightarrow
  A) = 3$.\par}

\medskip

\itempoint Le degré d'une coupure est le degré de la formule coupée.
Le \textbf{degré de coupure} d'une démonstration en calcul des
séquents est le plus grand degré d'une coupure (ou $-1$ s'il n'y en a
pas).  Une \textbf{coupure critique} est une coupure de degré
maximal.

\medskip

\itempoint La preuve de l'élimination des coupures se fait par une
\alert{double récurrence} :
\begin{itemize}
\item récurrence principale sur le degré de coupure de la
  démonstration,
\item récurrence secondaire, à degré donné, sur la somme des hauteurs
  des deux branches de la coupure à éliminer.
\end{itemize}

\medskip

\itempoint La démonstration est constructive (algorithmique) ;
l'algorithme esquissé est p.r. (même « élémentaire »).  Il n'est pas
déterministe {\footnotesize (ni même confluent)}.

\end{frame}
%
\begin{frame}
\frametitle{Élimination des coupures : conséquences}

{\footnotesize\textcolor{blue}{Quel intérêt de pouvoir éliminer les
    coupures ?}  Les preuves sans coupure ne peuvent pas faire
  apparaître de formule inattendue.\par}

\medskip

\itempoint\textbf{Propriété de la disjonction :} si $\vdash Q_1\lor
Q_2$ (\alert{sans hypothèse !}) alors $\vdash Q_1$ ou bien $\vdash
Q_2$.  {\footnotesize (\underline{Preuve :} une démonstration sans
  coupure de $\vdash Q_1\lor Q_2$ doit finir par la règle $\lor$R$_1$
  ou $\lor$R$_2$. \qedsymbol)}

\medskip

\itempoint\textbf{Corollaire :} $\vdash A\lor\neg A$ n'est pas
prouvable en logique intuitionniste.

\medskip

\itempoint\textbf{Propriété de la sous-formule :} si $P_1,\ldots,P_r
\vdash Q$ alors une preuve sans coupure ne fait intervenir que des
formules qui sont des \alert{sous-formules} de $P_1,\ldots,P_r$ ou
$Q$.  {\footnotesize (\underline{Preuve :} c'est clair sur chacune des
  règles. \qedsymbol)}

\medskip

\itempoint\textbf{Corollaire :} on peut \alert{décider
  algorithmiquement} si $P_1,\ldots,P_r \vdash Q$.

\smallskip

{\footnotesize\underline{Algorithme :} construire l'ensemble $\Phi$ de
  toutes les sous-formules d'une de $P_1,\ldots,P_r$ ou $Q$, et
  l'ensemble de tous les séquents possibles $\Gamma\vdash R$ avec
  $\Gamma\subseteq\Phi$ et $R\in\Phi$.  Marquer comme « valables »
  ceux qui découlent de séquents déjà marqués comme valables par
  application d'une des règles du calcul des séquents (autre que la
  coupure), et répéter jusqu'à ce que le séquent recherché ait été
  marqué ou que plus aucun séquent ne soit marqué.\qed\par}

\end{frame}
%
\begin{frame}
\frametitle{Calcul des séquents et Curry-Howard}

\itempoint On a vu que les preuves en déduction naturelle
correspondent aux termes du $\lambda$-calcul simplement typé
(+extensions).

\smallskip

On peut trouver un correspondant aux preuves en calcul des séquents
{\footnotesize (avec de petites modifications)} : le
$\overline{\lambda}$-calcul de Herbelin.

\smallskip

{\footnotesize Le $\overline{\lambda}$-calcul a deux sortes d'objets :
  les \emph{termes} $t$ et les \emph{listes d'arguments} (notées
  $[t_1;\ldots;t_r]$).\par}

\medskip

\itempoint En travaillant un peu, on peut déduire l'élimination des
coupures de la conversion en $\overline{\lambda}$-calcul de termes
normaux du $\lambda$-calcul (donc de sa normal\textsuperscript{ion}).

\begin{center}
\scalebox{0.56}{
$\inferrule*{
\inferrule*{
\inferrule*{ }{y : C\vdash y : C}
\\
\inferrule*{ }{x_1:A_1\vdash x_1 : A_1}
}{y:C, x_1:A_1\vdash yx_1 : A_2\to B}
\\
\inferrule*{ }{x_2:A_2\vdash x_2 : A_2}
}{y:C, x_1:A_1, x_2:A_2\vdash yx_1 x_2 : B}$
}
$\rightsquigarrow$
\scalebox{0.56}{
$\inferrule*{
\inferrule*{
\inferrule*{ }{x_1:A_1\vdash x_1:A_1}
\\
\inferrule*{
\inferrule*{ }{x_2:A_2\vdash x_2:A_2}
\\
\inferrule*{ }{\_:B \vdash \_ [] : B}
}{x_2:A_2;\_:A_2\to B \vdash \_ [x_2] : B}
}{x_1:A_1, x_2:A_2;\_:C \vdash \_ [x_1; x_2] : B}
}{y:C, x_1:A_1, x_2:A_2\vdash y [x_1; x_2] : B}$
}

\footnotesize

(où $C := A_1\to A_2 \to B$)

(à gauche le $\lambda$-calcul, à droite le
$\overline{\lambda}$-calcul ; remarquez l'inversion des arguments)

($y$ est une \emph{variable} ici, pas une application ni une abstraction)
\end{center}

\end{frame}
%
\begin{frame}
\frametitle{Combinateurs $\mathsf{S},\mathsf{K},\mathsf{I}$}

\itempoint\textbf{Axiomes de Hilbert :} le calcul
prop\textsuperscript{nel} intuitionniste peut être défini par la
\alert{seule} règle du modus ponens (« si $P\Rightarrow Q$ et $P$
alors $Q$ ») à partir des schémas d'axiomes suivants, où
$A,B,C,\ldots$ sont \alert{remplacés par des formules qcqes} :
{\footnotesize
\begin{itemize}
\item $(\mathsf{I}): A\Rightarrow A$
\quad\itempoint $(\mathsf{K}): A\Rightarrow B\Rightarrow A$
\quad\itempoint $(\mathsf{S}): (A\Rightarrow B\Rightarrow C) \Rightarrow
  (A\Rightarrow B) \Rightarrow (A \Rightarrow C)$
\item $A\land B \Rightarrow A$
\quad\itempoint $A\land B \Rightarrow B$
\quad\itempoint $A\Rightarrow B \Rightarrow A\land B$
\quad\itempoint $\top$ \quad\itempoint $\bot \Rightarrow C$
\item $A \Rightarrow A \lor B$
\quad\itempoint $B \Rightarrow A \lor B$
\quad\itempoint $(A\Rightarrow C)\Rightarrow(B\Rightarrow C)\Rightarrow (A\lor B\Rightarrow C)$
\end{itemize}
\par}

\medskip

\itempoint Via Curry-Howard, ceci correspond au fait qu'on peut
définir un $\lambda$-terme quelconque (à $\beta$-réduction près) par
application des trois combinateurs $\mathsf{I} := \lambda x.x$,
$\mathsf{K} := \lambda x.\lambda y.x$ et $\mathsf{S} := \lambda
x. \lambda y. \lambda z. xz(yz)$.  Ceci vaut pour le $\lambda$-calcul
non typé comme simplement typé.  {\footnotesize (On n'a même pas besoin de
  $\mathsf{I}$ qui peut s'écrire $\mathsf{SKK}$.)}

\medskip

\itempoint Idée de la conversion : remplacer $\lambda x.x$ par
$\mathsf{I}$, $\lambda x.y$ par $\mathsf{K}y$ (si $x$ n'apparaît pas
dans $y$) et $\lambda x.uv$ par $\mathsf{S}(\lambda x.u)(\lambda
x.v)$.  {\footnotesize (On peut aussi « $\eta$-convertir » $\lambda
  x.fx$ en $f$.)}

{\footnotesize P.ex. : $\lambda f.\lambda x.f(fx)$ se réécrit $\lambda
  f. \mathsf{S}(\mathsf{K}f)f$ donc
  $\mathsf{S}(\mathsf{S}(\mathsf{K}\mathsf{S})\mathsf{K})\mathsf{I}$.\par}

\smallskip

{\footnotesize\itempoint On peut donc en théorie imaginer un langage
  de programmation fonctionnel Turing-complet sans variables, basé sur
  les seules fonctions $\mathsf{S}$, $\mathsf{K}$, $\mathsf{I}$.
  {\tiny(Mais personne ne serait assez fou pour faire ça.)}\par}

\end{frame}
%
\section{Le call/cc et la logique classique}
\begin{frame}
\frametitle{Qu'est-ce qu'une continuation ?}

\itempoint Dans un langage de programmation, la \textbf{continuation}
de l'appel d'une fonction $f$ est « l'état de la machine qui attend
que $f$ renvoie une valeur ».

\medskip

On peut y penser comme (une copie de) la \alert{pile d'appels} jusqu'à
l'appel de $f$.

\medskip

\itempoint Certains langages donnent la \alert{citoyenneté de première
  classe} aux continuations, i.e., permettent de les passer, stocker
et invoquer explicitement :
\begin{itemize}
\item\alert{capturer} la contin\textsuperscript{ion} de $f$ revient à
  garder une copie de la pile d'appels de $f$,
\item\alert{invoquer} la contin\textsuperscript{ion} avec une valeur
  $v$ a pour effet de faire retourner à $f$ la valeur $v$ (même si
  elle avait \alert{déjà} terminé),
\item c'est une sorte de \texttt{goto} jusqu'au point à $f$ termine,
  avec restauration de la pile (en C : \texttt{getcontext} pour
  capturer, \texttt{setcontext} pour restaurer).
\end{itemize}

\medskip

{\footnotesize

\itempoint Variante : capturer la continuation revient à créer un fil
d'exécution (\textit{thread}) mis en attente au moment du retour de
$f$, l'invoquer revient à terminer le fil actuel et réactiver le fil
en attente, avec en lui transmettant $v$ : il va de nouveau créer un
fil en attente, puis considérer $v$ comme valeur de retour de $f$.

\medskip

\itempoint Fondement théorique : le $\lambda\mu$-calcul de Parigot
(exten\textsuperscript{ion} du $\lambda$-calcul avec
continuations).

\par}

\end{frame}
%
\begin{frame}
\frametitle{Applications des continuations}

\textcolor{blue}{À quoi sert} d'avoir des continuations réifiées (= de
première classe) ?
\begin{itemize}
\item elles peuvent prendre la place d'un mécanisme
  d'\alert{exceptions} ou de sortie de boucles (au lieu de soulever
  une exception, on invoque la continuation « exceptionnelle » pour
  abandonner le calcul normal),
\item mais les continuations, contrairement aux exceptions, peuvent
  \alert{reprendre} un calcul interrompu (si sa continuation a été
  capturée),
\item elles permettent donc d'implémenter notamment : des générateurs
  ou du multitâche coopératif.
\end{itemize}

\bigskip

\textcolor{blue}{Le coût d'implémentation} réside essentiellement dans
la gestion de la pile :
\begin{itemize}
\item soit on recopie toute la pile à chaque
  capture/invocation de continuation,
\item soit le séquencement d'appels cesse d'être une pile et doit être
  géré par le \textit{garbage-collector} (= ramasse-miettes), c'est ce
  qui se passera avec CPS (cf. plus loin).
\end{itemize}

\medskip

Informellement, les \emph{continuations} sont aux appels de fonctions
ce que les \emph{clôtures} sont aux données locales.

\end{frame}
%
\begin{frame}
\frametitle{La fonction call/cc}

\itempoint « call/cc » = « call-with-current-continuation »

\medskip

\itempoint La fonction call/cc existe dans plusieurs langages de
programmation (notam\textsuperscript{t} : Scheme, SML/NJ, Ruby) : elle
prend en argument une fonction $g$ et
\begin{itemize}
\item\alert{capture} sa propre continuation (= celle du retour du
  call/cc),
\item\alert{passe} celle-ci en argument de la fonction $g$, et renvoie
  soit la valeur de retour de $g$ soit celle passée à la continuation.
\end{itemize}

\medskip

\itempoint En Scheme, la continuation se présente comme une fonction,
qu'on peut appeler avec un argument $v$ :
\begin{itemize}
\item la continuation elle-même ne termine jamais (puisque c'est,
  justement, une continuation : elle \alert{remplace} la pile d'appels
  par celle qui a été capturée),
\item elle a pour effet de faire retourner $v$ au call/cc qui l'a
  créée.
\end{itemize}

\medskip

{\footnotesize

\itempoint En SML/NJ, la fonction s'appelle
\texttt{SMLofNJ.Cont.callcc} et les continuations doivent être
invoquées avec la fonction \texttt{throw} ; mais on peut facilement
créer une fonction analogue à celle du Scheme :

{\tt
val callcc = fn g => SMLofNJ.Cont.callcc (fn k => g (fn v => SMLofNJ.Cont.throw k v))
}

\par}

\end{frame}
%
\begin{frame}
\frametitle{Exemples en Scheme}

{\tt
(define call/cc call-with-current-continuation)\hfill ;; Shorthand\newline
(call/cc (lambda (k) 42))\hfill ;; Return normally\newline
$\rightarrow$ 42\newline
(call/cc (lambda (k) (k 42)))\hfill ;; Return by invoking continuation\newline
$\rightarrow$ 42\newline
(call/cc (lambda (k) (+ (k 42) 1)))\hfill ;; (+ \_ 1) never reached\newline
$\rightarrow$ 42\newline
(* (call/cc (lambda (k) (k 42))) 2)\hfill ;; Nothing weird here\newline
$\rightarrow$ 84\newline
((call/cc (lambda (k) k))\newline
\ (call/cc (lambda (k) k)))\hfill ;; Endless loop: why?
\par}

\bigskip

{\footnotesize

Très difficile à comprendre :

{\tt
((lambda (yin)\newline
\ \ \ ((lambda (yang) (yin yang))\newline
\ \ \ \ ((lambda (kk) (display \#\textbackslash *) kk) (call/cc (lambda (k) k)))))\newline
\ ((lambda (kk) (newline) kk) (call/cc (lambda (k) k))))
\par}

}

\end{frame}
%
\begin{frame}
\frametitle{Quel est le type de call/cc ?}

\itempoint\alert{Une continuation ne retourne jamais}.  On peut donc
la typer comme $\alpha \to \bot$ ou bien $\alpha \to \beta$ avec
$\beta$ un type arbitraire.

\medskip

\itempoint La fonction $g$ passée au call/cc doit renvoyer le même
type $\alpha$ qu'accepte la continuation qu'on lui a passée : donc
$(\alpha\to\beta)\to\alpha$.

\medskip

\itempoint La fonction call/cc elle-même renvoie ce même type
$\alpha$ : donc elle a pour type
$((\alpha\to\beta)\to\alpha)\to\alpha$.

\medskip

\itempoint C'est le type correspondant à la \textbf{loi de Peirce}
$((A\Rightarrow B)\Rightarrow A)\Rightarrow A$, une des formulations
de la \alert{logique classique}.

\medskip

\textcolor{blue}{\textbf{Moralité :}} la présence du call/cc
transforme la logique du typage de logique intuitionniste en logique
classique.

\medskip

{\footnotesize Ceci est dit de façon informelle, mais on peut
  introduire une variante du $\lambda$-calcul, le $\lambda\mu$-calcul,
  qui rend précise cette idée.  Le $\lambda\mu$-calcul simplement typé
  garantit encore la terminaison des programmes : on ne peut pas faire
  de boucles avec le seul call/cc \emph{bien typé}.\par}

\end{frame}
%
\begin{frame}
\frametitle{Oui mais il y a de la triche !}

{\footnotesize Le tiers exclu $A \lor \neg A$ dit moralement « soit je
  te donne une valeur de type $\alpha$ soit je te donne une promesse
  qu'il n'en existe pas (type $\alpha\to 0$) ».  En $\lambda$CST on ne
  peut pas faire ça !}

\bigskip

Regardons comment le call/cc permet d'implémenter le tiers exclu $A
\lor \neg A$ :
\[
\texttt{callcc}\;
\bigg(\lambda (k:(\alpha+(\alpha\to 0))\to 0).\;\iota^{(\alpha,\alpha\to 0)}_2\Big(\lambda (v:\alpha).\,k(\iota^{(\alpha,\alpha\to 0)}_1 v)\Big)\bigg)
\]
\[
\text{a pour type~:~}\quad
\alpha + (\alpha\to 0)
\]

{\footnotesize La fonction en argument du call/cc est a pour type
  $((\alpha + (\alpha\to 0)) \to 0) \to (\alpha + (\alpha\to 0))$,
  c'est-à-dire qu'elle correspond à une preuve de $(\neg (A\lor\neg
  A)) \Rightarrow (A\lor\neg A)$ (intuitionist\textsuperscript{t}
  valable).\par}

\bigskip

Que fait ce code ?
\begin{itemize}
\item il renvoie \alert{provisoirement} $\iota_2^{(\alpha,\alpha\to
  0)}(\cdots)$, où $(\cdots)$ définit une promesse qu'il n'y a pas de
  valeur de type $\alpha$,
\item si on invoque cette promesse (avec une valeur $v$ de type
  $\alpha$, donc !), il utilise la continuation $k$ pour « revenir
  dans le temps » et changer d'avis et renvoie finalement
  $\iota_1^{(\alpha,\alpha\to 0)}(v)$.
\end{itemize}

\end{frame}
%
\begin{frame}
\frametitle{Oui mais il y a de la triche : quelle est la morale ?}

{\footnotesize

En SML/NJ :

{\tt
val callcc = fn g => SMLofNJ.Cont.callcc (fn k => g (fn v => SMLofNJ.Cont.throw k v))\newline
datatype ('a,'b) sum = Inj1 of 'a | Inj2 of 'b\newline
val exclmiddle = callcc (fn k => Inj2 (fn v => k (Inj1 v)))
;;
}

\par}

\medskip

\centerline{*}

\textcolor{blue}{\textbf{Moralité :}}

\medskip

\itempoint Il est facile de tenir ses promesses quand on peut voyager
dans le temps avec l'aide de call/cc.

\medskip

\itempoint Le call/cc change la logique du typage en logique
classique, mais l'intérêt des types somme est grandement diminué : les
valeurs ne sont pas stables.

\medskip

\itempoint La logique classique n'a pas la propriété de la
disjonction : on a $\vdash A\lor \neg A$ en logique classique, mais ni
$\vdash A$ ni $\vdash \neg A$ (pour $A$ opaque).  Elle n'est pas
constructive.  On peut lui appliquer Curry-Howard, mais c'est moins
intéressant.

\end{frame}
%
\begin{frame}
\frametitle{Continuation Passing Style}

\textcolor{brown}{Et si le langage n'a pas de call/cc ?}

\medskip

\itempoint Tous les langages n'ont pas de continuations de première
classe (« réifiées »), mais dans un langage fonctionnel, on peut
réécrire le code en « Continuation Passing Style » :
\begin{itemize}
\item au lieu d'écrire des fonctions qui \alert{renvoient une valeur}
  $v$, on les fait \alert{prendre en argument} une autre fonction $k$,
  qui est une continuation-de-fait, et appellent cette fonction sur la
  valeur $v$ ;
\item donc aucune fonction ne renvoie jamais rien : elle termine en
  invoquant son argument continuation-de-fait ou, plus souvent, en
  invoquant une autre fonction en lui passant une continuation-de-fait
  construite exprès pour continuer le calcul {\footnotesize (ceci rend
    le style excessivement lourd et pénible)} ;
\item ceci ne fonctionne bien que dans un langage supportant la
  récursion terminale propre (car \alert{tous} les appels deviennent
  terminaux).
\end{itemize}

\smallskip

{\footnotesize \itempoint C'est plus ou moins le concept des
  « promesses » de JavaScript, par exemple.\par}

\end{frame}
%
\begin{frame}
\frametitle{Continuation Passing Style : exemple en OCaml}

{\footnotesize\tt
let sum\_cps = fun x -> fun y -> fun k -> k(x+y) ;;\newline
{\color{purple}val sum\_cps : int -> int -> (int -> 'a) -> 'a = <fun>}\newline
let minus\_cps = fun x -> fun y -> fun k -> k(x-y) ;;\newline
{\color{purple}val minus\_cps : int -> int -> (int -> 'a) -> 'a = <fun>}\newline
let rec fibonacci\_cps = fun n -> fun k -> if n <= 1 then k n\newline
\ \ else minus\_cps n 1 (fun n1 -> minus\_cps n 2 (fun n2 ->\newline
\ \ \ \ fibonacci\_cps n1 (fun v1 -> fibonacci\_cps n2 (fun v2 ->\newline
\ \ \ \ \ \ sum\_cps v1 v2 k)))) ;;\newline
{\color{purple}val fibonacci\_cps : int -> (int -> 'a) -> 'a = <fun>}\newline
fibonacci\_cps 8 (fun x -> x) ;;\newline
{\color{purple}- : int = 21}\newline
let callcc\_cps = fun g -> fun k -> g (fun v -> fun k0 -> k v) k ;;\newline
(* translation of: callcc (fun kf -> ((kf 42) + 1)) *)\newline
callcc\_cps (fun kf -> fun k -> kf 42 (fun v -> sum\_cps v 1 k)) (fun x -> x) ;;\newline
{\color{purple}- : int = 42}
\par}

\end{frame}
%
\begin{frame}
\frametitle{Continuation Passing Style : systématisation}

\textcolor{blue}{Définissons la transformation CPS de façon
  systématique.}

\medskip

{\footnotesize Prenons ici les notations logiques pour les types :
  notamment, $\Rightarrow$ désigne le type fonction.\par}

\medskip

\itempoint Fixons un type $Z$ de « retour ultime ».  On pose
\alert{$\mathop{\sim}P := (P\Rightarrow Z)$} pour le type d'« une
continuation qui attend une valeur de type $P$ » et
$\mathop{\sim}\mathop{\sim}P$ pour « une valeur $P$ passée par
continuation ».

\medskip

\itempoint Noter que $x:P \;\vdash\; \lambda(k:\mathop{\sim}P).kx :
\mathop{\sim}\mathop{\sim}P$ (transformation d'une valeur « directe »
en valeur passée par continuation).

\medskip

{\footnotesize On si on préfère voir ça comme une fonction :
  $\lambda(x:P).\,\lambda(k:\mathop{\sim}P).\,kx$ a pour type $P
  \Rightarrow \mathop{\sim}\mathop{\sim}P$.\par}

\medskip

\itempoint On va passer \alert{toutes} les valeurs par continuation :
tous les types transformés prendront la forme
$\mathop{\sim}\mathop{\sim}T$.

\smallskip

Mais en plus de ça, les fonctions renvoient leur valeur par
continuation : une fonction de type $P \Rightarrow Q$ va devenir $P
\Rightarrow \mathop{\sim}Q \Rightarrow Z$, c'est-à-dire $P \Rightarrow
\mathop{\sim}\mathop{\sim}Q$, et sera elle-même passée par
continuation, donc $\mathop{\sim}\mathop{\sim}(P \Rightarrow
\mathop{\sim}\mathop{\sim}Q)$ {\footnotesize (sans compter que $P$ et
  $Q$ peuvent eux-mêmes changer)}.

\end{frame}
%
\begin{frame}
\frametitle{Continuation Passing Style : l'essence de la transformation}

\itempoint Définissons la transformation CPS d'abord dans le
$\lambda$-calcul \alert{non typé} : par induction sur la complexité du
terme :
\begin{itemize}
\item $v^{\cps} = \lambda k.kv$ si $v$ est une variable (c'est la
  transformation de $P$ en $\mathop{\sim}\mathop{\sim}P$ définie
  ci-dessus).
\item $(\lambda v.t)^{\cps} = \lambda k.\, k(\lambda v.\, t^{\cps})$
  (idem pour une fonction, dont le corps est CPS-ifié).
\item pour l'application : \[(fx)^{\cps} = \lambda
  k.\,f^{\cps}(\lambda f_0.\, x^{\cps}(\lambda x_0.\, f_0 x_0 k))\] ce
  code se comprend ainsi : on invoque $f^{\cps}$ pour recevoir sa
  valeur « directe » $f_0$ (qui est quand même une fonction dans le
  style CPS), puis $x^{\cps}$ pour recevoir sa valeur « directe »
  $x_0$, puis on appelle la fonction $f_0$ avec la valeur $x_0$ et la
  continuation $k$ de l'ensemble de l'expression,
\item $\texttt{callcc}^{\cps} = \lambda \ell.\, \ell(\lambda g.\,
  \lambda k.\, g(\lambda v.\, \lambda k_0.\, kv)\,k)$
\end{itemize}

\medskip

\itempoint Ce code \alert{porte les graines d'un interpréteur}
(eval/apply) du $\lambda$-calcul dans le $\lambda$-calcul : il impose
d'ailleurs l'évaluation en appel-par-valeurs.

\end{frame}
%
\begin{frame}
\frametitle{Continuation Passing Style : transformation des types}

{\footnotesize Il reste à typer tout ça !\par}

\smallskip

\itempoint Rappelons qu'on a fixé $Z$ et posé $\mathop{\sim}P :=
(P\Rightarrow Z)$.  On définit une transformation $P \mapsto
P^\diamond$ des types (=propositions) par (inductivement) :
\begin{itemize}
\item $A^\diamond = A$ si $A$ est une variable de type,
\item $(P\Rightarrow Q)^\diamond = (P^\diamond \Rightarrow \mathop{\sim}\mathop{\sim}Q^\diamond)$,
\item $(P\land Q)^\diamond = P^\diamond \land Q^\diamond$,
\quad\itempoint $(P\lor Q)^\diamond = P^\diamond \lor Q^\diamond$,
\item $\top^\diamond = \top$,
\quad\itempoint $\bot^\diamond = \bot$.
\end{itemize}

\smallskip

\itempoint On pose enfin $P^{\cps} :=
\mathop{\sim}\mathop{\sim}P^\diamond$ {\footnotesize (comprendre :
  $\mathop{\sim}\mathop{\sim}(P^\diamond)$)}.

\smallskip

En gros :
\begin{itemize}
\item $P^\diamond$ est le type $P$ dans lequel toutes les fonctions
  ont été réécrites en style CPS (= renvoient leurs valeurs par
  continuation),
\item $P^{\cps} := \mathop{\sim}\mathop{\sim}P^\diamond$ est le type
  en question lui-même passé par continuation.
\end{itemize}

\smallskip

{\footnotesize \itempoint Noter : $(P\Rightarrow Q)^{\cps} =
  \mathop{\sim}\mathop{\sim}(P^\diamond \Rightarrow
  \mathop{\sim}\mathop{\sim}Q^\diamond) =
  \mathop{\sim}\mathop{\sim}(P^\diamond \Rightarrow Q^{\cps})$.}

%% {\footnotesize
%% 
%% \itempoint On vérifie : $\mathop{\sim} P \Leftrightarrow
%% \mathop{\sim}\mathop{\sim}\mathop{\sim} P$, donc $P^{\cps}
%% \Leftrightarrow \mathop{\sim}\mathop{\sim} P^{\cps}$ ; remarquer
%% aussi : $\bot^{\cps} \Leftrightarrow Z$ et $(\neg P)^{\cps}
%% \Leftrightarrow \mathop{\sim} P^{\cps}$.
%% 
%% \par}

\end{frame}
%
\begin{frame}
\frametitle{Continuation Passing Style : transformation des termes}

{\footnotesize

\itempoint On définit maintenant la transformation $t \mapsto
t^{\cps}$ sur les termes du $\lambda$CST de manière à ce que si
$\vdash t:P$ alors $\vdash t^{\cps}:P^{\cps}$ (rappel :
$\mathop{\sim}P := (P\Rightarrow Z)$ et $P^{\cps} :=
\mathop{\sim}\mathop{\sim}P^\diamond$) :
\begin{itemize}
\item $v^{\cps} = \lambda(k:\mathop{\sim}P^\diamond).\,kv$ lorsque $v$
  est une variable de type $P$,
\item $(fx)^{\cps} = \lambda (k:\mathop{\sim}Q^\diamond).\,
  f^{\cps}(\lambda (f_0:P^\diamond\Rightarrow Q^{\cps}).\,
  x^{\cps}(\lambda (x_0:P^\diamond).\, f_0 x_0 k))$ lorsque $f :
  P\Rightarrow Q$ et $x : Q$ (rappel : $(P\Rightarrow Q)^{\cps} =
  \mathop{\sim}\mathop{\sim}(P^\diamond \Rightarrow Q^{\cps})$),
\item $(\lambda (v:P).t)^{\cps} = \lambda k.\, k(\lambda
  (v:P^\diamond).\, t^{\cps})$ lorsque $\Gamma, v:P \vdash t:Q$,
\item $\langle x,y\rangle^{\cps} = \lambda
  (k:\mathop{\sim}(Q_1^\diamond \land Q_2^\diamond)).\,
  x^{\cps}(\lambda (x_0:Q_1^\diamond).\, y^{\cps} (\lambda
  (y_0:Q_2^\diamond).\, k\langle x_0,y_0\rangle))$,
\item $(\pi_i z)^{\cps} = \lambda (k:\mathop{\sim} Q_i^\diamond).\,
  z^{\cps}(\lambda (z_0:Q_1^\diamond\land Q_2^\diamond).\, k(\pi_i
  z_0))$ (pour $i\in\{1,2\}$),
\item $(\iota^{(Q_1,Q_2)}_i z)^{\cps} = \lambda
  (k:\mathop{\sim}(Q_1^\diamond\lor Q_2^\diamond)).\, z^{\cps}(\lambda
  (z_0:Q_i^\diamond).\, k(\iota^{(Q_1^\diamond,Q_2^\diamond)}_i z_0))$
  (pour $i\in\{1,2\}$),
\item $(\texttt{match~}r\texttt{~with~}\iota_1 v_1 \mapsto t_1,\;
  \iota_2 v_2 \mapsto t_2)^{\cps} = \lambda
  (k:\mathop{\sim}Q^\diamond).\, r^{\cps}(\lambda
  (r_0:P_1^\diamond\lor P_2^\diamond).\penalty0\,
  (\texttt{match}\penalty500\ r_0\ \texttt{with}\ \iota_1 v_1 \mapsto
  t_1^{\cps} k,\; \iota_2 v_2 \mapsto t_2^{\cps} k))$
\item $\bullet^{\cps} = \lambda (k:\mathop{\sim}\top).k{\bullet}$
\item $(\texttt{exfalso}^{(Q)} r)^{\cps} = \lambda
  (k:\mathop{\sim}Q^\diamond).\, r^{\cps}(\lambda (r_0:\bot).\,
  (\texttt{exfalso}^{(Q^\diamond)} r_0))$
\end{itemize}

\itempoint $\texttt{callcc}^{\cps} = \lambda (\ell:\cdots).\,
\ell(\lambda (g : (P^\diamond \Rightarrow Q^{\cps}) \Rightarrow
P^{\cps}).\, \lambda(k:\mathop{\sim}P^\diamond).\,
g(\lambda(v:P^\diamond).\, \lambda(k_0:\mathop{\sim}Q^\diamond).\,
kv)\,k)$, de type $(((P\Rightarrow Q)\Rightarrow P)\Rightarrow
P)^{\cps}$.

\par}

\end{frame}
%
\begin{frame}
\frametitle{Continuation Passing Style : remarques informatiques}

\itempoint Programmer en CPS fait disparaître l'utilisation de la
pile : \alert{tous} les appels de fonctions deviennent terminaux, donc
seront remplacés par des sauts dans un langage supportant la récursion
terminale propre.

\medskip

\itempoint La conversion vers CPS fixe l'ordre d'évaluation des
valeurs.  Par exemple, on a défini $\langle x,y\rangle^{\cps} =
\lambda k.\, x^{\cps}(\lambda x_0.\, y^{\cps} (\lambda y_0.\, k\langle
x_0,y_0\rangle))$ qui évalue $x$ avant $y$, mais on pouvait aussi
$\lambda k.\, y^{\cps} (\lambda y_0.\, x^{\cps}(\lambda x_0.\,
k\langle x_0,y_0\rangle))$ pour inverser l'ordre.

\medskip

{\footnotesize

\itempoint La conversion vers CPS fixe même une stratégie
d'évaluation : on a choisi « call-by-value » ici, mais on aurait pu
prendre « call-by-name » :

\begin{tabular}{c|c}
« call-by-value »&« call-by-name »\\\hline
$v^{\cps} = \lambda k.kv$&$v^{\cps} = \lambda k.vk$ (ou juste $v$)\\
&mais pour une \emph{valeur} : $x^{\cps} = \lambda k.kx$\\
$(\lambda v.t)^{\cps} = \lambda k.\, k(\lambda v.\, t^{\cps})$
&$(\lambda v.t)^{\cps} = \lambda k.\, k(\lambda v.\, t^{\cps})$\\
$(fx)^{\cps} = \lambda
  k.\,f^{\cps}(\lambda f_0.\, x^{\cps}(\lambda x_0.\, f_0 x_0 k))$
&$(fx)^{\cps} = \lambda
  k.\,f^{\cps}(\lambda f_0.\, f_0 x^{\cps} k)$\\
$(P\Rightarrow Q)^\diamond = (P^\diamond\Rightarrow \mathop{\sim}\mathop{\sim}Q^\diamond)$
&$(P\Rightarrow Q)^\diamond = (\mathop{\sim}\mathop{\sim}P^\diamond\Rightarrow \mathop{\sim}\mathop{\sim}Q^\diamond)$
\end{tabular}

\par}

\end{frame}
%
\begin{frame}
\frametitle{Continuation Passing Style : remarques informatiques (suite)}

\itempoint Le CPS donne le call/cc « pour rien ».  En fait, en style
CPS, chaque fonction a contrôle complet sur l'exécution de tout le
programme (il n'y a plus de pile !).

\medskip

\itempoint Le CPS peut servir de point de départ pour la compilation
(Appel, \textit{Compiling with Continuations}, 1992).

\medskip

\itempoint On peut concevoir un langage où chaque fonction a
p.ex. \alert{deux} continuations : une « continuation de succès » et
une « continuation d'échec », qu'on chaîne avec des opérateurs « et »
et « ou » : c'est essentiellement le Prolog (où ces opérateurs sont
notés ‘\texttt{,}’ et ‘\texttt{;}’).

\medskip

\itempoint Le CPS est utile pour séquencer explicit\textsuperscript{t}
les évaluations.  Voir notamment les promesses de JavaScript, et la
monade \texttt{Control.Monad.Cont} de Haskell.

\end{frame}
%
\begin{frame}
\frametitle{Continuation Passing Style : remarques logiques}

\itempoint Par la transformation $t \mapsto t^{\cps}$ et le fait
d'avoir trouvé un terme de type $(((P\Rightarrow Q)\Rightarrow
P)\Rightarrow P)^{\cps}$, on a montré que :
\[
\text{si~}\mathsf{CPC} \vdash P\text{~alors~}\mathsf{IPC} \vdash P^{\cps}
\]
où « $\mathsf{IPC} \vdash P$ » signifie « $P$ est prouvable en logique
intuitionniste » et « $\mathsf{CPC} \vdash P$ » signifie « $P$ est
prouvable en logique classique ».

\medskip

\itempoint \alert{Si on prend $Z := \bot$}, la proposition $P^{\cps}$
ajoute simplement des $\neg\neg$ un peu partout, ce qui est équivalent
en logique classique, donc on a évidemment
\[
\mathsf{CPC} \vdash P\text{~ssi~}\mathsf{CPC} \vdash P^{\cps}
\]
Comme par ailleurs $\mathsf{CPC} \vdash Q$ implique $\mathsf{IPC} \vdash
Q$, on déduit des deux affirmations ci-dessus que (pour $Z=\bot$) :
\[
\mathsf{CPC} \vdash P\text{~\alert{ssi}~}\mathsf{IPC} \vdash P^{\cps}
\]
On dit qu'on a « \alert{interprété} » la logique (propositionnelle)
classique en logique (propositionnelle) intuitionniste par la
traduction « double négation ».

\smallskip

{\footnotesize

\itempoint En fait, on peut se contenter de mettre un $\neg\neg$
devant les disjonctions et formules atomiques (traduction de
Gödel-Gentzen), ou bien (en calcul \alert{propositionnel} !)  d'en
mettre un seul devant toute la formule (traduction de Glivenko).

\par}

\end{frame}
%
\section{Sémantique du calcul propositionnel intuitionniste}
\begin{frame}
\frametitle{Qu'est-ce qu'une « sémantique » en logique ?}

De façon très (trop ?) vague :

\smallskip

\itempoint Une \textbf{logique} $\mathsf{L}$ est définie par un
certain nombre de règles (axiomes, modes d'inférence)
\alert{syntactique} opérant sur des « formules » et définit une notion
de \textbf{théorème} comme les formules ayant une \textbf{preuve}
selon ces règles.  On note (qqch comme) $\mathsf{L} \vdash \varphi$
pour « $\varphi$ est un théorème [= est démontrable]
selon $\mathsf{L}$ ».

\medskip

\itempoint Un \textbf{modèle} pour $\mathsf{L}$ est une structure
mathématique générale qui donne un « sens » au symboles utilisés dans
$\mathsf{L}$ et qui notamment déclare certaines formules comme
\textbf{vraies} (noté $\mathcal{M} \vDash \varphi$).  Une
\textbf{sémantique} est un ensemble de modèles (formés sur le même
schéma).  On dit que la sémantique $\mathsf{S}$ \textbf{valide} une
formule $\varphi$ lorsque $\varphi$ est vraie dans chacun de ses
modèles (on note $\mathsf{S} \vDash \varphi$).

\medskip

On dit que
\begin{itemize}
\item $\mathsf{S}$ est \textbf{correcte} pour $\mathsf{L}$ lorsque
  $\mathsf{L} \vdash \varphi$ implique $\mathsf{S} \vDash \varphi$
  (« tout théorème est vrai dans tout modèle »).  On veut toujours
  ça !
\item $\mathsf{S}$ est \textbf{complète} pour $\mathsf{L}$ lorsque
  $\mathsf{S} \vDash \varphi$ implique $\mathsf{L} \vdash \varphi$
  (« toute formule vraie dans tout modèle est un théorème »).
\end{itemize}

\end{frame}
%
\begin{frame}
\frametitle{Exemples de sémantiques}

\itempoint La sémantique des \textbf{tableaux de vérité booléens} est
\alert{correcte et complète} pour le calcul propositionnel
\alert{classique} : $\varphi$ est démontrable \alert{ssi} toute
affectation de « vrai » ou « faux » à chacune de ses variables donne
« vrai ».

\smallskip

Pour le calcul propositionnel \alert{intuitionniste}, elle n'est
\alert{que correcte}.

\medskip

\itempoint Le plan réel $\mathbb{R}^2$ est un modèle des axiomes de la
géométrie euclidienne.  Ce modèle est une sémantique correcte et
complète à lui seul : un énoncé géométrique est démontrable à partir
des axiomes \alert{ssi} il est vrai dans $\mathbb{R}^2$.

\medskip

\itempoint $\mathbb{N}$ est un modèle de ce qu'on appellera
l'« \textbf{arithmétique de Peano du premier ordre} » $\mathsf{PA}$
(c'est le « modèle souhaité » de $\mathsf{PA}$).  Ce modèle pris tout
seul est une sémantique correcte mais \alert{non complète} (théorème
d'\alert{incomplétude} de Gödel : il y a des énoncés vrais dans
$\mathbb{N}$ mais non démontrables dans $\mathsf{PA}$).

\medskip

\itempoint En revanche, si on élargit les modèles aux « modèles du
premier ordre » (« modèles au sens de Tarski »), la sémantique devient
complète (théorème de \alert{complétude} de Gödel pour la logique du
premier ordre : tout énoncé vrai dans tout modèle est démontrable).

\end{frame}
%
\begin{frame}
\frametitle{Sémantiques du calcul propositionnel intuitionniste ?}

\itempoint On a vu les \alert{règles syntaxiques} du calcul
propositionnel intuitionniste (déduction naturelle, ou calcul des
séquents).  Mais quel est le \textbf{sens} des connecteurs ?

\medskip

\begin{itemize}
\item En logique classique, c'est facile : la sémantique des tableaux
  de vérité est correcte et complète, donc on peut considérer qu'elle
  définit le sens.
\item En logique intuitionniste, on n'a donné pour l'instant qu'une
  interprétation intuitive (Brouwer-Heyting-Kolmogorov) des
  connecteurs.
\end{itemize}

\medskip

\itempoint Avoir une sémantique (correcte) complète, ou « moins
incomplète » que celle des tableaux de vérité permet de prouver qu'une
formule logique n'est \alert{pas démontrable} (si elle n'est pas
validée par cette sémantique).

\medskip

{\footnotesize

\itempoint On peut considérer \alert{Curry-Howard} comme une
sémantique : le « sens » de $\Rightarrow,\land,\lor,\top,\bot$ est
donné par les opérations sur les types d'un langage de programmation,
et les énoncés validés par le modèle sont ceux dont le type est
habité.  Elle est complète si le langage est le $\lambda$CST.  Mais
elle est peu maniable et/ou un peu triviale !

\par}

\end{frame}
%
\begin{frame}
\label{truth-table-semantics}
\frametitle{Sémantique 0 : tableaux de vérité}

{\footnotesize (Reprise de ce qui a déjà été dit.)\par}

\medskip

\itempoint On définit $\mathbb{B} := \{0,1\}$, $\dottedtop = 1$,
$\dottedbot = 0$ et $\dottedland, \dottedlor, \dottedlimp \colon
\mathbb{B}^2 \to \mathbb{B}$ par les tableaux de vérité usuels :

\smallskip

\begin{tabular}{c@{\hskip 30bp}c@{\hskip 30bp}c}
$
\begin{array}{c|cc}
\dottedland&0&1\\\hline
0&0&0\\
1&0&1\\
\end{array}
$
&
$
\begin{array}{c|cc}
\dottedlor&0&1\\\hline
0&0&1\\
1&1&1\\
\end{array}
$
&
$
\begin{array}{c|cc}
A \dottedlimp B&B=0&B=1\\\hline
A=0&1&1\\
A=1&0&1\\
\end{array}
$
\end{tabular}

\medskip

\itempoint Si $\varphi$ est une formule propositionnelle en $r$
variables, ceci définit une fonction $\dot\varphi \colon \mathbb{B}^r
\to \mathbb{B}$ par composition.

\medskip

\itempoint Un \textbf{modèle booléen} du calcul propositionnel est une
affectation $\{\mathrm{variables}\} \to \mathbb{B}$ : chaque formule
$\varphi$ a donc une valeur de vérité ($0$ ou $1$) dans le modèle,
donnée par $\dot\varphi$.  On dit que $\mathcal{M} \vDash \varphi$
lorsque c'est $1$.

\medskip

\itempoint On dit que la sémantique booléenne valide $\varphi$ lorsque
$\mathcal{M} \vDash \varphi$ pour tout modèle (i.e., $\dot\varphi$
vaut constamment $1$).

\end{frame}
%
\begin{frame}
\frametitle{Correction et complétude classique des tableaux de vérité}

\textbf{Théorème :} la sémantique booléenne est \alert{correcte et
  complète} pour le calcul propositionnel classique.

\medskip

\underline{Esquisse de preuve :}

\smallskip

\itempoint\underline{Correction :} on montre par induction sur la
preuve que si $\eta_1,\ldots,\eta_r \vdash \varphi$ alors $\dot\varphi
= 1$ dans tout modèle où $\dot\eta_1 = \cdots = \dot\eta_r = 1$ : la
vérification est très facile sur chaque règle du calcul
propositionnel.

\medskip

\itempoint\underline{Complétude :} on démontre \alert{classiquement}
$\vdash A_i \lor \neg A_i$ pour chaque variable $A_i$, puis on utilise
l'élimination du $\lor$ pour se placer dans chacun des $2^r$ cas
possibles (i.e., avec $A_i$ ou $\neg A_i$ dans les hypothèses).  Dans
chaque cas on a $A_i \Leftrightarrow \top$ ou $A_i \Leftrightarrow
\bot$ pour chaque variable, donc en suivant les tableaux de vérité on
a $\varphi$ (vraie).  L'élimination du $\lor$ montre alors $\varphi$
vraie.\qed

\bigskip

\alert{Intuitionnistement}, la correction vaut toujours, mais plus la
complétude ($A\lor\neg A$ n'est pas démontrable).

\end{frame}
%
\begin{frame}
\frametitle{Sémantique 1 : cadres de Kripke}

\textbf{Définitions :}
\begin{itemize}
\item un \textbf{cadre de Kripke} est {\footnotesize (dans ce
  contexte)} un ensemble partiellement ordonné $(W,\leq)$ ; les
  éléments de $W$ s'appellent les \textbf{mondes} ; si $w\leq w'$, on
  dit que $w'$ est \textbf{accessible} depuis $w$ ;
\item dans un cadre de Kripke, une \textbf{affectation de vérité} est
  une fonction $p\colon W \to \mathbb{B}$ (où $\mathbb{B} = \{0,1\}$)
  telle que si $p(w) = 1$ et $w\leq w'$ alors $p(w') = 1$ (i.e., $p$
  est croissante, on dit aussi « permanente ») ;
\item on définit des opérations $\dottedland, \dottedlor, \dottedlimp$
  sur les affectations de vérité :
\begin{itemize}
\item $(q_1\dottedland q_2)(w) = 1$ ssi $q_1(w) = 1$ et $q_2(w) = 1$ ;
\item $(q_1\dottedlor q_2)(w) = 1$ ssi $q_1(w) = 1$ ou $q_2(w) = 1$ ;
\item $(q_1\dottedlimp q_2)(w) = 1$ ssi pour tout $w' \geq w$ t.q.
  $q_1(w') = 1$, on a $q_2(w') = 1$ ;
\item $\dottedtop(w) = 1$ partout ;
\quad\itempoint $\dottedbot(w) = 0$ partout.
\end{itemize}
\item un \textbf{modèle de Kripke} est {\footnotesize (dans ce
  contexte)} un cadre de Kripke $(W,\leq)$ muni d'une affectation de
  vérité pour chaque variable propositionnelle ;
\item $\varphi$ est \textbf{validée} par le modèle de Kripke lorsque
  $\dot\varphi$ vaut $1$ dans tout monde $w$.
\end{itemize}

\end{frame}
%
\begin{frame}
\frametitle{Cadres de Kripke (suite)}

{\footnotesize Recopie : \itempoint $(q_1\dottedland q_2)(w) = 1$ ssi
  $q_1(w) = 1$ et $q_2(w) = 1$ ;\quad\itempoint $(q_1\dottedlor
  q_2)(w) = 1$ ssi $q_1(w) = 1$ ou $q_2(w) = 1$ ;\quad\itempoint
  $(q_1\dottedlimp q_2)(w) = 1$ ssi pour tout $w' \geq w$ t.q.
  $q_1(w') = 1$, on a $q_2(w') = 1$.\par}

\medskip

Par exemple :
\begin{itemize}
\item $\dottedneg p$ (c'est-à-dire $p\dottedlimp \dottedbot$) vaut $1$
  dans le monde $w$ lorsque $p$ vaut $0$ dans \alert{tout monde
    accessible} depuis $w$ ;
\item $p\dottedlor\dottedneg p$ vaut $1$ dans le monde $w$ lorsque $p$
  vaut $0$ dans \alert{tout monde accessible} depuis $w$ ou bien $1$
  dans $w$ (donc dans tout monde accessible depuis $w$, par
  permanence) : \textcolor{teal}{$p$ est « décidé » à vrai ou à
    faux} ;
\item $\dottedneg \dottedneg p$ vaut $1$ dans le monde $w$ lorsque
  $\forall w'\geq w.\, \exists w''\geq w'.\, p(w'')=1$
  \textcolor{teal}{(« dans tout futur possible, $p$ finit par devenir
    vrai »)}.
\end{itemize}

\bigskip

\itempoint La sémantique de Kripke modélise plus ou moins l'intuition
selon laquelle « $A \Rightarrow B$ » ne signifie pas juste « soit $A$
est vrai soit $B$ est faux » (sens en logique classique) mais bien
« dans tout monde possible où $A$ est vrai, $B$ l'est ».

\end{frame}
%
\begin{frame}
\frametitle{Sémantique des cadres de Kripke}

\textbf{Définitions} (suite) :
\begin{itemize}
\item {\footnotesize (déjà dit :)} $\varphi$ est \textbf{validée} par
  le \alert{modèle} de Kripke (= valide dedans, = valable) lorsque
  $\dot\varphi$ vaut $1$ dans tout monde $w$ ;
\item $\varphi$ est \textbf{validée} par le \alert{cadre} de Kripke
  lorsqu'elle est validée par toute affectation de vérité pour chacune
  de ses variables (= tout modèle sur ce cadre) ;
\item $\varphi$ est \textbf{validée} par \alert{la sémantique de
  Kripke} lorsque $\varphi$ est validée par \alert{tout} cadre (i.e.,
  tout modèle de Kripke).
\end{itemize}

\bigskip

\textbf{Théorème} (Kripke) : la sémantique de Kripke est correcte et
complète pour le calcul propositionnel intuitionniste :
\begin{itemize}
\item\underline{correction :} tout théorème propositionnel
  intuitionniste est valide dans tout modèle de Kripke,
\item\underline{complétude :} toute formule propositionnelle valide
  dans tout modèle de Kripke est démontrable en calcul propositionnel
  intuitionniste.
\end{itemize}

\end{frame}
%
\begin{frame}
\frametitle{Cadres de Kripke : remarques}

La sémantique de Kripke n'est complète que quand on considèe
\alert{tous les cadres}.  Si on se limite aux modèles sur un cadre
donné, ils peuvent valider des formules non démontrables.
\begin{itemize}
\item Le cadre réduit à un singleton valide la logique classique :
  c'est exactement la sémantique booléenne ;
\item Si $(W,\leq)$ est totalement ordonné, il valide la formule
  $(A\Rightarrow B) \lor (B\Rightarrow A)$ {\footnotesize
  (\underline{preuve :} sinon il y a un monde $w$ dans lequel
  $p_A(w)=1$ mais $p_B(w)=0$ et un $w'$ dans lequel $p_A(w')=0$ mais
  $p_B(w')=1$ ; mais si $w\leq w'$ ceci contredit la permanence de
  $p_A$ et si $w\geq w'$ ceci contredit la permanence de
  $p_B$~\qedsymbol)} qui n'est pas prouvable en logique
  propositionnelle intuitionniste.
\item Le cadre le plus simple après un singleton est $\{u,v\}$ avec
  $u\leq v$, qui correspond à une logique à 3 valeurs de vérité, aux
  tableaux suivants.
\end{itemize}

{\footnotesize

\begin{tabular}{c|@{\hskip 15bp}c@{\hskip 15bp}c@{\hskip 15bp}c}
$
\begin{array}{c|cc}
p&u&v\\\hline
0&0&0\\
\mathsf{q}&0&1\\
1&1&1\\
\end{array}
$
&
$
\begin{array}{c|ccc}
\dottedland&0&\mathsf{q}&1\\\hline
0&0&0&0\\
\mathsf{q}&0&\mathsf{q}&\mathsf{q}\\
1&0&\mathsf{q}&1\\
\end{array}
$
&
$
\begin{array}{c|ccc}
\dottedlor&0&\mathsf{q}&1\\\hline
0&0&\mathsf{q}&1\\
\mathsf{q}&\mathsf{q}&\mathsf{q}&1\\
1&1&1&1\\
\end{array}
$
&
$
\begin{array}{c|ccc}
A \dottedlimp B&B=0&B={\mathsf{q}}&B=1\\\hline
A=0&1&1&1\\
A={\mathsf{q}}&0&1&1\\
A=1&0&\mathsf{q}&1\\
\end{array}
$
\end{tabular}

\par}

\end{frame}
%
\begin{frame}
\frametitle{Sémantique 2 : les ouverts en topologie}

Fixons $X$ un espace topologique : si on ne sait pas ce que c'est,
penser à un espace métrique ou simplement $\mathbb{R}^m$ avec sa
topologie ordinaire.  On note $\mathcal{O}(X)$ la topologie de $X$,
c'est-à-dire l'ensemble des ouverts de $X$.

\medskip

On définit les opérations suivantes sur $\mathcal{O}(X)$ :
\begin{itemize}
\item $U \dottedland V := U\cap V$
\quad\itempoint $U \dottedlor V := U\cup V$
\quad\itempoint $\dottedtop := X$
\quad\itempoint $\dottedbot := \varnothing$
\item $(U \dottedlimp V) := \operatorname{int}((X\setminus U) \cup V)$
  (où $\operatorname{int}$ désigne l'intérieur) est l'ensemble des
  points $x \in X$ tels qu'il y ait un ouvert $W \ni x$ t.q.  $(U\cap
  W) \subseteq (V\cap W)$ {\footnotesize (= les points « au voisinage
    desquels » $U \subseteq V$)}, ou, si on préfère, le plus grand
  ouvert $W$ tel que $(U\cap W) \subseteq (V\cap W)$ ;
\item notamment, $\dottedneg U := \operatorname{int}(X\setminus U)$,
  plus grand ouvert disjoint de $U$ ;
\item notamment, $\dottedneg \dottedneg U :=
  \operatorname{int}(\operatorname{adh}(U))$ (intérieur de l'adhérence
  de $U$, ou « régularisé » de $U$).
\end{itemize}

\medskip

\itempoint Si $\varphi(A_1,\ldots,A_r)$ est une formule
propositionnelle et $U_1,\ldots,U_r \in \mathcal{O}(X)$, on définit
$\dot\varphi(U_1,\ldots,U_r)$ de façon évidente (par induction).

\end{frame}
%
\begin{frame}
\frametitle{Sémantique des ouverts : correction et complétude}

De façon surprenante, les opérations qu'on a définies sur les ouverts
de $X$ fournissent un modèle du calcul propositionnel intuitionniste :

\medskip

\textbf{Théorème} (Tarski) : la sémantique des ouverts est correcte et
complète pour le calcul propositionnel intuitionniste :
\begin{itemize}
\item\underline{correction :} si $\varphi(A_1,\ldots,A_r)$ est un
  théorème du calcul propositionnel intuitionniste, alors quel que
  soit $X$ et $U_1,\ldots,U_r$ ouverts de $X$, on a
  $\dot\varphi(U_1,\ldots,U_r) = \dottedtop$ (l'espace tout entier) ;
\item\underline{complétude :} réciproquement, si
  $\dot\varphi(U_1,\ldots,U_r) = \dottedtop$ pour tous ouverts
  $U_1,\ldots,U_r$ de tout espace topologique $X$, \alert{ou même
  simplement de $\mathbb{R}^n$} pour un $n \geq 1$ donné, alors
  $\varphi$ est démontrable en calcul propositionnel intuitionniste.
\end{itemize}

\medskip

\itempoint La sémantique des ouverts modélise plus ou moins
l'intuition que la vérité est une « notion locale » (si quelque chose
est vrai en un point, c'est vrai autour de ce point).

\end{frame}
%
\begin{frame}
\frametitle{Sémantique des ouverts : exemples}

\itempoint Soit $U = \mathopen{]}0,1\mathclose{[}$ dans $X =
    \mathbb{R}$.  Alors
\begin{itemize}
\item $\dottedneg U = \mathopen{]}-\infty,0\mathclose{[} \cup
  \mathopen{]}1,+\infty\mathclose{[}$
\quad\itempoint $\dottedneg\dottedneg U = \mathopen{]}0,1\mathclose{[} = U$
\quad\itempoint $(\dottedneg\dottedneg U \dottedlimp U) = \mathbb{R}$
\item $U \dottedlor \dottedneg U = \mathbb{R}\setminus\{0,1\}$
\quad\itempoint $((\dottedneg\dottedneg U \dottedlimp U) \dottedlimp
(U \dottedlor \dottedneg U)) = \mathbb{R}\setminus\{0,1\}$
\end{itemize}
Par correction, ceci montre que $((\neg\neg A \Rightarrow A)
\Rightarrow (A \lor \neg A))$ n'est pas prouvable en logique
intuitionniste.

\bigskip

{\footnotesize

\itempoint Dans $X = \mathbb{R}^2$, soit $U = \{(x_1,x_2) : x_1<0
\text{~et~} x_2<0\}$ et $V_1 = \{x_1>0\}$ et $V_2 = \{x_2>0$\}.
\begin{itemize}
\item $\dottedneg U = \{(x_1,x_2) : x_1>0 \text{~ou~} x_2>0\} =
  V_1\dottedlor V_2$
\quad\itempoint Donc $(\dottedneg U \dottedlimp (V_1\dottedlor V_2)) =
\mathbb{R}^2$
\item $(\dottedneg U \dottedlimp V_1) = \{(x_1,x_2) : x_1<0
  \text{~ou~} x_2>0\}$
\item $(\dottedneg U \dottedlimp V_2) = \{(x_1,x_2) : x_1>0
  \text{~ou~} x_2<0\}$
\item $(\dottedneg U \dottedlimp V_1) \dottedlor (\dottedneg U
  \dottedlimp V_2) = \mathbb{R}^2 \setminus\{(0,0)\}$
\end{itemize}
Ceci montre que $(\neg A \Rightarrow (B_1\lor B_2)) \Rightarrow (\neg
A \Rightarrow B_1) \lor (\neg A \Rightarrow B_2)$ (« axiome de
Kreisel-Putnam ») n'est pas prouvable en logique intuitionniste (et
\textit{a fortiori} $(C \Rightarrow B_1) \lor (C \Rightarrow B_2)
\Rightarrow (C \Rightarrow (B_1\lor B_2))$ ne l'est pas).

\par}

\end{frame}
%
\begin{frame}
\frametitle{Sémantique 3 : la réalisabilité propositionnelle}

On reprend les notations de la calculabilité, mais on notera $\Phi_e
\colon \mathbb{N} \dasharrow \mathbb{N}$ pour la $e$-ième fonction
générale récursive (pour éviter un conflit de notation).

\medskip

On définit les opérations suivantes sur l'ensemble
$\mathcal{P}(\mathbb{N})$ des parties de $\mathbb{N}$ :
\begin{itemize}
\item $P \dottedland Q = \{\langle m,n\rangle : m\in P, n\in Q\}$
  (codage de Gödel du produit $P\times Q$)
\item $P \dottedlor Q = \{\langle 0,m\rangle : m\in P\} \cup \{\langle
  1,n\rangle : n\in Q\}$ (sorte de réunion disjointe)
\item $(P \dottedlimp Q) = \{e \in \mathbb{N} : \Phi_e(P){\downarrow}
  \subseteq Q\}$, ce qui signifie $\forall m\in P.\,
  \Phi_e(m){\downarrow}\in Q$ (programmes définis sur tout $P$ et
  l'envoyant dans $Q$)
\item $\dottedtop = \mathbb{N}$\quad\itempoint $\dottedbot = \varnothing$
\item notamment, $\dottedneg P$ vaut $\mathbb{N}$ si $P = \varnothing$
  et $\varnothing$ sinon (« promesses » que $P=\varnothing$) ;
\item notamment, $\dottedneg \dottedneg P$ vaut $\varnothing$ si $P =
  \varnothing$ et $\mathbb{N}$ sinon (« promesses » que
  $P\neq\varnothing$).
\end{itemize}

\medskip

\itempoint Si $\varphi(A_1,\ldots,A_r)$ est une formule
propositionnelle et $P_1,\ldots,P_r \subseteq \mathbb{N}$, on définit
$\dot\varphi(P_1,\ldots,P_r)$ de façon évidente (par induction).

\smallskip

\itempoint Si $n \in \dot\varphi(P_1,\ldots,P_r)$, on dit aussi que
$n$ \textbf{réalise} $\varphi(P_1,\ldots,P_r)$.

\end{frame}
%
\begin{frame}
\frametitle{La réalisabilité propositionnelle}

\itempoint Les définitions de $\dottedlimp,\dottedland,\dottedlor$
suivent précisément l'idée informelle de l'interprétation de
Brouwer-Heyting-Kolmogorov (transp. \ref{bhk-interpretation}) en les
rendant précises avec des parties de $\mathbb{N}$ et des fonctions
générales récursives.

\medskip

\itempoint On dit qu'une formule propositionnelle
$\varphi(A_1,\ldots,A_r)$ est \textbf{réalisable} (plus précisément,
« uniformément réalisable ») lorsqu'il existe un \alert{même entier}
$n$ qui réalise $\varphi(P_1,\ldots,P_r)$ quels que soient
$P_1,\ldots,P_r \subseteq \mathbb{N}$ :
\[
n \in \bigcap_{P_1,\ldots,P_r} \dot\varphi(P_1,\ldots,P_r)
\]

\medskip

\itempoint\textbf{Théorème} (Dragalin, Troelstra, Nelson) : la
sémantique de la réalisabilité est correcte pour le calcul
propositionnel intuitionniste : si $\varphi(A_1,\ldots,A_r)$ est un
théorème du calcul propositionnel intuitionniste, alors il existe $n$
qui réalise $\varphi(P_1,\ldots,P_r)$ quels que soient $P_1,\ldots,P_r
\subseteq \mathbb{N}$.

\medskip

{\footnotesize Mieux : ce $n$ se construit à partir de la preuve de
  $\varphi$ : c'est une forme d'\alert{extraction d'algorithme}.}

\smallskip

\textcolor{brown}{En fait on peut voir ça comme une conséquence de
  Curry-Howard.}

\end{frame}
%
\begin{frame}
\frametitle{La réalisabilité propositionnelle : exemple}

\itempoint Essayons de réaliser $A\land B \Rightarrow B\land A$ : on
cherche donc un même entier dans $P\dottedland Q \dottedlimp
Q\dottedland P$ pour \emph{tous} $P,Q\subseteq\mathbb{N}$.

\medskip

\itempoint Autrement dit, on veut trouver un programme $e$ tel que
$\Phi_e$ soit défini sur $P\dottedland Q = \{\langle m,n\rangle : m\in
P, n\in Q\}$ et l'envoie dans $Q\dottedland P = \{\langle n,m\rangle :
n\in Q, m\in P\}$ (sans connaître $P,Q$).

\smallskip

Il suffit de permuter les coordonnées du couple !

\smallskip

Mais c'est ce que faisait le programme $\lambda z.\langle \pi_2 z,
\pi_1 z\rangle$ associé à la preuve via Curry-Howard
(cf. transp. \ref{curry-howard-example-with-conjunction}) !

\medskip

\itempoint Ceci marche en général : pour réaliser $\varphi$ où
$\varphi$ est prouvable, on prend le $\lambda$-terme de preuve, on
oublie les types (tout est entier !), on interprète l'application
d'une fonction $f$ sur $n$ comme $\Phi_f(n)$, les couples et les
sommes comme dans $\dottedland$ et $\dottedlor$, et le programme en
question réalise $\varphi$ quels que soient $P_1,\ldots,P_r$
puisqu'ils correspondent à des types dans le terme de preuve.

\end{frame}
%
\begin{frame}
\frametitle{La réalisabilité propositionnelle : contre-exemples}

\itempoint La formule $A\lor\neg A$ n'est pas réalisable : il
s'agirait de trouver un \alert{même} entier qui \alert{quel que soit
  $P$} soit de la forme $\langle 0,n\rangle$ avec $n\in P$ ou bien de
la forme $\langle 1,n\rangle$ si $P=\varnothing$.  Visiblement c'est
impossible (sans information sur $P$) !

\smallskip

{\footnotesize\textbf{Remarque :} en fait, si $\varphi\lor\psi$ est
  réalisable, l'une de $\varphi$ ou $\psi$ l'est (selon que l'entier
  réalisant $\varphi\lor\psi$ est de la forme $\langle 0,m\rangle$ ou
  $\langle 1,m\rangle$).\par}

\medskip

\itempoint La formule $\neg\neg A \Rightarrow A$ n'est pas
réalisable : il s'agirait de trouver un \alert{même} programme qui, si
$P$ est non-vide (si bien que $\dottedneg\dottedneg P = \mathbb{N}$),
termine sur n'importe quel entier et renvoie un élément de $P$.
Visiblement c'est impossible (sans information sur $P$) !

\medskip

\itempoint Plus subtilement, $(\neg\neg A \Rightarrow A) \Rightarrow
(A\lor\neg A)$ n'est pas réalisable : il s'agirait de trouver un
programme qui transforme une solution du 2\textsuperscript{e} problème
en une solution du 1\textsuperscript{er}.

\medskip

De nouveau, ceci montre que ces formules ne sont pas prouvables.

\medskip

{\footnotesize En revanche, $(A\lor\neg A) \Rightarrow (\neg\neg A
  \Rightarrow A)$ est réalisable car démontrable : le $\lambda$-terme
  $\lambda (z:A\lor\neg A). \lambda(u:\neg\neg A).\,
  (\texttt{match}\ z\ \texttt{with}\ \iota_1 v_1 \mapsto v_1,\;
  \iota_2 v_2 \mapsto \texttt{exfalso}^{(A)}(u v_2))$ le prouve et
  explique comment le réaliser.\par}

\end{frame}
%
\begin{frame}
\frametitle{La réalisabilité propositionnelle : incomplétude}

La réalisabilité suit tellement près les idées de Curry-Howard qu'on
pourrait naturellement croire que la réciproque est vraie, que toute
proposition réalisable donnera un $\lambda$-terme de preuve, i.e., que
la sémantique est complète.

\medskip

\alert{Surprise : non !}

\medskip

On connaît des formules propositionnelles, p.ex. (transp. suivant) :
\[
\begin{aligned}
&\Big(\neg (A_1 \land A_2) \land (\neg A_1 \Rightarrow (B_1 \lor B_2))
  \land (\neg A_2 \Rightarrow (B_1 \lor B_2))\Big)\\
\Rightarrow\,& \Big((\neg A_1 \Rightarrow B_1)\lor(\neg A_2 \Rightarrow B_1)\lor(\neg
  A_1 \Rightarrow B_2)\lor(\neg A_2 \Rightarrow B_2)\Big)
\end{aligned}
\]
qui sont réalisables mais non prouvables en logique intuitionniste.

\smallskip

{\footnotesize C'est-à-dire que bien qu'on ait un algorithme
  (transp. suivant) qui réalise cette formule pour toutes parties
  $A_1,A_2,B_1,B_2$ de $\mathbb{N}$ (et moralement : toutes données),
  on ne peut pas typer un tel algorithme, lequel dépend de la
  possibilité de faire temporairement des « fausses promesses » pour
  donner finalement un résultat juste.\par}

\medskip

\textbf{Problèmes ouverts :} l'ensemble des formules réalisables
est-il décidable ?  Semi-décidable ?

\end{frame}
%
\begin{frame}
\frametitle{Réalisabilité de la formule de Tseitin}

{\footnotesize\textcolor{gray}{Ceci est une digression mais je pense
    que c'est très instructif pour la calculabilité de comprendre la
    différence entre typage et réalisabilité.}\par}

\smallskip

La formule suivant, bien que non démontrable, est réalisable :
\[
\begin{aligned}
&(\neg (A_1 \land A_2) \land (\neg A_1 \Rightarrow (B_1 \lor B_2))
  \land (\neg A_2 \Rightarrow (B_1 \lor B_2)))\\
\Rightarrow\,& ((\neg A_1 \Rightarrow B_1)\lor(\neg A_2 \Rightarrow B_1)\lor(\neg
  A_1 \Rightarrow B_2)\lor(\neg A_2 \Rightarrow B_2))
\end{aligned}
\]

{\footnotesize

\textbf{Algorithme} (appelons $A_1,A_2,B_1,B_2$ les parties
concernées) :

\smallskip

\itempoint On reçoit en entrée une promesse que l'un de $A_1$ et $A_2$
est vide, ainsi que deux algorithmes, l'un $e_1 \in (\dottedneg A_1
\dottedlimp (B_1 \dottedlor B_2))$ prenant en entrée une promesse que
$A_1=\varnothing$ et renvoyant un élément de $B_1$ ou $B_2$, et
l'autre $e_2$ prennant promesse $A_2=\varnothing$ et renvoyant un
él\textsuperscript{t} de $B_1$ ou $B_2$.

\smallskip

\itempoint On \alert{lance en parallèle} $e_1$ resp. $e_2$ sur un
entier qcque promettant (peut-être faussement !)
$A_1=\varnothing$ resp. $A_2=\varnothing$.  Au moins l'une de ces
promesses est vraie (par la promesse en entrée) donc l'un de $e_1$ ou
$e_2$ va terminer, et renvoyer soit un élément annoncé de $B_1$ soit
de $B_2$.

\smallskip

Disons sans perte de généralité que $e_1$ renvoie un élément $n$
annoncé de $B_1$.

\smallskip

\itempoint On renvoie alors comme élément de $\dottedneg A_1
\dottedlimp B_1$ un programme renvoyant constamment $n$.  Il est bien
dans $\dottedneg A_1 \dottedlimp B_1$ car s'il reçoit une promesse que
$A_1=\varnothing$, à l'étape précédente le programme a tourné sur une
vraie promesse et donc a vraiment renvoyé un élément de $B_1$.

\par}

\end{frame}
%
\begin{frame}
\frametitle{Sémantique 4 : les problèmes finis}

\itempoint Un \textbf{problème fini} est un couple $(X,S)$ où $X$ est
un ensemble fini non vide appelé les \textbf{candidats} du problème et
$S \subseteq X$ est un sous-ensemble appelé les \textbf{solutions} du
problème.

\medskip

On définit les opérations suivantes sur les problèmes finis :
\begin{itemize}
\item $(X,S)\dottedland (Y,T) = (X\times Y, S\times T)$ (produit cartésien)
\item $(X,S)\dottedlor (Y,T) = (X\uplus Y, S\uplus T)$ (réunion disjointe)
\item $(X,S)\dottedlimp (Y,T) = (Y^X, U)$ où $U := \{f\colon X\to Y :
  f(S) \subseteq T\}$ (fonctions envoyant une solution de $(X,S)$ en
  une solution de $(Y,T)$)
\item $\dottedtop = (\{\bullet\},\{\bullet\})$
\quad\itempoint $\dottedbot = (\{\bullet\}, \varnothing)$
\end{itemize}

\medskip

\itempoint Si $\varphi(A_1,\ldots,A_r)$ est une formule
propositionnelle et $(X_1,S_1),\ldots,(X_r,S_r)$ des problèmes finis,
on définit le problème $\dot\varphi((X_1,S_1),\ldots,(X_r,S_r))$ de
façon évidente (par induction).

\end{frame}
%
\begin{frame}
\frametitle{Sémantique des problèmes finis}

\itempoint Si $\varphi(A_1,\ldots,A_r)$ est une formule
propositionnelle et $(X_1,S_1),\ldots,(X_r,S_r)$ des problèmes finis,
noter que l'ensemble $Z$ de candidats de du problème
$\dot\varphi((X_1,S_1),\ldots,(X_r,S_r))$ ne dépend que des ensembles
$X_1,\ldots,X_r$ de candidats donnés.

\medskip

\itempoint On dit que $\varphi$ est Medvedev-valide si pour tous
$X_1,\ldots,X_r$ il existe \alert{un même} $z \in Z$ tel que $z \in
\dot\varphi((X_1,S_1),\ldots,(X_r,S_r))$ quels que soient
$S_1,\ldots,S_r$ (sous-ensembles respectifs de $X_1,\ldots,X_r$).

\medskip

\itempoint Il s'agit d'une autre façon de donner un sens précis à
l'interprétation de Brouwer-Heyting-Kolmogorov, cette fois sur des
ensembles finis (donc pas de question de calculabilité).

\medskip

\itempoint\textbf{Théorème} (Medvedev) : cette sémantique est correcte
pour le calcul propositionnel intuitionniste.

\medskip

{\footnotesize Elle n'est pas complète : p.ex., $(\neg A \Rightarrow
  (B_1\lor B_2)) \Rightarrow (\neg A \Rightarrow B_1) \lor (\neg A
  \Rightarrow B_2)$ (Kreisel-Putnam, déjà évoqué) ou bien $((\neg\neg
  A \Rightarrow A) \Rightarrow (A\lor\neg A)) \Rightarrow (\neg
  A\lor\neg\neg A)$ (Scott) sont Medvedev-valides mais non
  démontrables.\par}

\end{frame}
%
\section{Inférence de type à la Hindley-Milner}
\begin{frame}
\frametitle{L'algorithme de Hindley-Milner : description sommaire}

{\footnotesize Rappel : le $\lambda$-calcul simplement typé
  ($\lambda$CST), pour nous, porte \emph{toutes} les annotations de
  type.\par}

\smallskip

\itempoint Donné un type du $\lambda$CST, on appelle
\textbf{désannotation} de celui-ci le type du $\lambda$-calcul non
typé obtenu en retirant toutes les annotations de type : p.ex.
\[
\lambda(x:\alpha).\,
\lambda(f:((\alpha\to\beta)\to\beta)\to\gamma).\,
f(\lambda(h:\alpha\to\beta).hx)
\]
de type $\alpha \to ((\alpha\to\beta)\to\beta)\to\gamma \to \gamma$ se
désannote en $\lambda xf.\, f(\lambda h.hx)$.

\medskip

\itempoint L'\textbf{inférence de type} pour le $\lambda$CST est le
problème, donné un terme du $\lambda$-calcul non typé, de :
\begin{itemize}
\item décider s'il est typable, i.e., est la désannotation d'un terme
  du $\lambda$CST,
\item si oui, retrouver son type et toutes les annotations, et même
\item trouver la solution la « plus générale » possible (p.ex.,
  $\lambda x.x$ doit retrouver $\lambda(x:\alpha).x$ et pas
  $\lambda(x:\alpha\to\beta).x$).
\end{itemize}

\medskip

\itempoint L'\textbf{algorithme de Hindley-Milner} (ou Damas-Milner,
car Damas a prouvé sa correction) ou \textbf{algorithme W} (ou
\textbf{J}, c'est quasi pareil) accomplit ces buts.

\medskip

\itempoint On fait ainsi des « types opaques » du $\lambda$CST un vrai
polymorphisme.

\end{frame}
%
\begin{frame}
\frametitle{L'algorithme de Hindley-Milner : esquisse}

L'algorithme de Hindley-Milner procède en deux phases :

\medskip

\itempoint\textbf{Première phase : collection des contraintes} :
donnée un terme à typer, on \alert{alloue une variable de type}
fraîches à chaque nouvelle variable introduite ou application et on
\alert{enregistre une contrainte} (équation) à chaque application.
Cette phase n'échoue pas.

\smallskip

{\footnotesize P.ex.: sur $\lambda x.\lambda y.xy$, on va enregistrer
  successivement les types $x:\eta_1$,\quad $y:\eta_2$,\quad
  $xy:\eta_3$ et la contrainte $\eta_1 = (\eta_2\to\eta_3)$, et
  renvoyer $\eta_1 \to \eta_2 \to \eta_3$.\par}

\medskip

\itempoint\textbf{Deuxième phase : résolution des contraintes par
  unification} : on cherche une « substitution » des variables de type
qui résout toutes les contraintes.

\smallskip

Pour ça, on prend de façon répétée une équation $\zeta_1 = \zeta_2$ du
jeu de contraintes et on cherche à la satisfaire : si $\zeta_1$ ou
$\zeta_2$ est une variable, c'est facile, sinon, on \alert{décompose}
chacune ($\zeta_i = (\xi_i \to \chi_i)$) en produisant de nouvelles
contraintes ($\xi_1=\xi_2$ et $\chi_1=\chi_2$).  \alert{Cette phase
  peut échouer.}

\medskip

\itempoint L'algorithme termine, donne une solution s'il en existe
une, et même la « \textbf{solution principale} » dont toutes les
autres se dérivent par substitution.

\end{frame}
%
\begin{frame}
\frametitle{Première phase : collection des contraintes}

La première phase de l'algorithme de H-M fonctionne comme la
vérification de type du $\lambda$CST
(transp. \ref{constructing-the-typing-derivation}) mais en
introduisant des variables de type « fraîches » au vol et en
collectant des contraintes au fur et à mesure.

\medskip

Soit $t$ le terme à typer : l'algo est écrit ici en style impératif et
opère sur un contexte $\Gamma$ et un jeu de contraintes $\mathscr{C}$,
il renvoie un type et modifie $\Gamma,\mathscr{C}$ :
\begin{itemize}
\item si $t = v$ est une variable, elle doit être dans le
  contexte : renvoyer son type ;
\item si $t = \lambda v.e$ est une abstraction, allouer une nouvelle
  variable de type $\eta$, ajouter $v:\eta$ au contexte, appliquer
  l'algorithme récursivement à $e$, qui donne $\zeta$, et renvoyer
  $\eta \to \zeta$ ;
\item si $t = fx$ est une application, appliquer l'algorithme
  récursivement à $f$ et $x$, qui donne $\zeta$ et $\xi$, allouer une
  nouvelle variable de type $\eta$, ajouter l'équation $\zeta =
  (\xi\to\eta)$ aux contraintes et renvoyer $\eta$.
\end{itemize}

\medskip

{\footnotesize\itempoint Ceci s'adapte sans difficulté au cas où on
  dispose d'annotations de type partielles.\par}

\medskip

\itempoint Le typage final s'obtiendra en appliquant la substitution
trouvée dans la deuxième phase.

\end{frame}
%
\begin{frame}
\frametitle{Collection des contraintes : exemples}

\itempoint Soit le terme à typer $\lambda xyz.xz(yz)$.  On collecte
les types et contraintes suivants :\quad
$x:\eta_1$,\quad $y:\eta_2$,\quad $z:\eta_3$,\quad $xz:\eta_4$ avec
$\eta_1 = (\eta_3\to\eta_4)$,\quad $yz:\eta_5$ avec $\eta_2 =
(\eta_3\to\eta_5)$,\quad $xz(yz):\eta_6$ avec $\eta_4 =
(\eta_5\to\eta_6)$,\quad et on renvoie finalement le type $\eta_1 \to
\eta_2 \to \eta_3 \to \eta_6$.

\medskip

\itempoint Soit le terme à typer $\lambda fx.f(\lambda g.gx)$.  On
collecte les types et contraintes suivants :\quad
$f:\eta_1$,\quad $x:\eta_2$,\quad $g:\eta_3$,\quad $gx:\eta_4$ avec
$\eta_3 = (\eta_2\to\eta_4)$,\quad $f(\lambda g.gx):\eta_5$ avec
$\eta_1 = ((\eta_3\to\eta_4) \to \eta_5)$,\quad et on renvoie
finalement le type $\eta_1 \to \eta_2 \to \eta_5$.

\medskip

\itempoint Soit le terme à typer $\lambda x.xx$.  On collecte
les types et contraintes suivants :\quad
$x:\eta_1$,\quad $xx:\eta_2$ avec $\eta_1 = (\eta_1 \to \eta_2)$,\quad
et on renvoie finalement le type $\eta_1 \to \eta_2$.

\medskip

\itempoint Soit le terme à typer $t := (\lambda xy.x) (\lambda
x'y'.x')$.  On collecte les types et contraintes suivants :\quad
$x:\eta_1$,\quad $y:\eta_2$,\quad $x':\eta_3$,\quad $y':\eta_4$\quad
enfin $t:\eta_5$ avec $(\eta_1 \to \eta_2 \to \eta_1) = ((\eta_3 \to
\eta_4 \to \eta_3) \to \eta_5)$.

\smallskip

{\footnotesize \textcolor{orange}{Attention :} ce terme aurait pu être
  écrit $(\lambda xy.x) (\lambda xy.x)$ : les deux $x$ \alert{ne sont
    pas le même} !\par}

\end{frame}
%
\begin{frame}
\frametitle{Seconde phase : résolution des contraintes}

L'algorithme d'\textbf{unification} prend en entrée un ensemble
$\mathscr{C}$ de contraintes (équations $\zeta_1 = \zeta_2$ avec
$\zeta_1,\zeta_2$ deux types) et renvoie une \textbf{substitution} des
variables de type ($\{\eta \mapsto \tau\}$ avec $\eta$ des variables
de type et $\tau$ des types \alert{ne faisant pas intervenir} les
variables substituées) qui vérifie les contraintes, ou « échec » :

\medskip

\itempoint Si $\mathscr{C}$ est vide, renvoyer la substitution vide.

\smallskip

\itempoint Sinon, soit $\zeta_1 = \zeta_2$ une contrainte, et
$\mathscr{C}' := \mathscr{C} \setminus \{(\zeta_1 = \zeta_2)\}$ le
reste des contraintes :
\begin{itemize}
\item si $\zeta_1 = \zeta_2$ déjà, unifier $\mathscr{C}'$,
\item si $\zeta_1$ est une \alert{variable} de type $\eta$ :
  \alert{si} $\zeta_2$ ne fait pas intervenir $\eta$, ajouter (et
  appliquer) $\eta \mapsto \zeta_2$ à la substitution, et unifier
  $\mathscr{C}'$ où $\eta$ est remplacé par $\zeta_2$ ; \alert{sinon},
  \alert{échouer} (« type récursif ») ;
\item si $\zeta_2$ est une variable de type : symétriquement ;
\item sinon, $\zeta_1 = (\xi_1 \to \chi_1)$ et $\zeta_2 = (\xi_2 \to
  \chi_2)$ : unifier $\mathscr{C}'' := \mathscr{C}' \cup \{(\xi_1 =
  \xi_2), (\chi_1 = \chi_2)\}$.
\end{itemize}

\end{frame}
%
\begin{frame}
\frametitle{Résolution des contraintes : exemple}

Reprenons l'exemple $t := (\lambda xy.x) (\lambda x'y'.x')$.  La
première phase a donné : $x:\eta_1$, $y:\eta_2$, $x':\eta_3$,
$y':\eta_4$ et le type final $\eta_5$ avec la constrainte $(\eta_1 \to
\eta_2 \to \eta_1) = ((\eta_3 \to \eta_4 \to \eta_3) \to \eta_5)$.

\smallskip

\begin{itemize}
\item On examine la seule contrainte $(\eta_1 \to \eta_2 \to \eta_1) =
  ((\eta_3 \to \eta_4 \to \eta_3) \to \eta_5)$ : ceci retire cette
  contrainte et on rajoute $\eta_1 = (\eta_3 \to \eta_4 \to \eta_3)$
  et $(\eta_2 \to \eta_1) = \eta_5$.
\item On examine $\eta_1 = (\eta_3 \to \eta_4 \to \eta_3)$ : comme
  $\eta_1$ est une variable, et n'apparaît pas à droite, on enregistre
  la substitution $\eta_1 \mapsto (\eta_3 \to \eta_4 \to \eta_3)$ et
  on l'applique à la contrainte restante qui devient $(\eta_2 \to
  \eta_3 \to \eta_4 \to \eta_3) = \eta_5$.
\item Comme $\eta_5$ est une variable et n'apparaît pas à gauche, on
  enregistre la substitution $\eta_5 \mapsto (\eta_2 \to \eta_3 \to
  \eta_4 \to \eta_3)$ et il ne reste plus de contrainte.
\end{itemize}

\medskip

Finalement, on a typé : $\vdash (\lambda (x:\eta_3 \to \eta_4 \to
\eta_3).\,\lambda(y:\eta_2).\,x) \; (\lambda
(x':\eta_3).\,\lambda(y':\eta_4).\,x') \; : \penalty-1000\; \eta_2 \to
\eta_3 \to \eta_4 \to \eta_3$.

\end{frame}
%
\begin{frame}
\frametitle{Propriétés de l'algorithme de Hindley-Milner}

La difficulté concerne essentiellement la seconde phase (résolution
des contraintes par unification).  On peut prouver :

\medskip

\itempoint\textbf{L'algorithme termine} (il est même p.r. et mieux).

\smallskip

{\footnotesize\underline{Idée de la preuve :} double récurrence, sur
  le nombre de variables de type (récurrence principale) et sur le
  degré total des formules dans les contraintes.\par}

\medskip

\itempoint\textbf{L'algorithme est correct} : s'il retourne une
substitution, celle-ci résout les contraintes. {\footnotesize (C'est à
  peu près évident.)\par}

\medskip

\itempoint (Damas) \textbf{L'algorithme est « complet » :} toute autre
solution des contraintes $\mathscr{C}$ s'obtient en effectuant une
substitution sur celle retournée par l'algorithme : on dit qu'il
trouve la \textbf{solution principale} {\footnotesize (notamment,
  celle-ci existe !)}.

\smallskip

{\footnotesize\underline{Idée de la preuve :} récurrence sur le nombre
  d'appels récursifs à la fonction d'unification.\par}

\medskip

\centerline{*}

\itempoint L'algorithme \alert{se généralise sans problème} à l'ajout
de types produits, sommes, $1$ et $0$ ($\rightarrow$ nvx cas d'échecs
d'unification, p.ex., si $\zeta_1 = ({?} \rightarrow {?})$ et $\zeta_2
= ({?} \times {?})$).

\end{frame}
%
\begin{frame}
\frametitle{Types récursifs ?}

\itempoint On a imposé lors de la résolution de la contrainte $\zeta_1
= \zeta_2$, lorsque $\zeta_1$ est une variable $\eta$, que $\zeta_2$
ne fasse pas intervenir $\eta$.

\medskip

\itempoint Ceci sert à empêcher d'inférer un type p.ex. à $\lambda
x.xx$ (il fallait résoudre $\eta_1 = (\eta_1 \to \eta_2)$).  De fait,
ce terme n'est pas typable dans le $\lambda$CST, sinon $(\lambda x.xx)
(\lambda x.xx)$ le serait, contredisant la normalisation forte du
$\lambda$CST.

\medskip

\itempoint On peut néanmoins étendre l'algorithme de H-M à de tels
\alert{types récursifs sans constructeur}, c'est ce que fait
\texttt{ocaml -rectypes} :

{\tt
\$ ocaml -rectypes\newline
fun x -> x x ;;\newline
- : ('a -> 'b as 'a) -> 'b = <fun>
}

\smallskip

{\footnotesize (Ceci doit faire de OCaml un interpréteur du
  $\lambda$-calcul \alert{non typé}, au moins pour une certaine
  notion d'évaluation.)\par}

\smallskip

Mais c'est une mauvaise idée de s'en servir : mieux vaut définir un
type récursif explicite.

\end{frame}
%
\begin{frame}
\frametitle{Le problème du polymorphisme du « let »}

\itempoint Dans les langages fonctionnels, « \texttt{let $v$=$x$ in
  $t$} » peut être vu comme un sucre syntaxique pour « \texttt{(fun $v
  \mapsto t$)$x$} » (i.e., $(\lambda v.t)x$,
cf. transp. \ref{typing-vs-beta-reduction}).

\medskip

\itempoint Ceci pose un problème au typage à la H-M : mettons qu'on
veuille utiliser l'entier de Church $\overline{2} := \lambda fx
. f(fx)$, typé par H-M comme $(\alpha\to\alpha) \to
(\alpha\to\alpha)$, avec deux types $\alpha$ différents, par exemple
en OCaml :

\smallskip

{\footnotesize\tt
let twice = fun f -> fun x -> f(f x) in (twice ((+)1) 42, twice not true) ;;\newline
{\color{purple}- : int * bool = (44, true)}\newline
(fun twice -> (twice ((+)1) 42, twice not true))(fun f -> fun x -> f(f x)) ;;\newline
{\color{purple}\alert{Error}: This expression has type bool -> bool\newline
       but an expression was expected of type int -> int\newline
       Type bool is not compatible with type int}
\par}

\smallskip

{\footnotesize Dans l'expression \texttt{fun twice -> ...}, le type de
  \texttt{twice} est fixé et ne peut pas être polymorphe.\par}

\medskip

\itempoint Pour réparer ce problème, très grossièrement :
\textcolor{blue}{(a)} on type d'abord $x$, puis
\textcolor{blue}{(b)} on « \alert{généralise} » chaque variable de
type (non présente dans le contexte d'ensemble) pour rendre $x$
polymorphe (cf. \ref{variables-are-effectively-polymorphic}),
c'est-à-dire que \textcolor{blue}{(c)} \alert{chaque} apparition de
$v$ dans $t$ reçoit des variables de type fraîches à ces places.

\end{frame}
%
\begin{frame}
\frametitle{La « restriction de valeur »}

\itempoint La solution trouvée pour rendre \texttt{let} polymorphe
(transp. précédent) apporte ses propres difficultés quand le langage
permet les effets de bord (variables mutables) : considérons
l'expression OCaml :

\smallskip

{\tt let r = ref (fun x -> x) in r := (fun x -> x+1) ; (!r) true ;;}

\smallskip

Ici, \alert{on ne veut pas} généraliser \texttt{r} à pouvoir servir à
la fois comme \texttt{(int -> int) ref} et comme \texttt{(bool ->
  bool) ref} car le code précédent causerait l'appel d'une fonction
\texttt{int -> int} sur une valeur de type \texttt{bool} (crash
potentiel !).

\bigskip

\itempoint La solution dans un langage comme OCaml est la
« restriction de valeur » : le $v$ dans « \texttt{let $v$=$x$ in
  $t$} » n'est rendu polymorphe (i.e., n'est « généralisé ») que
lorsque $x$ est une « valeur statique » déterminable à la compilation.

\bigskip

\itempoint Dans un langage comme Haskell, le problème ne se pose pas,
car toutes les fonctions et valeurs sont pures.

\end{frame}
%
%% \mode<article>{
%%     \ifodd\value{page}
%%         \clearpage\hbox{}\newpage
%%         \thispagestyle{empty}
%%     \fi}
%
\end{document}
